% Options for packages loaded elsewhere
% Options for packages loaded elsewhere
\PassOptionsToPackage{unicode}{hyperref}
\PassOptionsToPackage{hyphens}{url}
\PassOptionsToPackage{dvipsnames,svgnames,x11names}{xcolor}
%
\documentclass[
  letterpaper,
  DIV=11,
  numbers=noendperiod]{scrartcl}
\usepackage{xcolor}
\usepackage{amsmath,amssymb}
\setcounter{secnumdepth}{-\maxdimen} % remove section numbering
\usepackage{iftex}
\ifPDFTeX
  \usepackage[T1]{fontenc}
  \usepackage[utf8]{inputenc}
  \usepackage{textcomp} % provide euro and other symbols
\else % if luatex or xetex
  \usepackage{unicode-math} % this also loads fontspec
  \defaultfontfeatures{Scale=MatchLowercase}
  \defaultfontfeatures[\rmfamily]{Ligatures=TeX,Scale=1}
\fi
\usepackage{lmodern}
\ifPDFTeX\else
  % xetex/luatex font selection
\fi
% Use upquote if available, for straight quotes in verbatim environments
\IfFileExists{upquote.sty}{\usepackage{upquote}}{}
\IfFileExists{microtype.sty}{% use microtype if available
  \usepackage[]{microtype}
  \UseMicrotypeSet[protrusion]{basicmath} % disable protrusion for tt fonts
}{}
\makeatletter
\@ifundefined{KOMAClassName}{% if non-KOMA class
  \IfFileExists{parskip.sty}{%
    \usepackage{parskip}
  }{% else
    \setlength{\parindent}{0pt}
    \setlength{\parskip}{6pt plus 2pt minus 1pt}}
}{% if KOMA class
  \KOMAoptions{parskip=half}}
\makeatother
% Make \paragraph and \subparagraph free-standing
\makeatletter
\ifx\paragraph\undefined\else
  \let\oldparagraph\paragraph
  \renewcommand{\paragraph}{
    \@ifstar
      \xxxParagraphStar
      \xxxParagraphNoStar
  }
  \newcommand{\xxxParagraphStar}[1]{\oldparagraph*{#1}\mbox{}}
  \newcommand{\xxxParagraphNoStar}[1]{\oldparagraph{#1}\mbox{}}
\fi
\ifx\subparagraph\undefined\else
  \let\oldsubparagraph\subparagraph
  \renewcommand{\subparagraph}{
    \@ifstar
      \xxxSubParagraphStar
      \xxxSubParagraphNoStar
  }
  \newcommand{\xxxSubParagraphStar}[1]{\oldsubparagraph*{#1}\mbox{}}
  \newcommand{\xxxSubParagraphNoStar}[1]{\oldsubparagraph{#1}\mbox{}}
\fi
\makeatother

\usepackage{color}
\usepackage{fancyvrb}
\newcommand{\VerbBar}{|}
\newcommand{\VERB}{\Verb[commandchars=\\\{\}]}
\DefineVerbatimEnvironment{Highlighting}{Verbatim}{commandchars=\\\{\}}
% Add ',fontsize=\small' for more characters per line
\usepackage{framed}
\definecolor{shadecolor}{RGB}{241,243,245}
\newenvironment{Shaded}{\begin{snugshade}}{\end{snugshade}}
\newcommand{\AlertTok}[1]{\textcolor[rgb]{0.68,0.00,0.00}{#1}}
\newcommand{\AnnotationTok}[1]{\textcolor[rgb]{0.37,0.37,0.37}{#1}}
\newcommand{\AttributeTok}[1]{\textcolor[rgb]{0.40,0.45,0.13}{#1}}
\newcommand{\BaseNTok}[1]{\textcolor[rgb]{0.68,0.00,0.00}{#1}}
\newcommand{\BuiltInTok}[1]{\textcolor[rgb]{0.00,0.23,0.31}{#1}}
\newcommand{\CharTok}[1]{\textcolor[rgb]{0.13,0.47,0.30}{#1}}
\newcommand{\CommentTok}[1]{\textcolor[rgb]{0.37,0.37,0.37}{#1}}
\newcommand{\CommentVarTok}[1]{\textcolor[rgb]{0.37,0.37,0.37}{\textit{#1}}}
\newcommand{\ConstantTok}[1]{\textcolor[rgb]{0.56,0.35,0.01}{#1}}
\newcommand{\ControlFlowTok}[1]{\textcolor[rgb]{0.00,0.23,0.31}{\textbf{#1}}}
\newcommand{\DataTypeTok}[1]{\textcolor[rgb]{0.68,0.00,0.00}{#1}}
\newcommand{\DecValTok}[1]{\textcolor[rgb]{0.68,0.00,0.00}{#1}}
\newcommand{\DocumentationTok}[1]{\textcolor[rgb]{0.37,0.37,0.37}{\textit{#1}}}
\newcommand{\ErrorTok}[1]{\textcolor[rgb]{0.68,0.00,0.00}{#1}}
\newcommand{\ExtensionTok}[1]{\textcolor[rgb]{0.00,0.23,0.31}{#1}}
\newcommand{\FloatTok}[1]{\textcolor[rgb]{0.68,0.00,0.00}{#1}}
\newcommand{\FunctionTok}[1]{\textcolor[rgb]{0.28,0.35,0.67}{#1}}
\newcommand{\ImportTok}[1]{\textcolor[rgb]{0.00,0.46,0.62}{#1}}
\newcommand{\InformationTok}[1]{\textcolor[rgb]{0.37,0.37,0.37}{#1}}
\newcommand{\KeywordTok}[1]{\textcolor[rgb]{0.00,0.23,0.31}{\textbf{#1}}}
\newcommand{\NormalTok}[1]{\textcolor[rgb]{0.00,0.23,0.31}{#1}}
\newcommand{\OperatorTok}[1]{\textcolor[rgb]{0.37,0.37,0.37}{#1}}
\newcommand{\OtherTok}[1]{\textcolor[rgb]{0.00,0.23,0.31}{#1}}
\newcommand{\PreprocessorTok}[1]{\textcolor[rgb]{0.68,0.00,0.00}{#1}}
\newcommand{\RegionMarkerTok}[1]{\textcolor[rgb]{0.00,0.23,0.31}{#1}}
\newcommand{\SpecialCharTok}[1]{\textcolor[rgb]{0.37,0.37,0.37}{#1}}
\newcommand{\SpecialStringTok}[1]{\textcolor[rgb]{0.13,0.47,0.30}{#1}}
\newcommand{\StringTok}[1]{\textcolor[rgb]{0.13,0.47,0.30}{#1}}
\newcommand{\VariableTok}[1]{\textcolor[rgb]{0.07,0.07,0.07}{#1}}
\newcommand{\VerbatimStringTok}[1]{\textcolor[rgb]{0.13,0.47,0.30}{#1}}
\newcommand{\WarningTok}[1]{\textcolor[rgb]{0.37,0.37,0.37}{\textit{#1}}}

\usepackage{longtable,booktabs,array}
\usepackage{calc} % for calculating minipage widths
% Correct order of tables after \paragraph or \subparagraph
\usepackage{etoolbox}
\makeatletter
\patchcmd\longtable{\par}{\if@noskipsec\mbox{}\fi\par}{}{}
\makeatother
% Allow footnotes in longtable head/foot
\IfFileExists{footnotehyper.sty}{\usepackage{footnotehyper}}{\usepackage{footnote}}
\makesavenoteenv{longtable}
\usepackage{graphicx}
\makeatletter
\newsavebox\pandoc@box
\newcommand*\pandocbounded[1]{% scales image to fit in text height/width
  \sbox\pandoc@box{#1}%
  \Gscale@div\@tempa{\textheight}{\dimexpr\ht\pandoc@box+\dp\pandoc@box\relax}%
  \Gscale@div\@tempb{\linewidth}{\wd\pandoc@box}%
  \ifdim\@tempb\p@<\@tempa\p@\let\@tempa\@tempb\fi% select the smaller of both
  \ifdim\@tempa\p@<\p@\scalebox{\@tempa}{\usebox\pandoc@box}%
  \else\usebox{\pandoc@box}%
  \fi%
}
% Set default figure placement to htbp
\def\fps@figure{htbp}
\makeatother


% definitions for citeproc citations
\NewDocumentCommand\citeproctext{}{}
\NewDocumentCommand\citeproc{mm}{%
  \begingroup\def\citeproctext{#2}\cite{#1}\endgroup}
\makeatletter
 % allow citations to break across lines
 \let\@cite@ofmt\@firstofone
 % avoid brackets around text for \cite:
 \def\@biblabel#1{}
 \def\@cite#1#2{{#1\if@tempswa , #2\fi}}
\makeatother
\newlength{\cslhangindent}
\setlength{\cslhangindent}{1.5em}
\newlength{\csllabelwidth}
\setlength{\csllabelwidth}{3em}
\newenvironment{CSLReferences}[2] % #1 hanging-indent, #2 entry-spacing
 {\begin{list}{}{%
  \setlength{\itemindent}{0pt}
  \setlength{\leftmargin}{0pt}
  \setlength{\parsep}{0pt}
  % turn on hanging indent if param 1 is 1
  \ifodd #1
   \setlength{\leftmargin}{\cslhangindent}
   \setlength{\itemindent}{-1\cslhangindent}
  \fi
  % set entry spacing
  \setlength{\itemsep}{#2\baselineskip}}}
 {\end{list}}
\usepackage{calc}
\newcommand{\CSLBlock}[1]{\hfill\break\parbox[t]{\linewidth}{\strut\ignorespaces#1\strut}}
\newcommand{\CSLLeftMargin}[1]{\parbox[t]{\csllabelwidth}{\strut#1\strut}}
\newcommand{\CSLRightInline}[1]{\parbox[t]{\linewidth - \csllabelwidth}{\strut#1\strut}}
\newcommand{\CSLIndent}[1]{\hspace{\cslhangindent}#1}



\setlength{\emergencystretch}{3em} % prevent overfull lines

\providecommand{\tightlist}{%
  \setlength{\itemsep}{0pt}\setlength{\parskip}{0pt}}



 


\usepackage{float}
\usepackage{tabularray}
\usepackage[normalem]{ulem}
\usepackage{graphicx}
\usepackage{rotating}
\UseTblrLibrary{siunitx}
\NewTableCommand{\tinytableDefineColor}[3]{\definecolor{#1}{#2}{#3}}
\newcommand{\tinytableTabularrayUnderline}[1]{\underline{#1}}
\newcommand{\tinytableTabularrayStrikeout}[1]{\sout{#1}}
\usepackage{booktabs}
\usepackage{caption}
\usepackage{longtable}
\usepackage{colortbl}
\usepackage{array}
\usepackage{anyfontsize}
\usepackage{multirow}
\KOMAoption{captions}{tableheading}
\usepackage{float}
\floatplacement{figure}{H}
\floatplacement{table}{H}
\usepackage{listings}
\lstset{breaklines=true, basicstyle=\small\ttfamily}
\makeatletter
\@ifpackageloaded{tcolorbox}{}{\usepackage[skins,breakable]{tcolorbox}}
\@ifpackageloaded{fontawesome5}{}{\usepackage{fontawesome5}}
\definecolor{quarto-callout-color}{HTML}{909090}
\definecolor{quarto-callout-note-color}{HTML}{0758E5}
\definecolor{quarto-callout-important-color}{HTML}{CC1914}
\definecolor{quarto-callout-warning-color}{HTML}{EB9113}
\definecolor{quarto-callout-tip-color}{HTML}{00A047}
\definecolor{quarto-callout-caution-color}{HTML}{FC5300}
\definecolor{quarto-callout-color-frame}{HTML}{acacac}
\definecolor{quarto-callout-note-color-frame}{HTML}{4582ec}
\definecolor{quarto-callout-important-color-frame}{HTML}{d9534f}
\definecolor{quarto-callout-warning-color-frame}{HTML}{f0ad4e}
\definecolor{quarto-callout-tip-color-frame}{HTML}{02b875}
\definecolor{quarto-callout-caution-color-frame}{HTML}{fd7e14}
\makeatother
\makeatletter
\@ifpackageloaded{caption}{}{\usepackage{caption}}
\AtBeginDocument{%
\ifdefined\contentsname
  \renewcommand*\contentsname{Table of contents}
\else
  \newcommand\contentsname{Table of contents}
\fi
\ifdefined\listfigurename
  \renewcommand*\listfigurename{List of Figures}
\else
  \newcommand\listfigurename{List of Figures}
\fi
\ifdefined\listtablename
  \renewcommand*\listtablename{List of Tables}
\else
  \newcommand\listtablename{List of Tables}
\fi
\ifdefined\figurename
  \renewcommand*\figurename{Figure}
\else
  \newcommand\figurename{Figure}
\fi
\ifdefined\tablename
  \renewcommand*\tablename{Table}
\else
  \newcommand\tablename{Table}
\fi
}
\@ifpackageloaded{float}{}{\usepackage{float}}
\floatstyle{ruled}
\@ifundefined{c@chapter}{\newfloat{codelisting}{h}{lop}}{\newfloat{codelisting}{h}{lop}[chapter]}
\floatname{codelisting}{Listing}
\newcommand*\listoflistings{\listof{codelisting}{List of Listings}}
\makeatother
\makeatletter
\makeatother
\makeatletter
\@ifpackageloaded{caption}{}{\usepackage{caption}}
\@ifpackageloaded{subcaption}{}{\usepackage{subcaption}}
\makeatother
\usepackage{bookmark}
\IfFileExists{xurl.sty}{\usepackage{xurl}}{} % add URL line breaks if available
\urlstyle{same}
\hypersetup{
  pdftitle={Predicting the effects of behavioral interventions on humans with LLMs},
  colorlinks=true,
  linkcolor={blue},
  filecolor={Maroon},
  citecolor={Blue},
  urlcolor={Blue},
  pdfcreator={LaTeX via pandoc}}


\title{Predicting the effects of behavioral interventions on humans with
LLMs}
\author{}
\date{2026-05-04}
\begin{document}
\maketitle
\begin{abstract}
We test the potential of predicting human behavior in survey experiments
with behavioral clones---large language model (LLM) generated artificial
agents mimicking human participants. To evaluate this potential, we
predict the outcomes of a large-scale field experiment (N =
\textasciitilde22,000) to be conducted in the United States. The
experiment will test 20 text-based interventions designed to increase
trust in climate scientists, assessing a broad range of secondary,
outcomes including self-reported attitudes (e.g., support for funding of
climate science) and behavioral measures (e.g., donations to a
scientific organization). Before the human study is fielded, we will
simulate it with behavioral clones using demographically representative
U.S. participant profiles. We will lock the simulated dataset and
analysis code. Once human data are available, we will compare
clone-based predictions to human results for average treatment effect
and subgroup heterogeneity. We will also compare how well the LLMS
produce outcome distributions, and whether they overemphasize
demographic associations (in other words, whether they stereotype). Our
study will provide a large-scale, credible benchmark test of LLMs'
potential---and limitations---as tools for simulating human subjects in
social science research. All analyses in this preregistration are
demonstrated on simulated placeholder data with identical structure to
the real datasets; no human outcome data have been observed at the time
of registration.
\end{abstract}

\clearpage

\renewcommand*\contentsname{Table of contents}
{
\hypersetup{linkcolor=}
\setcounter{tocdepth}{5}
\tableofcontents
}

\begin{Shaded}
\begin{Highlighting}[]
\CommentTok{\# packages}
\FunctionTok{library}\NormalTok{(tidyverse)}
\FunctionTok{library}\NormalTok{(sandwich)}
\FunctionTok{library}\NormalTok{(lmtest)}
\FunctionTok{library}\NormalTok{(marginaleffects)}
\FunctionTok{library}\NormalTok{(broom)}
\FunctionTok{library}\NormalTok{(MetBrewer)}
\FunctionTok{library}\NormalTok{(patchwork)}
\FunctionTok{library}\NormalTok{(gt)}
\FunctionTok{library}\NormalTok{(tinytable)}
\FunctionTok{library}\NormalTok{(here)}

\CommentTok{\# custom functions}
\FunctionTok{source}\NormalTok{(}\FunctionTok{here}\NormalTok{(}\StringTok{"R/functions/statistics.R"}\NormalTok{))}
\FunctionTok{source}\NormalTok{(}\FunctionTok{here}\NormalTok{(}\StringTok{"R/functions/plots.R"}\NormalTok{))}
\FunctionTok{source}\NormalTok{(}\FunctionTok{here}\NormalTok{(}\StringTok{"R/functions/tables.R"}\NormalTok{))}
\FunctionTok{source}\NormalTok{(}\FunctionTok{here}\NormalTok{(}\StringTok{"R/functions/reporting.R"}\NormalTok{))}

\CommentTok{\# simulated (random) data}
\NormalTok{human\_data }\OtherTok{\textless{}{-}} \FunctionTok{readRDS}\NormalTok{(}\FunctionTok{here}\NormalTok{(}\StringTok{"data/simulation\_preregistration/human\_data\_preregistration.rds"}\NormalTok{))}
\NormalTok{llm\_data   }\OtherTok{\textless{}{-}} \FunctionTok{readRDS}\NormalTok{(}\FunctionTok{here}\NormalTok{(}\StringTok{"data/simulation\_preregistration/llm\_data\_preregistration.rds"}\NormalTok{))}

\CommentTok{\# Preregistered random split of human sample into two halves.}
\CommentTok{\# human\_data\_1 is the reference half used in all primary analyses.}
\CommentTok{\# human\_data\_2 is the baseline comparator (human{-}human ceiling for all metrics).}
\FunctionTok{set.seed}\NormalTok{(}\DecValTok{42}\NormalTok{)}
\NormalTok{split\_ids    }\OtherTok{\textless{}{-}} \FunctionTok{sample}\NormalTok{(}\FunctionTok{unique}\NormalTok{(human\_data}\SpecialCharTok{$}\NormalTok{id), }\AttributeTok{size =} \FunctionTok{floor}\NormalTok{(}\FunctionTok{length}\NormalTok{(}\FunctionTok{unique}\NormalTok{(human\_data}\SpecialCharTok{$}\NormalTok{id)) }\SpecialCharTok{/} \DecValTok{2}\NormalTok{))}
\NormalTok{human\_data\_1 }\OtherTok{\textless{}{-}}\NormalTok{ human\_data }\SpecialCharTok{|\textgreater{}} \FunctionTok{filter}\NormalTok{(id }\SpecialCharTok{\%in\%}\NormalTok{ split\_ids)}
\NormalTok{human\_data\_2 }\OtherTok{\textless{}{-}}\NormalTok{ human\_data }\SpecialCharTok{|\textgreater{}} \FunctionTok{filter}\NormalTok{(}\SpecialCharTok{!}\NormalTok{id }\SpecialCharTok{\%in\%}\NormalTok{ split\_ids)}

\CommentTok{\# interventions}
\NormalTok{interventions }\OtherTok{\textless{}{-}} \FunctionTok{read\_csv}\NormalTok{(}\FunctionTok{here}\NormalTok{(}\StringTok{"data/interventions.csv"}\NormalTok{))}

\CommentTok{\# actual data later}
\CommentTok{\# human\_data  \textless{}{-} readRDS("data/cleaned\_llm.rds")}
\CommentTok{\# llm\_data \textless{}{-} readRDS("data/cleaned\_llm.rds")}
\end{Highlighting}
\end{Shaded}

\begin{tcolorbox}[enhanced jigsaw, left=2mm, rightrule=.15mm, titlerule=0mm, colframe=quarto-callout-important-color-frame, opacitybacktitle=0.6, colback=white, opacityback=0, breakable, bottomtitle=1mm, coltitle=black, leftrule=.75mm, arc=.35mm, toprule=.15mm, title=\textcolor{quarto-callout-important-color}{\faExclamation}\hspace{0.5em}{Important}, colbacktitle=quarto-callout-important-color!10!white, bottomrule=.15mm, toptitle=1mm]

\textbf{Placeholder data.} Both datasets used throughout are simulated
random noise with the correct structure. The analysis code will be
applied to real human and LLM clone data without modification once both
are available.

\end{tcolorbox}

\subsection{Motivation}\label{motivation}

The aim of this project is to understand how well large language models
(LLMs) can simulate human behavior in survey experiments.

Large language models (LLMs) are trained on extensive human-generated
data. Consequently, they reproduce patterns of human judgment and
reasoning (Binz and Schulz 2023; Aher, Arriaga, and Kalai 2023). In
controlled, text-based interactions, humans sometimes struggle to
distinguish between responses generated by LLMs and those produced by
other humans (Jones et al. 2025; but see Pagan et al. 2025).

This apparent capability of LLMs to mimic human behavior has generated
substantial interest among social scientists. Running surveys and
experiments with human subjects is costly and time-consuming. If LLMs
could reliably approximate human behavior, they would offer a fast and
low-cost alternative. While it seems unrealistic at this point that LLMs
could substitute human subjects entirely (see, e.g., Schröder et al.
2025), they might still open up substantial new opportunities for social
science research. Academic publishing systems tend to incentivize
positive results (Nosek, Spies, and Motyl 2012; Smaldino and McElreath
2016), which can discourage the testing of novel or high-risk ideas.
Approximative LLM simulations could serve as a low-risk tool for
pretesting such novel ideas before committing resources to human data
collection. LLMs could also help testing the boundaries of established
theories. Scientific theories---at least in the social sciences---are
often highly context-dependent. To improve theory building, researchers
have called for more systematic exploration of these contextual
variations (Almaatouq et al. 2022). LLMs could make such exploration
more feasible: they could simulate a multitude of variations of an
experiment, thereby pointing researchers to critical context factors
that could be tested in follow-up studies on human participants.

Recent studies have demonstrated promising results for LLMs in
simulating public opinion (Kaiser et al. 2025; Wang, Morgenstern, and
Dickerson 2025; Argyle et al. 2023; Durmus et al. 2024; Santurkar et al.
2023; Bisbee et al. 2024; Shrestha et al. 2024) and in predicting the
outcomes of experiments with human participants (Hewitt et al. 2025;
Chen, Hu, and Lu 2025; Hu et al. 2025; Guo et al. 2025; Lippert et al.
2024; Bisbee et al. 2024). At the same time, other work has identified
important limitations of LLM simulations, including systematic
misrepresentation of certain identity groups (Wang, Morgenstern, and
Dickerson 2025), a general tendency of overestimating experimental
effects (Doudkin, Pataranutaporn, and Maes 2025; Cui, Li, and Zhou
2024), and high volatility of simulation results, depending on model and
prompt choice (Schröder et al. 2025).

Currently, interpreting these mixed findings is difficult, as most
existing studies face at least one of the following limitations. First,
some studies rely on already published experiments, raising concerns
that LLMs may have been trained on the very data they are asked to
simulate. These studies might only test LLMs capacity to reproduce, but
not to predict. Second, even studies that plausibly rely on unpublished
data are rarely preregistered. Given that LLM performance is highly
sensitive to model choice and prompting, preregistered predictions are
crucial for credible evaluation. Third, some studies operate at the
experiment level, providing LLMs with a full description of the study
and requesting a single aggregate prediction. While informative about
average effects, this approach precludes the analysis of
individual-level heterogeneity.

In this project, we aim to provide a robust test of how well LLMs can
simulate human behavior in survey experiments. We address limitations of
previous studies by leveraging a unique opportunity to evaluate LLM
simulations in the context of a large-scale field experiment, involving
approximately 21,000 participants from the United States (US), testing
20 text-based interventions designed to increase trust in climate
scientists. The study assesses a broad range of outcomes, including
self-reported attitudes (e.g., trust in scientists) and behavioral
measures (e.g., donations to a scientific organization). Before running
the study with human participants, we will use multiple LLMs and
prompting strategies to simulate individual-level survey responses
representative of the US population. We will then collect the human data
and test how closely they match the data simulated by the LLMs. This
design allows for a large-scale, credible benchmark test of LLMs'
potential---and limitations---as tools for simulating human subjects in
social science research.

\subsection{Research Questions}\label{research-questions}

The main outcome of the megastudy is trust in climate scientists, but
several secondary and tertiary outcomes are preregistered. Analyses are
organized in three sections that mirror one another: we first examine
average effects, then subgroup differences, and finally demographic
calibration at the individual level.

\textbf{Section 1 --- Average Treatment Effects.} We ask (a) how well
LLM clones predict the average treatment effect (ATE) of each
intervention across all outcomes, and (b) to which extent clones produce
similar response distributions for the different outcomes (in the
control condition).

\textbf{Section 2 --- Subgroup Effects.} We ask (a) how closely clones
predict differences in intervention effects across six demographic
moderators (gender, age, race, education, income, partisan identity) for
all outcomes, and (b) how closely within-subgroup response distributions
in the control condition match between humans and clones for the primary
outcome.

\textbf{Section 3 --- Demographic Baseline Calibration.} Independent of
treatment effects, we ask whether clone baseline outcome levels are
calibrated to the assigned demographic profile: (a) do clones reproduce
the correct group mean for each demographic group in the control
condition, and (b) do demographic variables predict clone outcomes with
the same magnitude as in humans, or do clones over-rely on demographic
cues? Analyses are run separately for all outcomes; results are
illustrated with the primary outcome (\texttt{trust\_multidimensional}).

\subsection{Data}\label{data}

\subsubsection{Human data}\label{human-data}

The human dataset will come from a large-scale between-subjects field
experiment to be conducted in the United States (N ≈ 22,000).
Participants will be randomly assigned to one of 20 text-based
interventions designed to increase trust in climate scientists, or to a
control condition. The human study is preregistered separately
(\url{https://doi.org/10.5281/zenodo.20160212}). \textbf{No human data
has been colleced at the moment of registering this study.}

To provide a reference baseline for interpreting LLM performance, we
split the human sample into two equal halves using a preregistered
random split (\texttt{set.seed(42)}). \textbf{Human 1} (the reference
half) is used in all primary analyses and serves as the ground-truth
against which both LLMs and the second human half are compared.
\textbf{Human 2} provides a human-human reference: the similarity
between the two halves quantifies how well human data predict other
human data, which is the upper bound for any predictive method. All
comparison tables and the pooled scatter plot report metrics for both
comparisons side by side.

\paragraph{Interventions}\label{interventions}

Table~\ref{tbl-interventions} provides a summary of the different
interventions tested in the megastudy. In this project, we only generate
LLM clone responses for non-interactive interventions. As a result, we
only test 16 interventions (none of the three LLM-chatbot interventions,
neither the `Value similarity' quiz intervention).

\begin{Shaded}
\begin{Highlighting}[]
\CommentTok{\# build table data with category headers manually}

\NormalTok{intervention\_table\_data }\OtherTok{\textless{}{-}}\NormalTok{ interventions }\SpecialCharTok{|\textgreater{}}
  \FunctionTok{select}\NormalTok{(title, summary, tag) }\SpecialCharTok{|\textgreater{}}
  \FunctionTok{mutate}\NormalTok{(}
    \AttributeTok{tag =} \FunctionTok{factor}\NormalTok{(tag, }\AttributeTok{levels =} \FunctionTok{c}\NormalTok{(}
      \StringTok{"Collaboration and peer{-}review"}\NormalTok{,}
      \StringTok{"Scientific methods and results"}\NormalTok{,}
      \StringTok{"Applications and impact"}\NormalTok{,}
      \StringTok{"Others\textquotesingle{} endorsement"}\NormalTok{,}
      \StringTok{"Values"}\NormalTok{,}
      \StringTok{"LLM{-}chatbot"}\NormalTok{,}
      \StringTok{"Other"}
\NormalTok{    ))}
\NormalTok{  ) }\SpecialCharTok{|\textgreater{}}
  \FunctionTok{arrange}\NormalTok{(tag) }\SpecialCharTok{|\textgreater{}}
  \FunctionTok{mutate}\NormalTok{(}\AttributeTok{number =} \FunctionTok{as.character}\NormalTok{(}\FunctionTok{row\_number}\NormalTok{())) }\SpecialCharTok{|\textgreater{}} 
  \FunctionTok{group\_by}\NormalTok{(tag) }\SpecialCharTok{|\textgreater{}}
  \FunctionTok{group\_modify}\NormalTok{(}\SpecialCharTok{\textasciitilde{}} \FunctionTok{bind\_rows}\NormalTok{(}
    \FunctionTok{tibble}\NormalTok{(}\AttributeTok{title =} \FunctionTok{as.character}\NormalTok{(.y}\SpecialCharTok{$}\NormalTok{tag), }\AttributeTok{summary =} \StringTok{""}\NormalTok{),}
\NormalTok{    .x}
\NormalTok{  )) }\SpecialCharTok{|\textgreater{}}
  \FunctionTok{ungroup}\NormalTok{() }\SpecialCharTok{|\textgreater{}}
  \FunctionTok{mutate}\NormalTok{(}\StringTok{\textasciigrave{}}\AttributeTok{ }\StringTok{\textasciigrave{}} \OtherTok{=} \FunctionTok{ifelse}\NormalTok{(title }\SpecialCharTok{\%in\%} \FunctionTok{levels}\NormalTok{(tag), }\StringTok{""}\NormalTok{, number)) }\SpecialCharTok{|\textgreater{}}
  \FunctionTok{select}\NormalTok{(}\StringTok{\textasciigrave{}}\AttributeTok{ }\StringTok{\textasciigrave{}}\NormalTok{, title, summary) }\SpecialCharTok{|\textgreater{}} 
  \FunctionTok{rename}\NormalTok{(}\StringTok{\textasciigrave{}}\AttributeTok{Intervention Title}\StringTok{\textasciigrave{}} \OtherTok{=}\NormalTok{ title, }
         \AttributeTok{Summary =}\NormalTok{ summary)}

\CommentTok{\# identify header rows}
\NormalTok{header\_rows }\OtherTok{\textless{}{-}} \FunctionTok{which}\NormalTok{(intervention\_table\_data}\SpecialCharTok{$}\StringTok{\textasciigrave{}}\AttributeTok{ }\StringTok{\textasciigrave{}} \SpecialCharTok{==} \StringTok{""}\NormalTok{)}

\NormalTok{intervention\_table }\OtherTok{\textless{}{-}}\NormalTok{ intervention\_table\_data }\SpecialCharTok{|\textgreater{}}
  \FunctionTok{tt}\NormalTok{(}\AttributeTok{width =} \FunctionTok{c}\NormalTok{(}\FloatTok{0.05}\NormalTok{, }\FloatTok{0.25}\NormalTok{, }\FloatTok{0.70}\NormalTok{)) }\SpecialCharTok{|\textgreater{}}
  \FunctionTok{format\_tt}\NormalTok{(}\AttributeTok{escape =} \ConstantTok{TRUE}\NormalTok{) }\SpecialCharTok{|\textgreater{}}
  \FunctionTok{style\_tt}\NormalTok{(}\AttributeTok{align =} \StringTok{"l"}\NormalTok{, }\AttributeTok{alignv =} \StringTok{"t"}\NormalTok{) }\SpecialCharTok{|\textgreater{}}
  \FunctionTok{style\_tt}\NormalTok{(}\AttributeTok{i =} \DecValTok{0}\NormalTok{, }\AttributeTok{bold =} \ConstantTok{TRUE}\NormalTok{) }\SpecialCharTok{|\textgreater{}}
  \FunctionTok{style\_tt}\NormalTok{(}\AttributeTok{i =}\NormalTok{ header\_rows, }\AttributeTok{bold =} \ConstantTok{TRUE}\NormalTok{, }\AttributeTok{j =} \DecValTok{2}\NormalTok{, }\AttributeTok{colspan =} \DecValTok{2}\NormalTok{) }\SpecialCharTok{|\textgreater{}}
  \FunctionTok{style\_tt}\NormalTok{(}\AttributeTok{fontsize =} \FloatTok{0.85}\NormalTok{)}
\end{Highlighting}
\end{Shaded}

\begin{table}

\caption{\label{tbl-interventions}Overview of interventions included in
the megastudy.}

\centering{

\begin{Shaded}
\begin{Highlighting}[]
\NormalTok{intervention\_table}
\end{Highlighting}
\end{Shaded}

}

\end{table}%

\begin{Shaded}
\begin{Highlighting}[]
\NormalTok{intervention\_table }\SpecialCharTok{|\textgreater{}} 
  \FunctionTok{theme\_latex}\NormalTok{(}\AttributeTok{multipage =} \ConstantTok{TRUE}\NormalTok{, }\AttributeTok{rowhead =} \DecValTok{1}\NormalTok{,}
              \AttributeTok{outer =} \StringTok{"caption=\{Overview of interventions included in the megastudy.\}, label=\{tbl{-}interventions\}"}\NormalTok{)}
\end{Highlighting}
\end{Shaded}

\begin{longtblr}[         %% tabularray outer open
caption={Overview of interventions included in the megastudy.}, label={tbl-interventions},
]                     %% tabularray outer close
{                     %% tabularray inner open
width={1\linewidth},
colspec={X[0.05]X[0.25]X[0.7]},
hline{2}={1-3}{solid, black, 0.05em},
hline{1}={1-3}{solid, black, 0.08em},
hline{29}={1-3}{solid, black, 0.08em},
row{1}={}{font=\bfseries\fontsize{0.85em}{1.15em}\selectfont, halign=l, valign=h},
row{3,4,6,7,8,10,11,13,14,16,17,19,20,21,23,24,25,26,27,28}={}{font=\fontsize{0.85em}{1.15em}\selectfont, halign=l, valign=h},
cell{2,5,9,12,15,18,22}{1}={}{font=\fontsize{0.85em}{1.15em}\selectfont, halign=l, valign=h},
cell{2,5,9,12,15,18,22}{2}={c=2}{font=\bfseries\fontsize{0.85em}{1.15em}\selectfont, halign=l, valign=h, wd=0.95\linewidth},
cell{2,5,9,12,15,18,22}{3}={}{font=\fontsize{0.85em}{1.15em}\selectfont, halign=l, valign=h},
rowhead=1,
}                     %% tabularray inner close
& Intervention Title & Summary \\
& Collaboration and peer-review &  \\
1 & Interview Prof. Maraun & Climate scientist Prof. Douglas Maraun at the University of Graz in Austria stresses the collaborative and self-correcting process of climate science. \\
2 & Peer-review & What makes climate science trustworthy is the process of independent peer-review. \\
& Scientific methods and results &  \\
3 & Measurement \& modeling (1) & Climate scientists use sophisticated measurement and computational modeling techniques to surveil climate and predict how it changes. \\
4 & Measurement \& modeling (2) & Climate scientists are primarily natural scientists (e.g., biologists, physicists). They use use sophisticated tools and quantiative methods to measure and predict climate change. \\
5 & Model accuracy & This is an edited version of a real news article showcasing that even old climate models, despite some flaws, were remarkably correct in predicting global warming. \\
& Applications and impact &  \\
6 & Portrait Prof. Cherry & Todd Cherry, a scientist focused on climate issues, is integrated in his local community, and does work that is relevant for this community. \\
7 & Extreme weather predictions & Showcases how climate science predicts and helps adapting to different extreme weather events (blizzards, floods, wildfires). Takes into account a participant's state and adresses the extreme weather event most common in that state. \\
& Others' endorsement &  \\
8 & Corporate reliance & Insurance companies and large corporations rely on climate scientists' projections. \\
9 & Former skeptics & Former climate change skeptics Jennifer Rukavina (television meteorologist) and Bob Inglis (former Repubican congressman) explain how they came to change their mind. \\
& Values &  \\
10 & Value similarity & The quiz 'Which Type of Climate Scientist Are You?' highlights dimensions of climate scientists' trustworthiness. By providing participants with their personalized climate scientists profile, the intervention intends to create perceptions of value similiarity and identification. \\
11 & Interview Prof. Sebille & Prof. Erik van Sebille, a climate scientist and oceanographer at Utrecht University, Netherlands, mentions harmful consequences of climate change on oceans and humans, and how he cares about preventing these consequences. \\
& LLM-chatbot &  \\
12 & LLM chatbot (1) & LLM-chatbot \\
13 & LLM chatbot (2) & LLM-chatbot \\
14 & LLM-chatbot (3) & LLM-chatbot \\
& Other &  \\
15 & Social justice & In the United States, the wealthiest 10\% of the population are responsible for roughly 40\% of the country’s total greenhouse gas emissions. Climate scientists provide evidence to hold the emitters accountable. \\
16 & Funding & Correcting potential misperceptions on the amount and sources of climate science funding. Showcases that climate science receives relatively little pubilc and private funding. \\
17 & Oil industry misinformation & Oil companies have spent decades financing large propaganda campaigns to cast doubt on the existence climate change and the credibility of climate scientists. \\
18 & High public trust & Correct potential misperceptions of how many Americans trusts climate scientists. A majority of Americans trusts climate scientists at least to some extent. \\
19 & Scientist community helpers & Climate scientists are members of local communities and their work helps local communities in times of cliamte disasters (e.g., floods and wildfires). \\
20 & Consensus & Correcting potential misperceptions on the level of agreement among climate scientists on cliamte change, and climate change related information. \\
\end{longtblr}

\paragraph{Outcomes}\label{outcomes}

Outcomes are organized in three tiers reflecting theoretical proximity
to the primary target of the interventions:

\begin{itemize}
\tightlist
\item
  \textbf{Primary}: \texttt{trust\_multidimensional} --- a composite
  measure of trust in climate scientists averaging four subdimensions
  (competence, integrity, benevolence, openness), each assessed with
  three items (12 items total, scored 0--100). This is the pre-specified
  main endpoint.
\item
  \textbf{Secondary} (five outcomes): single-item trust
  (\texttt{trust\_post}), single-item distrust
  (\texttt{distrust\_post}), perceptions of science funding
  (\texttt{funding\_perceptions}), perceived appropriate policy role of
  scientists (\texttt{policy\_role\_mean}), and institutional trust
  (\texttt{inst\_trust\_mean}).
\item
  \textbf{Tertiary} (five outcomes): climate belief
  (\texttt{belief\_post}), climate concern (\texttt{concern\_mean}),
  general climate policy support (\texttt{policy\_general}), specific
  climate policy support (\texttt{policy\_specific\_mean}), and
  pro-climate behavior intentions (\texttt{behavior\_mean}).
\item
  \textbf{Binary}: newsletter sign-up (\texttt{newsletter\_signup}) ---
  a behavioral outcome recording whether the participant subscribed to a
  climate science newsletter. Analyzed separately with logistic
  regression; excluded from distribution analyses.
\end{itemize}

ATE analyses cover all 13 outcomes. Analyses that break results down by
outcome --- response distribution comparisons (Section 1) and subgroup
heterogeneity (Section 2) --- use an illustrative subset of four
outcomes spanning the tiers: multidimensional trust (primary), donation
to AMS (secondary), funding perceptions (secondary), and general climate
policy support (tertiary). Subgroup analyses use six demographic
moderators: gender, age band, race, education, income, and partisan
identity.

\subsubsection{LLM generated data}\label{llm-generated-data}

\paragraph{Profiles}\label{profiles}

Clone profiles were constructed by drawing 9,000 individuals from four
recent waves of the General Social Survey (GSS; 2018, 2021, 2022, 2024).
Within-year normalized sampling weights were applied so that each wave
contributes equal sampling probability to the pool, while preserving
within-year relative weights. Each profile was assigned to exactly one
experimental condition (between-subjects, matching the human study
design). The exact procedure is documented in the project repository.

Profiles encode the following demographic characteristics: age, sex,
race/ethnicity (racial self-identification and Hispanic origin),
educational attainment, household income, household size,
self-identified social class, U.S. region, partisan identification,
political ideology, religious affiliation, religious intensity, and
confidence in the scientific community. Each clone then completed the
survey experiment under the same conditions as human participants.

\begin{Shaded}
\begin{Highlighting}[]
\NormalTok{profiles }\OtherTok{\textless{}{-}} \FunctionTok{read\_csv}\NormalTok{(}\FunctionTok{here}\NormalTok{(}\StringTok{"clone\_profiles"}\NormalTok{, }\StringTok{"profiles.csv"}\NormalTok{), }\AttributeTok{show\_col\_types =} \ConstantTok{FALSE}\NormalTok{)}

\NormalTok{meta\_vars }\OtherTok{\textless{}{-}} \FunctionTok{c}\NormalTok{(}\StringTok{"profile\_id"}\NormalTok{, }\StringTok{"year"}\NormalTok{, }\StringTok{"id"}\NormalTok{, }\StringTok{"wtssps"}\NormalTok{, }\StringTok{"w\_pooled"}\NormalTok{,}
               \StringTok{"profile\_basic"}\NormalTok{, }\StringTok{"profile\_detailed"}\NormalTok{)}
\NormalTok{prof\_vars }\OtherTok{\textless{}{-}} \FunctionTok{setdiff}\NormalTok{(}\FunctionTok{names}\NormalTok{(profiles), meta\_vars)}
\NormalTok{yrs       }\OtherTok{\textless{}{-}} \FunctionTok{sort}\NormalTok{(}\FunctionTok{unique}\NormalTok{(profiles}\SpecialCharTok{$}\NormalTok{year))}

\NormalTok{yr\_ns }\OtherTok{\textless{}{-}}\NormalTok{ profiles }\SpecialCharTok{|\textgreater{}} \FunctionTok{count}\NormalTok{(year) }\SpecialCharTok{|\textgreater{}} \FunctionTok{deframe}\NormalTok{()}

\NormalTok{comp\_tbl }\OtherTok{\textless{}{-}}\NormalTok{ prof\_vars }\SpecialCharTok{|\textgreater{}}
  \FunctionTok{map\_dfr}\NormalTok{(\textbackslash{}(v) \{}
\NormalTok{    row }\OtherTok{\textless{}{-}} \FunctionTok{tibble}\NormalTok{(}\AttributeTok{Variable =}\NormalTok{ v)}
    \ControlFlowTok{for}\NormalTok{ (yr }\ControlFlowTok{in}\NormalTok{ yrs) \{}
\NormalTok{      d }\OtherTok{\textless{}{-}}\NormalTok{ profiles }\SpecialCharTok{|\textgreater{}} \FunctionTok{filter}\NormalTok{(year }\SpecialCharTok{==}\NormalTok{ yr) }\SpecialCharTok{|\textgreater{}} \FunctionTok{pull}\NormalTok{(}\FunctionTok{all\_of}\NormalTok{(v))}
\NormalTok{      row[[}\FunctionTok{paste0}\NormalTok{(yr, }\StringTok{" (N="}\NormalTok{, yr\_ns[[}\FunctionTok{as.character}\NormalTok{(yr)]], }\StringTok{")"}\NormalTok{)]] }\OtherTok{\textless{}{-}} \FunctionTok{sum}\NormalTok{(}\FunctionTok{is.na}\NormalTok{(d))}
\NormalTok{    \}}
\NormalTok{    row}
\NormalTok{  \})}

\NormalTok{comp\_tbl }\SpecialCharTok{|\textgreater{}}
  \FunctionTok{tt}\NormalTok{() }\SpecialCharTok{|\textgreater{}}
  \FunctionTok{style\_tt}\NormalTok{(}\AttributeTok{j =} \DecValTok{1}\NormalTok{, }\AttributeTok{align =} \StringTok{"l"}\NormalTok{) }\SpecialCharTok{|\textgreater{}}
  \FunctionTok{style\_tt}\NormalTok{(}\AttributeTok{j =} \FunctionTok{seq}\NormalTok{(}\DecValTok{2}\NormalTok{, }\FunctionTok{ncol}\NormalTok{(comp\_tbl)), }\AttributeTok{align =} \StringTok{"r"}\NormalTok{) }\SpecialCharTok{|\textgreater{}}
  \FunctionTok{style\_tt}\NormalTok{(}\AttributeTok{i =} \DecValTok{0}\NormalTok{, }\AttributeTok{bold =} \ConstantTok{TRUE}\NormalTok{)}
\end{Highlighting}
\end{Shaded}

\begin{table}

\caption{\label{tbl-profile-completeness}Missing values per variable by
GSS wave. Column headers show the wave year and total N for that wave in
brackets.}

\centering{

\centering
\begin{tblr}[         %% tabularray outer open
]                     %% tabularray outer close
{                     %% tabularray inner open
colspec={Q[]Q[]Q[]Q[]Q[]},
hline{2}={1-5}{solid, black, 0.05em},
hline{1}={1-5}{solid, black, 0.08em},
hline{18}={1-5}{solid, black, 0.08em},
cell{1}{1}={}{font=\bfseries, halign=l},
cell{1}{2}={}{font=\bfseries, halign=r},
cell{1}{3}={}{font=\bfseries, halign=r},
cell{1}{4}={}{font=\bfseries, halign=r},
cell{1}{5}={}{font=\bfseries, halign=r},
cell{2-17}{1}={}{halign=l},
cell{2-17}{2}={}{halign=r},
cell{2-17}{3}={}{halign=r},
cell{2-17}{4}={}{halign=r},
cell{2-17}{5}={}{halign=r},
}                     %% tabularray inner close
Variable & 2018 (N=1867) & 2021 (N=2687) & 2022 (N=2214) & 2024 (N=2232) \\
age & 0 & 0 & 0 & 0 \\
sex & 0 & 0 & 0 & 0 \\
racecen1 & 0 & 0 & 0 & 0 \\
hispanic & 0 & 0 & 0 & 0 \\
degree & 0 & 0 & 0 & 0 \\
income16 & 0 & 0 & 0 & 0 \\
hompop & 0 & 60 & 1112 & 1098 \\
class & 0 & 0 & 0 & 0 \\
region & 0 & 0 & 0 & 0 \\
srcbelt & 0 & 0 & 40 & 2232 \\
partyid & 0 & 0 & 0 & 0 \\
polviews & 0 & 0 & 0 & 0 \\
relig & 0 & 0 & 0 & 0 \\
reborn & 22 & 909 & 29 & 38 \\
reliten & 7 & 2687 & 1092 & 1300 \\
consci & 652 & 907 & 753 & 783 \\
\end{tblr}

}

\end{table}%

\begin{Shaded}
\begin{Highlighting}[]
\NormalTok{race\_var }\OtherTok{\textless{}{-}} \ControlFlowTok{if}\NormalTok{ (}\StringTok{"racecen1"} \SpecialCharTok{\%in\%} \FunctionTok{names}\NormalTok{(profiles)) }\StringTok{"racecen1"} \ControlFlowTok{else} \StringTok{"race"}

\NormalTok{age\_tbl }\OtherTok{\textless{}{-}}\NormalTok{ profiles }\SpecialCharTok{|\textgreater{}}
  \FunctionTok{mutate}\NormalTok{(}\AttributeTok{Category =} \FunctionTok{cut}\NormalTok{(age,}
    \AttributeTok{breaks =} \FunctionTok{c}\NormalTok{(}\DecValTok{17}\NormalTok{, }\DecValTok{29}\NormalTok{, }\DecValTok{44}\NormalTok{, }\DecValTok{59}\NormalTok{, }\ConstantTok{Inf}\NormalTok{),}
    \AttributeTok{labels =} \FunctionTok{c}\NormalTok{(}\StringTok{"18–29"}\NormalTok{, }\StringTok{"30–44"}\NormalTok{, }\StringTok{"45–59"}\NormalTok{, }\StringTok{"60+"}\NormalTok{)}
\NormalTok{  )) }\SpecialCharTok{|\textgreater{}}
  \FunctionTok{count}\NormalTok{(Category) }\SpecialCharTok{|\textgreater{}}
  \FunctionTok{mutate}\NormalTok{(}\AttributeTok{pct =} \FunctionTok{round}\NormalTok{(}\DecValTok{100} \SpecialCharTok{*}\NormalTok{ n }\SpecialCharTok{/} \FunctionTok{sum}\NormalTok{(n), }\DecValTok{1}\NormalTok{), }\AttributeTok{group =} \StringTok{"Age"}\NormalTok{)}

\NormalTok{sex\_tbl }\OtherTok{\textless{}{-}}\NormalTok{ profiles }\SpecialCharTok{|\textgreater{}}
  \FunctionTok{count}\NormalTok{(}\AttributeTok{Category =}\NormalTok{ sex) }\SpecialCharTok{|\textgreater{}}
  \FunctionTok{mutate}\NormalTok{(}\AttributeTok{pct =} \FunctionTok{round}\NormalTok{(}\DecValTok{100} \SpecialCharTok{*}\NormalTok{ n }\SpecialCharTok{/} \FunctionTok{sum}\NormalTok{(n), }\DecValTok{1}\NormalTok{), }\AttributeTok{group =} \StringTok{"Sex"}\NormalTok{)}

\NormalTok{race\_tbl }\OtherTok{\textless{}{-}}\NormalTok{ profiles }\SpecialCharTok{|\textgreater{}}
  \FunctionTok{count}\NormalTok{(}\AttributeTok{Category =}\NormalTok{ .data[[race\_var]]) }\SpecialCharTok{|\textgreater{}}
  \FunctionTok{mutate}\NormalTok{(}\AttributeTok{pct =} \FunctionTok{round}\NormalTok{(}\DecValTok{100} \SpecialCharTok{*}\NormalTok{ n }\SpecialCharTok{/} \FunctionTok{sum}\NormalTok{(n), }\DecValTok{1}\NormalTok{), }\AttributeTok{group =} \StringTok{"Race / Ethnicity"}\NormalTok{)}

\NormalTok{demo\_tbl }\OtherTok{\textless{}{-}} \FunctionTok{bind\_rows}\NormalTok{(age\_tbl, sex\_tbl, race\_tbl) }\SpecialCharTok{|\textgreater{}}
  \FunctionTok{mutate}\NormalTok{(}\AttributeTok{Category =} \FunctionTok{as.character}\NormalTok{(Category))}

\NormalTok{demo\_tbl }\SpecialCharTok{|\textgreater{}}
  \FunctionTok{select}\NormalTok{(Category, n, pct) }\SpecialCharTok{|\textgreater{}}
  \FunctionTok{rename}\NormalTok{(}\AttributeTok{N =}\NormalTok{ n, }\StringTok{"\%"} \OtherTok{=}\NormalTok{ pct) }\SpecialCharTok{|\textgreater{}}
  \FunctionTok{tt}\NormalTok{(}\AttributeTok{width =} \FloatTok{0.5}\NormalTok{) }\SpecialCharTok{|\textgreater{}}
  \FunctionTok{group\_tt}\NormalTok{(}\AttributeTok{i =} \FunctionTok{list}\NormalTok{(}
    \StringTok{"Age"}              \OtherTok{=} \FunctionTok{which}\NormalTok{(demo\_tbl}\SpecialCharTok{$}\NormalTok{group }\SpecialCharTok{==} \StringTok{"Age"}\NormalTok{)[}\DecValTok{1}\NormalTok{],}
    \StringTok{"Sex"}              \OtherTok{=} \FunctionTok{which}\NormalTok{(demo\_tbl}\SpecialCharTok{$}\NormalTok{group }\SpecialCharTok{==} \StringTok{"Sex"}\NormalTok{)[}\DecValTok{1}\NormalTok{],}
    \StringTok{"Race / Ethnicity"} \OtherTok{=} \FunctionTok{which}\NormalTok{(demo\_tbl}\SpecialCharTok{$}\NormalTok{group }\SpecialCharTok{==} \StringTok{"Race / Ethnicity"}\NormalTok{)[}\DecValTok{1}\NormalTok{]}
\NormalTok{  )) }\SpecialCharTok{|\textgreater{}}
  \FunctionTok{style\_tt}\NormalTok{(}\AttributeTok{align =} \StringTok{"l"}\NormalTok{) }\SpecialCharTok{|\textgreater{}}
  \FunctionTok{style\_tt}\NormalTok{(}\AttributeTok{i =} \DecValTok{0}\NormalTok{, }\AttributeTok{bold =} \ConstantTok{TRUE}\NormalTok{) }\SpecialCharTok{|\textgreater{}}
  \FunctionTok{style\_tt}\NormalTok{(}\AttributeTok{i =} \StringTok{"groupi"}\NormalTok{, }\AttributeTok{bold =} \ConstantTok{TRUE}\NormalTok{) }\SpecialCharTok{|\textgreater{}}
  \FunctionTok{format\_tt}\NormalTok{(}\AttributeTok{escape =} \ConstantTok{TRUE}\NormalTok{)}
\end{Highlighting}
\end{Shaded}

\begin{table}

\caption{\label{tbl-profile-demographics}Distribution of key quota
variables in the LLM clone profile pool (N = 9,000).}

\centering{

\centering
\begin{tblr}[         %% tabularray outer open
]                     %% tabularray outer close
{                     %% tabularray inner open
width={0.5\linewidth},
colspec={X[]X[]X[]},
hline{2}={1-3}{solid, black, 0.05em},
hline{1}={1-3}{solid, black, 0.08em},
hline{18}={1-3}{solid, black, 0.08em},
row{1}={}{font=\bfseries, halign=l},
row{3,4,5,6,8,9,11,12,13,14,15,16,17}={}{halign=l},
cell{2,7,10}{1}={c=3}{font=\bfseries, halign=l},
cell{2,7,10}{2}={}{font=\bfseries, halign=l},
cell{2,7,10}{3}={}{font=\bfseries, halign=l},
}                     %% tabularray inner close
Category & N & \% \\
Age & Age & Age \\
18–29 & 1302 & 14.5 \\
30–44 & 2494 & 27.7 \\
45–59 & 2204 & 24.5 \\
60+ & 3000 & 33.3 \\
Sex & Sex & Sex \\
female & 4747 & 52.7 \\
male & 4253 & 47.3 \\
Race / Ethnicity & Race / Ethnicity & Race / Ethnicity \\
american indian or alaska native & 119 & 1.3 \\
asian & 358 & 4.0 \\
black or african american & 1230 & 13.7 \\
hispanic & 409 & 4.5 \\
other pacific islander & 21 & 0.2 \\
some other race & 80 & 0.9 \\
white & 6783 & 75.4 \\
\end{tblr}

}

\end{table}%

\paragraph{Prompt}\label{prompt}

\begin{lstlisting}
"You are an expert in simulating participants in behavioral and social science studies. You will be given a detailed profile of a specific U.S. resident and asked to respond to survey questions exactly as that person would.\n\n### Participant Profile\n- Age: [AGE] (current year: 2026)\n- Sex: [SEX]\n- Race: [RACECEN1]\n- Highest educational degree: [DEGREE]\n- Household income (GSS bracket): [INCOME16]\n- Household size: [HOMPOP]\n- Self-described social class: [CLASS]\n- Census region: [REGION]\n- Living area: [SRCBELT]\n- Party identification: [PARTYID]\n- Political views (liberal-conservative): [POLVIEWS]\n- Religion: [RELIG]\n- Strength of religious affiliation: [RELITEN]\n- Identify as a born-again or Evangelical Christian: [REBORN]\n- Trust in the scientific community: [CONSCI]\n\nGive the answer this specific person would actually give, even if it is blunt, contradictory, or reflects limited information."
\end{lstlisting}

\paragraph{Simulation}\label{simulation}

The simulation uses the Qualtrics file of the human study to walk the
each LLM through the study as a participant would experience it, with
minor deviations. For details, see
\url{github.com/yarakyrychenko/llm-participants}.

\paragraph{Data}\label{data-1}

This preregistration demonstrates the full analysis workflow using data
from a single model. At the time of preregistration, complete
simulations exist for two models: OpenAI GPT-4o mini and Google Gemini
2.5 Flash Lite. The data for these two models is registered on Zenodo
before data collection (\url{https://doi.org/10.5281/zenodo.20167059}).

Due to API request limitations, simulation data from other models are
not available at the point of registration. These simulations need to be
be conducted in smaller batches. However, the code for running these
simulations is preregistered. For a time-stamped release of the code,
see
\url{https://github.com/yarakyrychenko/llm-participants/releases/tag/v1.0.0}.

\subsection{Measures}\label{measures}

\begin{Shaded}
\begin{Highlighting}[]
\NormalTok{outcomes\_primary }\OtherTok{\textless{}{-}} \StringTok{"trust\_multidimensional"}

\NormalTok{outcomes\_secondary }\OtherTok{\textless{}{-}} \FunctionTok{c}\NormalTok{(}
  \StringTok{"trust\_post"}\NormalTok{, }\StringTok{"distrust\_post"}\NormalTok{, }
  \CommentTok{\#"donation\_ams",}
  \StringTok{"funding\_perceptions"}\NormalTok{, }\StringTok{"policy\_role\_mean"}\NormalTok{, }\StringTok{"inst\_trust\_mean"}
\NormalTok{)}

\NormalTok{outcomes\_binary }\OtherTok{\textless{}{-}} \StringTok{"newsletter\_signup"}

\NormalTok{outcomes\_tertiary }\OtherTok{\textless{}{-}} \FunctionTok{c}\NormalTok{(}
  \StringTok{"belief\_post"}\NormalTok{, }\StringTok{"concern\_mean"}\NormalTok{, }\StringTok{"policy\_general"}\NormalTok{,}
  \StringTok{"policy\_specific\_mean"}\NormalTok{, }\StringTok{"behavior\_mean"}
\NormalTok{)}

\NormalTok{outcomes\_continuous }\OtherTok{\textless{}{-}} \FunctionTok{c}\NormalTok{(outcomes\_primary, outcomes\_secondary, outcomes\_tertiary)}

\CommentTok{\# Illustrative subset for by{-}outcome breakdowns (distributions, moderator detail)}
\NormalTok{outcomes\_illustrative }\OtherTok{\textless{}{-}} \FunctionTok{c}\NormalTok{(}
  \StringTok{"trust\_multidimensional"}\NormalTok{,}
  \StringTok{"donation\_ams"}\NormalTok{,}
  \StringTok{"funding\_perceptions"}\NormalTok{,}
  \StringTok{"policy\_general"}
\NormalTok{)}

\NormalTok{trust\_dimensions }\OtherTok{\textless{}{-}} \FunctionTok{c}\NormalTok{(}
  \StringTok{"trust\_competence"}\NormalTok{, }\StringTok{"trust\_integrity"}\NormalTok{, }\StringTok{"trust\_benevolence"}\NormalTok{, }\StringTok{"trust\_openness"}
\NormalTok{)}

\NormalTok{trust\_items }\OtherTok{\textless{}{-}} \FunctionTok{c}\NormalTok{(}
  \FunctionTok{paste0}\NormalTok{(}\StringTok{"trust\_competence\_"}\NormalTok{,  }\DecValTok{1}\SpecialCharTok{:}\DecValTok{3}\NormalTok{),}
  \FunctionTok{paste0}\NormalTok{(}\StringTok{"trust\_integrity\_"}\NormalTok{,   }\DecValTok{1}\SpecialCharTok{:}\DecValTok{3}\NormalTok{),}
  \FunctionTok{paste0}\NormalTok{(}\StringTok{"trust\_benevolence\_"}\NormalTok{, }\DecValTok{1}\SpecialCharTok{:}\DecValTok{3}\NormalTok{),}
  \FunctionTok{paste0}\NormalTok{(}\StringTok{"trust\_openness\_"}\NormalTok{,    }\DecValTok{1}\SpecialCharTok{:}\DecValTok{3}\NormalTok{)}
\NormalTok{)}

\NormalTok{moderators\_cat }\OtherTok{\textless{}{-}} \FunctionTok{c}\NormalTok{(}\StringTok{"gender"}\NormalTok{, }\StringTok{"age\_band"}\NormalTok{, }\StringTok{"race"}\NormalTok{, }\StringTok{"education"}\NormalTok{, }\StringTok{"income"}\NormalTok{, }\StringTok{"party"}\NormalTok{)}

\NormalTok{outcome\_label\_map }\OtherTok{\textless{}{-}} \FunctionTok{c}\NormalTok{(}
  \AttributeTok{trust\_multidimensional =} \StringTok{"Multidimensional trust (primary)"}\NormalTok{,}
  \AttributeTok{trust\_post             =} \StringTok{"Overall trust"}\NormalTok{,}
  \AttributeTok{distrust\_post          =} \StringTok{"Distrust"}\NormalTok{,}
  \AttributeTok{donation\_ams           =} \StringTok{"Donation to AMS ($)"}\NormalTok{,}
  \AttributeTok{newsletter\_signup      =} \StringTok{"Newsletter sign{-}up"}\NormalTok{,}
  \AttributeTok{funding\_perceptions    =} \StringTok{"Funding perceptions"}\NormalTok{,}
  \AttributeTok{policy\_role\_mean       =} \StringTok{"Policy role"}\NormalTok{,}
  \AttributeTok{inst\_trust\_mean        =} \StringTok{"Institutional trust"}\NormalTok{,}
  \AttributeTok{belief\_post            =} \StringTok{"Climate belief"}\NormalTok{,}
  \AttributeTok{concern\_mean           =} \StringTok{"Climate concern"}\NormalTok{,}
  \AttributeTok{policy\_general         =} \StringTok{"General climate policy"}\NormalTok{,}
  \AttributeTok{policy\_specific\_mean   =} \StringTok{"Specific climate policy"}\NormalTok{,}
  \AttributeTok{behavior\_mean          =} \StringTok{"Pro{-}climate behavior"}
\NormalTok{)}
\end{Highlighting}
\end{Shaded}

\begin{tcolorbox}[enhanced jigsaw, left=2mm, rightrule=.15mm, titlerule=0mm, colframe=quarto-callout-note-color-frame, opacitybacktitle=0.6, colback=white, opacityback=0, breakable, bottomtitle=1mm, coltitle=black, leftrule=.75mm, arc=.35mm, toprule=.15mm, title=\textcolor{quarto-callout-note-color}{\faInfo}\hspace{0.5em}{Note}, colbacktitle=quarto-callout-note-color!10!white, bottomrule=.15mm, toptitle=1mm]

\textbf{Illustrative outcomes.} Analyses that break down results by
outcome --- response distributions (Section 1) and subgroup
heterogeneity (Section 2) --- report results for an illustrative subset
of four outcomes spanning the primary, secondary, and tertiary tiers:
multidimensional trust (primary), donation to AMS, funding perceptions,
and general climate policy. The preregistered ATE comparison (Section 1,
ATE Recovery) covers all outcomes.

\end{tcolorbox}

\subsection{Comparison Functions}\label{comparison-functions}

Two helper functions operationalize the comparison metrics. They are
applied symmetrically to human and LLM clone model outputs throughout
the document.

\subsubsection{Comparison metrics}\label{comparison-metrics}

\texttt{compare\_estimates()} is used throughout --- for ATEs (RQ1),
moderator interaction estimates (RQ3), and secondary analyses. It
computes six metrics, reported in increasing order of strictness. The
rationale for this ordering is to locate precisely \emph{where} clone
performance breaks down: a clone dataset might pass the weakest test
(getting directions right) while failing a stricter one (magnitude
equivalence), and knowing at which rung the ladder breaks is more
informative than any single pass/fail verdict.

\begin{itemize}
\tightlist
\item
  \textbf{Directional agreement (\%)} --- \emph{least strict.} \% of
  estimates with the same sign as the human estimate. Chance performance
  is 50\%. Answers: do clones at least identify the direction of each
  effect?
\item
  \textbf{Spearman ρ} --- rank correlation of point estimates. Answers:
  do clones rank interventions in the same order as humans? Insensitive
  to scale or offset; a clone that inflates all effects equally would
  still score ρ = 1.
\item
  \textbf{Pearson r} --- linear correlation of point estimates. Captures
  proportionality between clone and human estimates but not absolute
  level.
\item
  \textbf{Inferential agreement (\%)} --- \% of estimates with the same
  BH-adjusted conclusion (significantly positive, significantly
  negative, or non-significant at α = .05). Requires getting both
  magnitude and uncertainty right; answers whether a researcher using
  clones would reach the same inferential decision as one using human
  data.
\item
  \textbf{RMSE} --- root mean squared error in outcome units. Captures
  average absolute magnitude error; scale-dependent and therefore not
  comparable across outcomes.
\item
  \textbf{TOST (\%)} --- \emph{most strict.} \% of individual estimate
  pairs passing an equivalence test with bound Δ\_k = 0.5 ×
  \textbar ATE\_human\_k\textbar. Equivalence is declared (two one-sided
  tests, α = .05) when the clone estimate falls within ±50\% of the
  human estimate. Note: near-zero human ATEs shrink Δ\_k proportionally,
  making equivalence harder to achieve.
\end{itemize}

\begin{Shaded}
\begin{Highlighting}[]
\FunctionTok{print\_a\_function\_from\_file}\NormalTok{(}\StringTok{"compare\_estimates"}\NormalTok{)}
\end{Highlighting}
\end{Shaded}

\begin{Shaded}
\begin{Highlighting}[]
\NormalTok{compare\_estimates }\OtherTok{\textless{}{-}} \ControlFlowTok{function}\NormalTok{(h\_result, l\_result,}
                              \AttributeTok{join\_by =} \StringTok{"condition"}\NormalTok{,}
                              \AttributeTok{tost\_delta\_pct =} \FloatTok{0.5}\NormalTok{) \{}

  \CommentTok{\# Standardize columns; SE derived from std.error or from 95\% CI width if absent}
\NormalTok{  prep }\OtherTok{\textless{}{-}} \ControlFlowTok{function}\NormalTok{(df, suffix) \{}
\NormalTok{    df }\SpecialCharTok{|\textgreater{}}
      \FunctionTok{select}\NormalTok{(}\FunctionTok{all\_of}\NormalTok{(join\_by), estimate) }\SpecialCharTok{|\textgreater{}}
      \FunctionTok{mutate}\NormalTok{(}
        \AttributeTok{se    =} \ControlFlowTok{if}\NormalTok{ (}\StringTok{"std.error"} \SpecialCharTok{\%in\%} \FunctionTok{names}\NormalTok{(df)) df}\SpecialCharTok{$}\NormalTok{std.error}
                \ControlFlowTok{else} \ControlFlowTok{if}\NormalTok{ (}\FunctionTok{all}\NormalTok{(}\FunctionTok{c}\NormalTok{(}\StringTok{"conf.low"}\NormalTok{, }\StringTok{"conf.high"}\NormalTok{) }\SpecialCharTok{\%in\%} \FunctionTok{names}\NormalTok{(df)))}
\NormalTok{                  (df}\SpecialCharTok{$}\NormalTok{conf.high }\SpecialCharTok{{-}}\NormalTok{ df}\SpecialCharTok{$}\NormalTok{conf.low) }\SpecialCharTok{/}\NormalTok{ (}\DecValTok{2} \SpecialCharTok{*} \FloatTok{1.96}\NormalTok{)}
                \ControlFlowTok{else} \ConstantTok{NA\_real\_}\NormalTok{,}
        \AttributeTok{p\_adj =} \ControlFlowTok{if}\NormalTok{ (}\StringTok{"p.value\_adjusted"} \SpecialCharTok{\%in\%} \FunctionTok{names}\NormalTok{(df)) df}\SpecialCharTok{$}\NormalTok{p.value\_adjusted}
                \ControlFlowTok{else} \ConstantTok{NA\_real\_}
\NormalTok{      ) }\SpecialCharTok{|\textgreater{}}
      \FunctionTok{rename\_with}\NormalTok{(}\SpecialCharTok{\textasciitilde{}} \FunctionTok{paste0}\NormalTok{(.x, }\StringTok{"\_"}\NormalTok{, suffix), }\FunctionTok{c}\NormalTok{(estimate, se, p\_adj))}
\NormalTok{  \}}

\NormalTok{  joined }\OtherTok{\textless{}{-}} \FunctionTok{inner\_join}\NormalTok{(}\FunctionTok{prep}\NormalTok{(h\_result, }\StringTok{"h"}\NormalTok{), }\FunctionTok{prep}\NormalTok{(l\_result, }\StringTok{"l"}\NormalTok{), }\AttributeTok{by =}\NormalTok{ join\_by)}

  \CommentTok{\# Three{-}level inferential category based on BH{-}adjusted p and direction}
\NormalTok{  infer\_cat }\OtherTok{\textless{}{-}} \ControlFlowTok{function}\NormalTok{(est, p\_adj) \{}
    \FunctionTok{case\_when}\NormalTok{(}
\NormalTok{      p\_adj }\SpecialCharTok{\textless{}} \FloatTok{0.05} \SpecialCharTok{\&}\NormalTok{ est }\SpecialCharTok{\textgreater{}} \DecValTok{0} \SpecialCharTok{\textasciitilde{}} \StringTok{"sig\_positive"}\NormalTok{,}
\NormalTok{      p\_adj }\SpecialCharTok{\textless{}} \FloatTok{0.05} \SpecialCharTok{\&}\NormalTok{ est }\SpecialCharTok{\textless{}} \DecValTok{0} \SpecialCharTok{\textasciitilde{}} \StringTok{"sig\_negative"}\NormalTok{,}
      \ConstantTok{TRUE}                   \SpecialCharTok{\textasciitilde{}} \StringTok{"not\_significant"}
\NormalTok{    )}
\NormalTok{  \}}

  \CommentTok{\# TOST: equivalence declared when |ATE\_h {-} ATE\_l| \textless{} delta = tost\_delta\_pct × |ATE\_h|}
  \CommentTok{\# p\_tost is the maximum of the two one{-}sided p{-}values; equivalent if p\_tost \textless{} 0.05}
\NormalTok{  joined }\SpecialCharTok{|\textgreater{}}
    \FunctionTok{mutate}\NormalTok{(}
      \AttributeTok{delta      =}\NormalTok{ tost\_delta\_pct }\SpecialCharTok{*} \FunctionTok{abs}\NormalTok{(estimate\_h),}
      \AttributeTok{diff       =}\NormalTok{ estimate\_h }\SpecialCharTok{{-}}\NormalTok{ estimate\_l,}
      \AttributeTok{se\_diff    =} \FunctionTok{sqrt}\NormalTok{(se\_h}\SpecialCharTok{\^{}}\DecValTok{2} \SpecialCharTok{+}\NormalTok{ se\_l}\SpecialCharTok{\^{}}\DecValTok{2}\NormalTok{),}
      \AttributeTok{p\_lower    =} \FunctionTok{pnorm}\NormalTok{((diff }\SpecialCharTok{+}\NormalTok{ delta) }\SpecialCharTok{/}\NormalTok{ se\_diff, }\AttributeTok{lower.tail =} \ConstantTok{FALSE}\NormalTok{),}
      \AttributeTok{p\_upper    =} \FunctionTok{pnorm}\NormalTok{((diff }\SpecialCharTok{{-}}\NormalTok{ delta) }\SpecialCharTok{/}\NormalTok{ se\_diff),}
      \AttributeTok{p\_tost     =} \FunctionTok{pmax}\NormalTok{(p\_lower, p\_upper),}
      \AttributeTok{equivalent =}\NormalTok{ p\_tost }\SpecialCharTok{\textless{}} \FloatTok{0.05}\NormalTok{,}
      \AttributeTok{same\_infer =} \FunctionTok{infer\_cat}\NormalTok{(estimate\_h, p\_adj\_h) }\SpecialCharTok{==} \FunctionTok{infer\_cat}\NormalTok{(estimate\_l, p\_adj\_l)}
\NormalTok{    ) }\SpecialCharTok{|\textgreater{}}
    \FunctionTok{summarise}\NormalTok{(}
      \AttributeTok{spearman\_rho    =} \FunctionTok{cor}\NormalTok{(estimate\_h, estimate\_l, }\AttributeTok{method =} \StringTok{"spearman"}\NormalTok{,}
                            \AttributeTok{use =} \StringTok{"pairwise.complete.obs"}\NormalTok{),}
      \AttributeTok{pearson\_r       =} \FunctionTok{cor}\NormalTok{(estimate\_h, estimate\_l, }\AttributeTok{use =} \StringTok{"pairwise.complete.obs"}\NormalTok{),}
      \AttributeTok{rmse            =} \FunctionTok{sqrt}\NormalTok{(}\FunctionTok{mean}\NormalTok{((estimate\_h }\SpecialCharTok{{-}}\NormalTok{ estimate\_l)}\SpecialCharTok{\^{}}\DecValTok{2}\NormalTok{, }\AttributeTok{na.rm =} \ConstantTok{TRUE}\NormalTok{)),}
      \AttributeTok{directional\_pct =} \FunctionTok{mean}\NormalTok{(}\FunctionTok{sign}\NormalTok{(estimate\_h) }\SpecialCharTok{==} \FunctionTok{sign}\NormalTok{(estimate\_l), }\AttributeTok{na.rm =} \ConstantTok{TRUE}\NormalTok{) }\SpecialCharTok{*} \DecValTok{100}\NormalTok{,}
      \AttributeTok{inferential\_pct =} \ControlFlowTok{if}\NormalTok{ (}\SpecialCharTok{!}\FunctionTok{anyNA}\NormalTok{(p\_adj\_h) }\SpecialCharTok{\&\&} \SpecialCharTok{!}\FunctionTok{anyNA}\NormalTok{(p\_adj\_l))}
        \FunctionTok{mean}\NormalTok{(same\_infer, }\AttributeTok{na.rm =} \ConstantTok{TRUE}\NormalTok{) }\SpecialCharTok{*} \DecValTok{100} \ControlFlowTok{else} \ConstantTok{NA\_real\_}\NormalTok{,}
      \AttributeTok{tost\_pct        =} \FunctionTok{mean}\NormalTok{(equivalent, }\AttributeTok{na.rm =} \ConstantTok{TRUE}\NormalTok{) }\SpecialCharTok{*} \DecValTok{100}
\NormalTok{    )}
\NormalTok{\}}
\end{Highlighting}
\end{Shaded}

\subsubsection{Distribution comparison
metrics}\label{distribution-comparison-metrics}

\texttt{compare\_distributions()} compares response distributions for a
single condition and outcome using three metrics, again in increasing
order of strictness:

\begin{itemize}
\tightlist
\item
  \textbf{Variance ratio} (clone / human variance) --- \emph{least
  strict.} Captures only one aspect of distributional similarity ---
  relative spread. A ratio of 1 means equal dispersion; does not assess
  shape.
\item
  \textbf{OVL (overlapping coefficient)} --- proportion of the total
  area shared between the two kernel density estimates; 1 = identical
  distributions, 0 = no overlap. Summarises overall shape similarity.
\item
  \textbf{KS D-statistic} --- \emph{most strict.} Maximum absolute gap
  between the two empirical CDFs at any point on the response scale; 0 =
  indistinguishable CDFs everywhere. Sensitive to local departures that
  OVL might average away.
\end{itemize}

\begin{Shaded}
\begin{Highlighting}[]
\FunctionTok{print\_a\_function\_from\_file}\NormalTok{(}\StringTok{"compute\_ovl"}\NormalTok{)}
\end{Highlighting}
\end{Shaded}

\begin{Shaded}
\begin{Highlighting}[]
\NormalTok{compute\_ovl }\OtherTok{\textless{}{-}} \ControlFlowTok{function}\NormalTok{(x, y, }\AttributeTok{n\_grid =} \DecValTok{512}\NormalTok{) \{}
\NormalTok{  lo }\OtherTok{\textless{}{-}} \FunctionTok{min}\NormalTok{(}\FunctionTok{c}\NormalTok{(x, y)); hi }\OtherTok{\textless{}{-}} \FunctionTok{max}\NormalTok{(}\FunctionTok{c}\NormalTok{(x, y))}
  \ControlFlowTok{if}\NormalTok{ (lo }\SpecialCharTok{==}\NormalTok{ hi) }\FunctionTok{return}\NormalTok{(}\ConstantTok{NA\_real\_}\NormalTok{)}
  \CommentTok{\# Evaluate both KDEs on the same grid, then integrate the overlap area}
\NormalTok{  d\_h }\OtherTok{\textless{}{-}} \FunctionTok{density}\NormalTok{(x, }\AttributeTok{from =}\NormalTok{ lo, }\AttributeTok{to =}\NormalTok{ hi, }\AttributeTok{n =}\NormalTok{ n\_grid)}
\NormalTok{  d\_l }\OtherTok{\textless{}{-}} \FunctionTok{density}\NormalTok{(y, }\AttributeTok{from =}\NormalTok{ lo, }\AttributeTok{to =}\NormalTok{ hi, }\AttributeTok{n =}\NormalTok{ n\_grid)}
  \FunctionTok{sum}\NormalTok{(}\FunctionTok{pmin}\NormalTok{(d\_h}\SpecialCharTok{$}\NormalTok{y, d\_l}\SpecialCharTok{$}\NormalTok{y)) }\SpecialCharTok{*}\NormalTok{ (hi }\SpecialCharTok{{-}}\NormalTok{ lo) }\SpecialCharTok{/}\NormalTok{ n\_grid}
\NormalTok{\}}
\end{Highlighting}
\end{Shaded}

\subsubsection{Calibration regression}\label{calibration-regression}

Correlation between human and clone ATEs tells us only whether the two
effect-size vectors \emph{covary} --- whether interventions that produce
large effects in humans also produce large effects in clones. However,
correlation does not capture baseline differences or multiplicative
scaling: clones could underestimate the control-condition baseline and
inflate every effect by a factor of two while still yielding Pearson r =
1. Calibration regression separates these two failure modes --- additive
bias and proportionality.

\texttt{run\_calibration()} regresses human ATEs on LLM clone ATEs:
\texttt{ATE\_human\ =\ α\ +\ β\ ×\ ATE\_clone}. The intercept α captures
additive bias: α \textgreater{} 0 means clones systematically
underestimate human ATEs (an offset must be added to recover them), α
\textless{} 0 means they overestimate. The slope β captures
proportionality: β = 1 is perfect scaling (clone effects move 1:1 with
human effects), β \textgreater{} 1 means clone effects are compressed
relative to humans (too small in absolute magnitude), β \textless{} 1
means they are exaggerated. We report 95\% confidence intervals for both
parameters and R² as the share of human-ATE variance linearly
recoverable from clone ATEs.

We then test calibration at two levels of strictness. The \textbf{joint
F-test of (α = 0, β = 1)} is the omnibus test of perfect calibration: it
asks whether clone ATEs can be read directly as estimates of human ATEs,
with no correction at all. A significant result means the null of
perfect calibration is rejected --- clone ATEs are not a direct
substitute for human ATEs. The \textbf{slope-only F-test of (β = 1)}
asks the weaker diagnostical question whether clone effects are
\emph{proportionally} calibrated, i.e.~whether a one-unit difference in
a clone ATE corresponds to a one-unit difference in a human ATE, while
permitting a constant additive offset α. This additional test is
meaningful because proportional calibration is the more practically
relevant criterion: in real-world applications, additive bias is
relatively easy to correct with a single human anchor condition, whereas
a slope distortion reflects a deeper failure in how clones translate to
humans and cannot be removed without re-estimating β.

\begin{Shaded}
\begin{Highlighting}[]
\FunctionTok{print\_a\_function\_from\_file}\NormalTok{(}\StringTok{"run\_calibration"}\NormalTok{)}
\end{Highlighting}
\end{Shaded}

\begin{Shaded}
\begin{Highlighting}[]
\NormalTok{run\_calibration }\OtherTok{\textless{}{-}} \ControlFlowTok{function}\NormalTok{(h\_result, l\_result) \{}
\NormalTok{  joined }\OtherTok{\textless{}{-}} \FunctionTok{inner\_join}\NormalTok{(}
\NormalTok{    h\_result }\SpecialCharTok{|\textgreater{}} \FunctionTok{select}\NormalTok{(condition, }\AttributeTok{estimate\_h =}\NormalTok{ estimate),}
\NormalTok{    l\_result }\SpecialCharTok{|\textgreater{}} \FunctionTok{select}\NormalTok{(condition, }\AttributeTok{estimate\_l =}\NormalTok{ estimate),}
    \AttributeTok{by =} \StringTok{"condition"}
\NormalTok{  )}

  \CommentTok{\# ATE\_h = α + β × ATE\_l}
\NormalTok{  fit   }\OtherTok{\textless{}{-}} \FunctionTok{lm}\NormalTok{(estimate\_h }\SpecialCharTok{\textasciitilde{}}\NormalTok{ estimate\_l, }\AttributeTok{data =}\NormalTok{ joined)}
\NormalTok{  coefs }\OtherTok{\textless{}{-}}\NormalTok{ broom}\SpecialCharTok{::}\FunctionTok{tidy}\NormalTok{(fit, }\AttributeTok{conf.int =} \ConstantTok{TRUE}\NormalTok{)}

  \FunctionTok{tibble}\NormalTok{(}
    \AttributeTok{alpha     =}\NormalTok{ coefs}\SpecialCharTok{$}\NormalTok{estimate[}\DecValTok{1}\NormalTok{],}
    \AttributeTok{alpha\_lo  =}\NormalTok{ coefs}\SpecialCharTok{$}\NormalTok{conf.low[}\DecValTok{1}\NormalTok{],}
    \AttributeTok{alpha\_hi  =}\NormalTok{ coefs}\SpecialCharTok{$}\NormalTok{conf.high[}\DecValTok{1}\NormalTok{],}
    \AttributeTok{beta      =}\NormalTok{ coefs}\SpecialCharTok{$}\NormalTok{estimate[}\DecValTok{2}\NormalTok{],}
    \AttributeTok{beta\_lo   =}\NormalTok{ coefs}\SpecialCharTok{$}\NormalTok{conf.low[}\DecValTok{2}\NormalTok{],}
    \AttributeTok{beta\_hi   =}\NormalTok{ coefs}\SpecialCharTok{$}\NormalTok{conf.high[}\DecValTok{2}\NormalTok{],}
    \AttributeTok{r\_squared =}\NormalTok{ broom}\SpecialCharTok{::}\FunctionTok{glance}\NormalTok{(fit)}\SpecialCharTok{$}\NormalTok{r.squared,}
    \CommentTok{\# Joint F{-}test  H0: α = 0 AND β = 1  (perfect calibration on the original scale)}
    \AttributeTok{p\_joint   =}\NormalTok{ car}\SpecialCharTok{::}\FunctionTok{linearHypothesis}\NormalTok{(}
\NormalTok{                  fit, }\FunctionTok{c}\NormalTok{(}\StringTok{"(Intercept) = 0"}\NormalTok{, }\StringTok{"estimate\_l = 1"}\NormalTok{)}
\NormalTok{                )}\SpecialCharTok{$}\StringTok{\textasciigrave{}}\AttributeTok{Pr(\textgreater{}F)}\StringTok{\textasciigrave{}}\NormalTok{[}\DecValTok{2}\NormalTok{],}
    \CommentTok{\# Slope{-}only F{-}test  H0: β = 1  (proportional calibration; constant offset allowed)}
    \AttributeTok{p\_beta\_1  =}\NormalTok{ car}\SpecialCharTok{::}\FunctionTok{linearHypothesis}\NormalTok{(fit, }\StringTok{"estimate\_l = 1"}\NormalTok{)}\SpecialCharTok{$}\StringTok{\textasciigrave{}}\AttributeTok{Pr(\textgreater{}F)}\StringTok{\textasciigrave{}}\NormalTok{[}\DecValTok{2}\NormalTok{]}
\NormalTok{  )}
\NormalTok{\}}
\end{Highlighting}
\end{Shaded}

\subsubsection{Demographic baseline and predictability
metrics}\label{demographic-baseline-and-predictability-metrics}

\texttt{compare\_demographic\_baselines()} computes, for each moderator,
the Pearson correlation and RMSE between human and clone cell means in
the control condition. \texttt{compare\_demographic\_predictability()}
fits separate OLS regressions of the outcome on each moderator plus
condition fixed effects, returning R² and dummy-coded coefficients for
humans and clones.

\begin{Shaded}
\begin{Highlighting}[]
\FunctionTok{print\_a\_function\_from\_file}\NormalTok{(}\StringTok{"compare\_demographic\_baselines"}\NormalTok{)}
\end{Highlighting}
\end{Shaded}

\begin{Shaded}
\begin{Highlighting}[]
\NormalTok{compare\_demographic\_baselines }\OtherTok{\textless{}{-}} \ControlFlowTok{function}\NormalTok{(human\_data, llm\_data,}
\NormalTok{                                           outcome,}
\NormalTok{                                           moderators,}
                                           \AttributeTok{condition\_var =} \StringTok{"condition"}\NormalTok{) \{}
\NormalTok{  control\_val }\OtherTok{\textless{}{-}} \FunctionTok{levels}\NormalTok{(human\_data[[condition\_var]])[}\DecValTok{1}\NormalTok{]}

  \FunctionTok{map\_dfr}\NormalTok{(moderators, }\ControlFlowTok{function}\NormalTok{(mod) \{}
\NormalTok{    h\_cells }\OtherTok{\textless{}{-}}\NormalTok{ human\_data }\SpecialCharTok{|\textgreater{}}
      \FunctionTok{filter}\NormalTok{(.data[[condition\_var]] }\SpecialCharTok{==}\NormalTok{ control\_val) }\SpecialCharTok{|\textgreater{}}
      \FunctionTok{group\_by}\NormalTok{(}\AttributeTok{cell =}\NormalTok{ .data[[mod]]) }\SpecialCharTok{|\textgreater{}}
      \FunctionTok{summarise}\NormalTok{(}\AttributeTok{mean\_h =} \FunctionTok{mean}\NormalTok{(.data[[outcome]], }\AttributeTok{na.rm =} \ConstantTok{TRUE}\NormalTok{), }\AttributeTok{.groups =} \StringTok{"drop"}\NormalTok{)}

\NormalTok{    l\_cells }\OtherTok{\textless{}{-}}\NormalTok{ llm\_data }\SpecialCharTok{|\textgreater{}}
      \FunctionTok{filter}\NormalTok{(.data[[condition\_var]] }\SpecialCharTok{==}\NormalTok{ control\_val) }\SpecialCharTok{|\textgreater{}}
      \FunctionTok{group\_by}\NormalTok{(}\AttributeTok{cell =}\NormalTok{ .data[[mod]]) }\SpecialCharTok{|\textgreater{}}
      \FunctionTok{summarise}\NormalTok{(}\AttributeTok{mean\_l =} \FunctionTok{mean}\NormalTok{(.data[[outcome]], }\AttributeTok{na.rm =} \ConstantTok{TRUE}\NormalTok{), }\AttributeTok{.groups =} \StringTok{"drop"}\NormalTok{)}

    \FunctionTok{inner\_join}\NormalTok{(h\_cells, l\_cells, }\AttributeTok{by =} \StringTok{"cell"}\NormalTok{) }\SpecialCharTok{|\textgreater{}}
      \FunctionTok{summarise}\NormalTok{(}
        \AttributeTok{moderator =}\NormalTok{ mod,}
        \AttributeTok{r         =} \FunctionTok{cor}\NormalTok{(mean\_h, mean\_l, }\AttributeTok{use =} \StringTok{"pairwise.complete.obs"}\NormalTok{),}
        \AttributeTok{rmse      =} \FunctionTok{sqrt}\NormalTok{(}\FunctionTok{mean}\NormalTok{((mean\_h }\SpecialCharTok{{-}}\NormalTok{ mean\_l)}\SpecialCharTok{\^{}}\DecValTok{2}\NormalTok{, }\AttributeTok{na.rm =} \ConstantTok{TRUE}\NormalTok{)),}
        \AttributeTok{n\_cells   =} \FunctionTok{n}\NormalTok{()}
\NormalTok{      )}
\NormalTok{  \})}
\NormalTok{\}}
\end{Highlighting}
\end{Shaded}

\begin{Shaded}
\begin{Highlighting}[]
\FunctionTok{print\_a\_function\_from\_file}\NormalTok{(}\StringTok{"compare\_demographic\_predictability"}\NormalTok{)}
\end{Highlighting}
\end{Shaded}

\begin{Shaded}
\begin{Highlighting}[]
\NormalTok{compare\_demographic\_predictability }\OtherTok{\textless{}{-}} \ControlFlowTok{function}\NormalTok{(human\_data, llm\_data,}
\NormalTok{                                                outcome,}
\NormalTok{                                                predictors,}
                                                \AttributeTok{condition\_var =} \StringTok{"condition"}\NormalTok{) \{}
\NormalTok{  results }\OtherTok{\textless{}{-}} \FunctionTok{map}\NormalTok{(predictors, }\ControlFlowTok{function}\NormalTok{(mod) \{}
\NormalTok{    formula }\OtherTok{\textless{}{-}} \FunctionTok{as.formula}\NormalTok{(}\FunctionTok{paste}\NormalTok{(outcome, }\StringTok{"\textasciitilde{}"}\NormalTok{, mod, }\StringTok{"+"}\NormalTok{, condition\_var))}
\NormalTok{    fit\_h   }\OtherTok{\textless{}{-}} \FunctionTok{lm}\NormalTok{(formula, }\AttributeTok{data =}\NormalTok{ human\_data)}
\NormalTok{    fit\_l   }\OtherTok{\textless{}{-}} \FunctionTok{lm}\NormalTok{(formula, }\AttributeTok{data =}\NormalTok{ llm\_data)}

\NormalTok{    rsq }\OtherTok{\textless{}{-}} \FunctionTok{tibble}\NormalTok{(}
      \AttributeTok{moderator =}\NormalTok{ mod,}
      \AttributeTok{source    =} \FunctionTok{c}\NormalTok{(}\StringTok{"human"}\NormalTok{, }\StringTok{"llm"}\NormalTok{),}
      \AttributeTok{r\_squared =} \FunctionTok{c}\NormalTok{(broom}\SpecialCharTok{::}\FunctionTok{glance}\NormalTok{(fit\_h)}\SpecialCharTok{$}\NormalTok{r.squared, broom}\SpecialCharTok{::}\FunctionTok{glance}\NormalTok{(fit\_l)}\SpecialCharTok{$}\NormalTok{r.squared)}
\NormalTok{    )}

\NormalTok{    coefs\_h }\OtherTok{\textless{}{-}}\NormalTok{ broom}\SpecialCharTok{::}\FunctionTok{tidy}\NormalTok{(fit\_h, }\AttributeTok{conf.int =} \ConstantTok{TRUE}\NormalTok{) }\SpecialCharTok{|\textgreater{}}
      \FunctionTok{filter}\NormalTok{(}\FunctionTok{str\_detect}\NormalTok{(term, }\FunctionTok{fixed}\NormalTok{(mod))) }\SpecialCharTok{|\textgreater{}}
      \FunctionTok{select}\NormalTok{(term, }\AttributeTok{est\_h =}\NormalTok{ estimate, }\AttributeTok{lo\_h =}\NormalTok{ conf.low, }\AttributeTok{hi\_h =}\NormalTok{ conf.high)}

\NormalTok{    coefs\_l }\OtherTok{\textless{}{-}}\NormalTok{ broom}\SpecialCharTok{::}\FunctionTok{tidy}\NormalTok{(fit\_l, }\AttributeTok{conf.int =} \ConstantTok{TRUE}\NormalTok{) }\SpecialCharTok{|\textgreater{}}
      \FunctionTok{filter}\NormalTok{(}\FunctionTok{str\_detect}\NormalTok{(term, }\FunctionTok{fixed}\NormalTok{(mod))) }\SpecialCharTok{|\textgreater{}}
      \FunctionTok{select}\NormalTok{(term, }\AttributeTok{est\_l =}\NormalTok{ estimate, }\AttributeTok{lo\_l =}\NormalTok{ conf.low, }\AttributeTok{hi\_l =}\NormalTok{ conf.high)}

\NormalTok{    coefs }\OtherTok{\textless{}{-}} \FunctionTok{inner\_join}\NormalTok{(coefs\_h, coefs\_l, }\AttributeTok{by =} \StringTok{"term"}\NormalTok{) }\SpecialCharTok{|\textgreater{}}
      \FunctionTok{mutate}\NormalTok{(}\AttributeTok{moderator =}\NormalTok{ mod)}

    \FunctionTok{list}\NormalTok{(}\AttributeTok{r\_squared =}\NormalTok{ rsq, }\AttributeTok{coefficients =}\NormalTok{ coefs)}
\NormalTok{  \})}

  \FunctionTok{list}\NormalTok{(}
    \AttributeTok{r\_squared    =} \FunctionTok{map\_dfr}\NormalTok{(results, }\StringTok{"r\_squared"}\NormalTok{),}
    \AttributeTok{coefficients =} \FunctionTok{map\_dfr}\NormalTok{(results, }\StringTok{"coefficients"}\NormalTok{)}
\NormalTok{  )}
\NormalTok{\}}
\end{Highlighting}
\end{Shaded}

\subsection{1. Average Treatment
Effects}\label{average-treatment-effects}

\subsubsection{ATE Recovery}\label{ate-recovery}

For each intervention, the average treatment effect (ATE) is the
difference in mean outcome between participants assigned to that
intervention and those in the control condition. We estimate ATEs by
running \texttt{run\_main\_treatment\_model()} --- OLS with HC2
heteroskedasticity-robust standard errors and BH-adjusted p-values ---
on human and LLM clone data separately for each preregistered outcome.
For \texttt{newsletter\_signup} (binary), we use
\texttt{run\_main\_treatment\_model\_binary()}, which returns marginal
effects on the probability scale.

We report the six comparison metrics in increasing order of strictness,
from directional agreement (do clones identify the correct direction of
each effect?) to TOST equivalence testing (does each clone ATE fall
within a preregistered tolerance of the corresponding human ATE?). We
first report metrics pooled across all outcomes --- 240 estimate pairs
in total --- to maximize statistical power, then break down by outcome
and zoom in on the primary outcome. All models are fit first.

\begin{Shaded}
\begin{Highlighting}[]
\CommentTok{\# Run and store all models (avoids re{-}fitting in later chunks)}
\NormalTok{ate\_model\_results }\OtherTok{\textless{}{-}} \FunctionTok{map}\NormalTok{(outcomes\_continuous, }\ControlFlowTok{function}\NormalTok{(out) \{}
  \FunctionTok{list}\NormalTok{(}
    \AttributeTok{h  =} \FunctionTok{run\_main\_treatment\_model}\NormalTok{(human\_data\_1, }\AttributeTok{outcome =}\NormalTok{ out),}
    \AttributeTok{l  =} \FunctionTok{run\_main\_treatment\_model}\NormalTok{(llm\_data,     }\AttributeTok{outcome =}\NormalTok{ out),}
    \AttributeTok{h2 =} \FunctionTok{run\_main\_treatment\_model}\NormalTok{(human\_data\_2, }\AttributeTok{outcome =}\NormalTok{ out)}
\NormalTok{  )}
\NormalTok{\}) }\SpecialCharTok{|\textgreater{}} \FunctionTok{set\_names}\NormalTok{(outcomes\_continuous)}

\CommentTok{\# Binary outcome: marginal effects on probability scale}
\NormalTok{h\_bin  }\OtherTok{\textless{}{-}} \FunctionTok{run\_main\_treatment\_model\_binary}\NormalTok{(human\_data\_1, }\StringTok{"newsletter\_signup"}\NormalTok{)}
\NormalTok{l\_bin  }\OtherTok{\textless{}{-}} \FunctionTok{run\_main\_treatment\_model\_binary}\NormalTok{(llm\_data,     }\StringTok{"newsletter\_signup"}\NormalTok{)}
\NormalTok{h2\_bin }\OtherTok{\textless{}{-}} \FunctionTok{run\_main\_treatment\_model\_binary}\NormalTok{(human\_data\_2, }\StringTok{"newsletter\_signup"}\NormalTok{)}

\CommentTok{\# References to primary outcome models for the primary subsection below}
\NormalTok{h\_primary  }\OtherTok{\textless{}{-}}\NormalTok{ ate\_model\_results[[}\StringTok{"trust\_multidimensional"}\NormalTok{]]}\SpecialCharTok{$}\NormalTok{h}
\NormalTok{l\_primary  }\OtherTok{\textless{}{-}}\NormalTok{ ate\_model\_results[[}\StringTok{"trust\_multidimensional"}\NormalTok{]]}\SpecialCharTok{$}\NormalTok{l}
\NormalTok{h2\_primary }\OtherTok{\textless{}{-}}\NormalTok{ ate\_model\_results[[}\StringTok{"trust\_multidimensional"}\NormalTok{]]}\SpecialCharTok{$}\NormalTok{h2}

\CommentTok{\# Extract ATE pairs with SE and adjusted p{-}value from a model result tibble}
\NormalTok{prep\_ate }\OtherTok{\textless{}{-}} \ControlFlowTok{function}\NormalTok{(df) \{}
\NormalTok{  df }\SpecialCharTok{|\textgreater{}}
    \FunctionTok{select}\NormalTok{(condition, estimate) }\SpecialCharTok{|\textgreater{}}
    \FunctionTok{mutate}\NormalTok{(}
      \AttributeTok{se    =} \ControlFlowTok{if}\NormalTok{ (}\StringTok{"std.error"}        \SpecialCharTok{\%in\%} \FunctionTok{names}\NormalTok{(df)) df}\SpecialCharTok{$}\NormalTok{std.error        }\ControlFlowTok{else} \ConstantTok{NA\_real\_}\NormalTok{,}
      \AttributeTok{p\_adj =} \ControlFlowTok{if}\NormalTok{ (}\StringTok{"p.value\_adjusted"} \SpecialCharTok{\%in\%} \FunctionTok{names}\NormalTok{(df)) df}\SpecialCharTok{$}\NormalTok{p.value\_adjusted }\ControlFlowTok{else} \ConstantTok{NA\_real\_}
\NormalTok{    )}
\NormalTok{\}}

\CommentTok{\# Pool ATE pairs across all outcomes (continuous + binary)}
\NormalTok{pooled\_ates }\OtherTok{\textless{}{-}} \FunctionTok{bind\_rows}\NormalTok{(}
  \FunctionTok{map\_dfr}\NormalTok{(outcomes\_continuous, }\ControlFlowTok{function}\NormalTok{(out) \{}
    \FunctionTok{inner\_join}\NormalTok{(}
      \FunctionTok{prep\_ate}\NormalTok{(ate\_model\_results[[out]]}\SpecialCharTok{$}\NormalTok{h) }\SpecialCharTok{|\textgreater{}}
        \FunctionTok{rename}\NormalTok{(}\AttributeTok{estimate\_h =}\NormalTok{ estimate, }\AttributeTok{se\_h =}\NormalTok{ se, }\AttributeTok{p\_adj\_h =}\NormalTok{ p\_adj),}
      \FunctionTok{prep\_ate}\NormalTok{(ate\_model\_results[[out]]}\SpecialCharTok{$}\NormalTok{l) }\SpecialCharTok{|\textgreater{}}
        \FunctionTok{rename}\NormalTok{(}\AttributeTok{estimate\_l =}\NormalTok{ estimate, }\AttributeTok{se\_l =}\NormalTok{ se, }\AttributeTok{p\_adj\_l =}\NormalTok{ p\_adj),}
      \AttributeTok{by =} \StringTok{"condition"}
\NormalTok{    ) }\SpecialCharTok{|\textgreater{}} \FunctionTok{mutate}\NormalTok{(}\AttributeTok{outcome =}\NormalTok{ out)}
\NormalTok{  \}),}
  \FunctionTok{inner\_join}\NormalTok{(}
    \FunctionTok{prep\_ate}\NormalTok{(h\_bin}\SpecialCharTok{$}\NormalTok{marginal\_effects) }\SpecialCharTok{|\textgreater{}}
      \FunctionTok{rename}\NormalTok{(}\AttributeTok{estimate\_h =}\NormalTok{ estimate, }\AttributeTok{se\_h =}\NormalTok{ se, }\AttributeTok{p\_adj\_h =}\NormalTok{ p\_adj),}
    \FunctionTok{prep\_ate}\NormalTok{(l\_bin}\SpecialCharTok{$}\NormalTok{marginal\_effects) }\SpecialCharTok{|\textgreater{}}
      \FunctionTok{rename}\NormalTok{(}\AttributeTok{estimate\_l =}\NormalTok{ estimate, }\AttributeTok{se\_l =}\NormalTok{ se, }\AttributeTok{p\_adj\_l =}\NormalTok{ p\_adj),}
    \AttributeTok{by =} \StringTok{"condition"}
\NormalTok{  ) }\SpecialCharTok{|\textgreater{}} \FunctionTok{mutate}\NormalTok{(}\AttributeTok{outcome =} \StringTok{"newsletter\_signup"}\NormalTok{)}
\NormalTok{) }\SpecialCharTok{|\textgreater{}}
  \FunctionTok{mutate}\NormalTok{(}
    \AttributeTok{tier =} \FunctionTok{case\_when}\NormalTok{(}
\NormalTok{      outcome }\SpecialCharTok{==}\NormalTok{ outcomes\_primary                          }\SpecialCharTok{\textasciitilde{}} \StringTok{"Primary"}\NormalTok{,}
\NormalTok{      outcome }\SpecialCharTok{\%in\%} \FunctionTok{c}\NormalTok{(outcomes\_secondary, outcomes\_binary) }\SpecialCharTok{\textasciitilde{}} \StringTok{"Secondary"}\NormalTok{,}
      \ConstantTok{TRUE}                                                \SpecialCharTok{\textasciitilde{}} \StringTok{"Tertiary"}
\NormalTok{    ),}
    \AttributeTok{outcome\_label =}\NormalTok{ outcome\_label\_map[outcome]}
\NormalTok{  )}

\CommentTok{\# Human 1 vs. Human 2 ATE pairs (same structure as pooled\_ates)}
\NormalTok{pooled\_ates\_h2 }\OtherTok{\textless{}{-}} \FunctionTok{bind\_rows}\NormalTok{(}
  \FunctionTok{map\_dfr}\NormalTok{(outcomes\_continuous, }\ControlFlowTok{function}\NormalTok{(out) \{}
    \FunctionTok{inner\_join}\NormalTok{(}
      \FunctionTok{prep\_ate}\NormalTok{(ate\_model\_results[[out]]}\SpecialCharTok{$}\NormalTok{h) }\SpecialCharTok{|\textgreater{}}
        \FunctionTok{rename}\NormalTok{(}\AttributeTok{estimate\_h =}\NormalTok{ estimate, }\AttributeTok{se\_h =}\NormalTok{ se, }\AttributeTok{p\_adj\_h =}\NormalTok{ p\_adj),}
      \FunctionTok{prep\_ate}\NormalTok{(ate\_model\_results[[out]]}\SpecialCharTok{$}\NormalTok{h2) }\SpecialCharTok{|\textgreater{}}
        \FunctionTok{rename}\NormalTok{(}\AttributeTok{estimate\_h2 =}\NormalTok{ estimate, }\AttributeTok{se\_h2 =}\NormalTok{ se, }\AttributeTok{p\_adj\_h2 =}\NormalTok{ p\_adj),}
      \AttributeTok{by =} \StringTok{"condition"}
\NormalTok{    ) }\SpecialCharTok{|\textgreater{}} \FunctionTok{mutate}\NormalTok{(}\AttributeTok{outcome =}\NormalTok{ out)}
\NormalTok{  \}),}
  \FunctionTok{inner\_join}\NormalTok{(}
    \FunctionTok{prep\_ate}\NormalTok{(h\_bin}\SpecialCharTok{$}\NormalTok{marginal\_effects) }\SpecialCharTok{|\textgreater{}}
      \FunctionTok{rename}\NormalTok{(}\AttributeTok{estimate\_h =}\NormalTok{ estimate, }\AttributeTok{se\_h =}\NormalTok{ se, }\AttributeTok{p\_adj\_h =}\NormalTok{ p\_adj),}
    \FunctionTok{prep\_ate}\NormalTok{(h2\_bin}\SpecialCharTok{$}\NormalTok{marginal\_effects) }\SpecialCharTok{|\textgreater{}}
      \FunctionTok{rename}\NormalTok{(}\AttributeTok{estimate\_h2 =}\NormalTok{ estimate, }\AttributeTok{se\_h2 =}\NormalTok{ se, }\AttributeTok{p\_adj\_h2 =}\NormalTok{ p\_adj),}
    \AttributeTok{by =} \StringTok{"condition"}
\NormalTok{  ) }\SpecialCharTok{|\textgreater{}} \FunctionTok{mutate}\NormalTok{(}\AttributeTok{outcome =} \StringTok{"newsletter\_signup"}\NormalTok{)}
\NormalTok{) }\SpecialCharTok{|\textgreater{}}
  \FunctionTok{mutate}\NormalTok{(}
    \AttributeTok{tier =} \FunctionTok{case\_when}\NormalTok{(}
\NormalTok{      outcome }\SpecialCharTok{==}\NormalTok{ outcomes\_primary                          }\SpecialCharTok{\textasciitilde{}} \StringTok{"Primary"}\NormalTok{,}
\NormalTok{      outcome }\SpecialCharTok{\%in\%} \FunctionTok{c}\NormalTok{(outcomes\_secondary, outcomes\_binary) }\SpecialCharTok{\textasciitilde{}} \StringTok{"Secondary"}\NormalTok{,}
      \ConstantTok{TRUE}                                                \SpecialCharTok{\textasciitilde{}} \StringTok{"Tertiary"}
\NormalTok{    ),}
    \AttributeTok{outcome\_label =}\NormalTok{ outcome\_label\_map[outcome]}
\NormalTok{  )}

\CommentTok{\# Per{-}outcome summary metrics}
\NormalTok{ate\_metrics }\OtherTok{\textless{}{-}} \FunctionTok{map\_dfr}\NormalTok{(outcomes\_continuous, }\ControlFlowTok{function}\NormalTok{(out) \{}
  \FunctionTok{compare\_estimates}\NormalTok{(ate\_model\_results[[out]]}\SpecialCharTok{$}\NormalTok{h, ate\_model\_results[[out]]}\SpecialCharTok{$}\NormalTok{l) }\SpecialCharTok{|\textgreater{}}
    \FunctionTok{mutate}\NormalTok{(}\AttributeTok{outcome =}\NormalTok{ out)}
\NormalTok{\})}

\NormalTok{calib\_metrics }\OtherTok{\textless{}{-}} \FunctionTok{map\_dfr}\NormalTok{(outcomes\_continuous, }\ControlFlowTok{function}\NormalTok{(out) \{}
  \FunctionTok{run\_calibration}\NormalTok{(ate\_model\_results[[out]]}\SpecialCharTok{$}\NormalTok{h, ate\_model\_results[[out]]}\SpecialCharTok{$}\NormalTok{l) }\SpecialCharTok{|\textgreater{}}
    \FunctionTok{mutate}\NormalTok{(}\AttributeTok{outcome =}\NormalTok{ out)}
\NormalTok{\})}

\NormalTok{ate\_metrics }\OtherTok{\textless{}{-}} \FunctionTok{bind\_rows}\NormalTok{(}
\NormalTok{  ate\_metrics,}
  \FunctionTok{compare\_estimates}\NormalTok{(h\_bin}\SpecialCharTok{$}\NormalTok{marginal\_effects, l\_bin}\SpecialCharTok{$}\NormalTok{marginal\_effects) }\SpecialCharTok{|\textgreater{}}
    \FunctionTok{mutate}\NormalTok{(}\AttributeTok{outcome =} \StringTok{"newsletter\_signup"}\NormalTok{)}
\NormalTok{) }\SpecialCharTok{|\textgreater{}}
  \FunctionTok{mutate}\NormalTok{(}
    \AttributeTok{tier =} \FunctionTok{case\_when}\NormalTok{(}
\NormalTok{      outcome }\SpecialCharTok{==}\NormalTok{ outcomes\_primary                          }\SpecialCharTok{\textasciitilde{}} \StringTok{"Primary"}\NormalTok{,}
\NormalTok{      outcome }\SpecialCharTok{\%in\%} \FunctionTok{c}\NormalTok{(outcomes\_secondary, outcomes\_binary) }\SpecialCharTok{\textasciitilde{}} \StringTok{"Secondary"}\NormalTok{,}
      \ConstantTok{TRUE}                                                \SpecialCharTok{\textasciitilde{}} \StringTok{"Tertiary"}
\NormalTok{    ),}
    \AttributeTok{outcome\_label =}\NormalTok{ outcome\_label\_map[outcome]}
\NormalTok{  )}

\NormalTok{calib\_metrics }\OtherTok{\textless{}{-}} \FunctionTok{bind\_rows}\NormalTok{(}
\NormalTok{  calib\_metrics,}
  \FunctionTok{run\_calibration}\NormalTok{(h\_bin}\SpecialCharTok{$}\NormalTok{marginal\_effects, l\_bin}\SpecialCharTok{$}\NormalTok{marginal\_effects) }\SpecialCharTok{|\textgreater{}}
    \FunctionTok{mutate}\NormalTok{(}\AttributeTok{outcome =} \StringTok{"newsletter\_signup"}\NormalTok{)}
\NormalTok{) }\SpecialCharTok{|\textgreater{}}
  \FunctionTok{mutate}\NormalTok{(}
    \AttributeTok{tier =} \FunctionTok{case\_when}\NormalTok{(}
\NormalTok{      outcome }\SpecialCharTok{==}\NormalTok{ outcomes\_primary                          }\SpecialCharTok{\textasciitilde{}} \StringTok{"Primary"}\NormalTok{,}
\NormalTok{      outcome }\SpecialCharTok{\%in\%} \FunctionTok{c}\NormalTok{(outcomes\_secondary, outcomes\_binary) }\SpecialCharTok{\textasciitilde{}} \StringTok{"Secondary"}\NormalTok{,}
      \ConstantTok{TRUE}                                                \SpecialCharTok{\textasciitilde{}} \StringTok{"Tertiary"}
\NormalTok{    ),}
    \AttributeTok{outcome\_label =}\NormalTok{ outcome\_label\_map[outcome]}
\NormalTok{  )}

\CommentTok{\# Calibration regression for Human 1 vs. Human 2 (baseline ceiling)}
\NormalTok{calib\_metrics\_h2 }\OtherTok{\textless{}{-}} \FunctionTok{map\_dfr}\NormalTok{(outcomes\_continuous, }\ControlFlowTok{function}\NormalTok{(out) \{}
  \FunctionTok{run\_calibration}\NormalTok{(ate\_model\_results[[out]]}\SpecialCharTok{$}\NormalTok{h, ate\_model\_results[[out]]}\SpecialCharTok{$}\NormalTok{h2) }\SpecialCharTok{|\textgreater{}}
    \FunctionTok{mutate}\NormalTok{(}\AttributeTok{outcome =}\NormalTok{ out)}
\NormalTok{\})}

\NormalTok{calib\_metrics\_h2 }\OtherTok{\textless{}{-}} \FunctionTok{bind\_rows}\NormalTok{(}
\NormalTok{  calib\_metrics\_h2,}
  \FunctionTok{run\_calibration}\NormalTok{(h\_bin}\SpecialCharTok{$}\NormalTok{marginal\_effects, h2\_bin}\SpecialCharTok{$}\NormalTok{marginal\_effects) }\SpecialCharTok{|\textgreater{}}
    \FunctionTok{mutate}\NormalTok{(}\AttributeTok{outcome =} \StringTok{"newsletter\_signup"}\NormalTok{)}
\NormalTok{) }\SpecialCharTok{|\textgreater{}}
  \FunctionTok{mutate}\NormalTok{(}
    \AttributeTok{tier =} \FunctionTok{case\_when}\NormalTok{(}
\NormalTok{      outcome }\SpecialCharTok{==}\NormalTok{ outcomes\_primary                          }\SpecialCharTok{\textasciitilde{}} \StringTok{"Primary"}\NormalTok{,}
\NormalTok{      outcome }\SpecialCharTok{\%in\%} \FunctionTok{c}\NormalTok{(outcomes\_secondary, outcomes\_binary) }\SpecialCharTok{\textasciitilde{}} \StringTok{"Secondary"}\NormalTok{,}
      \ConstantTok{TRUE}                                                \SpecialCharTok{\textasciitilde{}} \StringTok{"Tertiary"}
\NormalTok{    ),}
    \AttributeTok{outcome\_label =}\NormalTok{ outcome\_label\_map[outcome]}
\NormalTok{  )}

\CommentTok{\# Pooled summary row — prepended to ate\_metrics so it appears at the top of tbl{-}rq1}
\NormalTok{pooled\_row }\OtherTok{\textless{}{-}}\NormalTok{ pooled\_ates }\SpecialCharTok{|\textgreater{}}
  \FunctionTok{mutate}\NormalTok{(}
    \AttributeTok{delta      =} \FloatTok{0.5} \SpecialCharTok{*} \FunctionTok{abs}\NormalTok{(estimate\_h),}
    \AttributeTok{diff       =}\NormalTok{ estimate\_h }\SpecialCharTok{{-}}\NormalTok{ estimate\_l,}
    \AttributeTok{se\_diff    =} \FunctionTok{sqrt}\NormalTok{(se\_h}\SpecialCharTok{\^{}}\DecValTok{2} \SpecialCharTok{+}\NormalTok{ se\_l}\SpecialCharTok{\^{}}\DecValTok{2}\NormalTok{),}
    \AttributeTok{p\_lower    =} \FunctionTok{pnorm}\NormalTok{((diff }\SpecialCharTok{+}\NormalTok{ delta) }\SpecialCharTok{/}\NormalTok{ se\_diff, }\AttributeTok{lower.tail =} \ConstantTok{FALSE}\NormalTok{),}
    \AttributeTok{p\_upper    =} \FunctionTok{pnorm}\NormalTok{((diff }\SpecialCharTok{{-}}\NormalTok{ delta) }\SpecialCharTok{/}\NormalTok{ se\_diff),}
    \AttributeTok{p\_tost     =} \FunctionTok{pmax}\NormalTok{(p\_lower, p\_upper),}
    \AttributeTok{equivalent =}\NormalTok{ p\_tost }\SpecialCharTok{\textless{}} \FloatTok{0.05}\NormalTok{,}
    \AttributeTok{same\_sign  =} \FunctionTok{sign}\NormalTok{(estimate\_h) }\SpecialCharTok{==} \FunctionTok{sign}\NormalTok{(estimate\_l),}
    \AttributeTok{infer\_h    =} \FunctionTok{case\_when}\NormalTok{(}
      \SpecialCharTok{!}\FunctionTok{is.na}\NormalTok{(p\_adj\_h) }\SpecialCharTok{\&}\NormalTok{ p\_adj\_h }\SpecialCharTok{\textless{}} \FloatTok{0.05} \SpecialCharTok{\&}\NormalTok{ estimate\_h }\SpecialCharTok{\textgreater{}} \DecValTok{0} \SpecialCharTok{\textasciitilde{}} \StringTok{"sig\_positive"}\NormalTok{,}
      \SpecialCharTok{!}\FunctionTok{is.na}\NormalTok{(p\_adj\_h) }\SpecialCharTok{\&}\NormalTok{ p\_adj\_h }\SpecialCharTok{\textless{}} \FloatTok{0.05} \SpecialCharTok{\&}\NormalTok{ estimate\_h }\SpecialCharTok{\textless{}} \DecValTok{0} \SpecialCharTok{\textasciitilde{}} \StringTok{"sig\_negative"}\NormalTok{,}
      \ConstantTok{TRUE} \SpecialCharTok{\textasciitilde{}} \StringTok{"not\_significant"}
\NormalTok{    ),}
    \AttributeTok{infer\_l    =} \FunctionTok{case\_when}\NormalTok{(}
      \SpecialCharTok{!}\FunctionTok{is.na}\NormalTok{(p\_adj\_l) }\SpecialCharTok{\&}\NormalTok{ p\_adj\_l }\SpecialCharTok{\textless{}} \FloatTok{0.05} \SpecialCharTok{\&}\NormalTok{ estimate\_l }\SpecialCharTok{\textgreater{}} \DecValTok{0} \SpecialCharTok{\textasciitilde{}} \StringTok{"sig\_positive"}\NormalTok{,}
      \SpecialCharTok{!}\FunctionTok{is.na}\NormalTok{(p\_adj\_l) }\SpecialCharTok{\&}\NormalTok{ p\_adj\_l }\SpecialCharTok{\textless{}} \FloatTok{0.05} \SpecialCharTok{\&}\NormalTok{ estimate\_l }\SpecialCharTok{\textless{}} \DecValTok{0} \SpecialCharTok{\textasciitilde{}} \StringTok{"sig\_negative"}\NormalTok{,}
      \ConstantTok{TRUE} \SpecialCharTok{\textasciitilde{}} \StringTok{"not\_significant"}
\NormalTok{    )}
\NormalTok{  ) }\SpecialCharTok{|\textgreater{}}
  \FunctionTok{summarise}\NormalTok{(}
    \AttributeTok{spearman\_rho    =} \FunctionTok{cor}\NormalTok{(estimate\_h, estimate\_l, }\AttributeTok{method =} \StringTok{"spearman"}\NormalTok{,}
                          \AttributeTok{use =} \StringTok{"pairwise.complete.obs"}\NormalTok{),}
    \AttributeTok{pearson\_r       =} \FunctionTok{cor}\NormalTok{(estimate\_h, estimate\_l, }\AttributeTok{use =} \StringTok{"pairwise.complete.obs"}\NormalTok{),}
    \AttributeTok{rmse            =} \ConstantTok{NA\_real\_}\NormalTok{,}
    \AttributeTok{directional\_pct =} \FunctionTok{mean}\NormalTok{(same\_sign,  }\AttributeTok{na.rm =} \ConstantTok{TRUE}\NormalTok{) }\SpecialCharTok{*} \DecValTok{100}\NormalTok{,}
    \AttributeTok{inferential\_pct =} \FunctionTok{mean}\NormalTok{(infer\_h }\SpecialCharTok{==}\NormalTok{ infer\_l, }\AttributeTok{na.rm =} \ConstantTok{TRUE}\NormalTok{) }\SpecialCharTok{*} \DecValTok{100}\NormalTok{,}
    \AttributeTok{tost\_pct        =} \FunctionTok{mean}\NormalTok{(equivalent, }\AttributeTok{na.rm =} \ConstantTok{TRUE}\NormalTok{) }\SpecialCharTok{*} \DecValTok{100}
\NormalTok{  ) }\SpecialCharTok{|\textgreater{}}
  \FunctionTok{mutate}\NormalTok{(}\AttributeTok{outcome =} \StringTok{"pooled"}\NormalTok{, }\AttributeTok{tier =} \StringTok{"All outcomes"}\NormalTok{, }\AttributeTok{outcome\_label =} \StringTok{"All outcomes"}\NormalTok{)}

\CommentTok{\# Per{-}outcome summary metrics for Human 1 vs. Human 2 (baseline ceiling)}
\NormalTok{ate\_metrics\_h2 }\OtherTok{\textless{}{-}} \FunctionTok{map\_dfr}\NormalTok{(outcomes\_continuous, }\ControlFlowTok{function}\NormalTok{(out) \{}
  \FunctionTok{compare\_estimates}\NormalTok{(ate\_model\_results[[out]]}\SpecialCharTok{$}\NormalTok{h, ate\_model\_results[[out]]}\SpecialCharTok{$}\NormalTok{h2) }\SpecialCharTok{|\textgreater{}}
    \FunctionTok{mutate}\NormalTok{(}\AttributeTok{outcome =}\NormalTok{ out)}
\NormalTok{\})}

\NormalTok{ate\_metrics\_h2 }\OtherTok{\textless{}{-}} \FunctionTok{bind\_rows}\NormalTok{(}
\NormalTok{  ate\_metrics\_h2,}
  \FunctionTok{compare\_estimates}\NormalTok{(h\_bin}\SpecialCharTok{$}\NormalTok{marginal\_effects, h2\_bin}\SpecialCharTok{$}\NormalTok{marginal\_effects) }\SpecialCharTok{|\textgreater{}}
    \FunctionTok{mutate}\NormalTok{(}\AttributeTok{outcome =} \StringTok{"newsletter\_signup"}\NormalTok{)}
\NormalTok{) }\SpecialCharTok{|\textgreater{}}
  \FunctionTok{mutate}\NormalTok{(}
    \AttributeTok{tier =} \FunctionTok{case\_when}\NormalTok{(}
\NormalTok{      outcome }\SpecialCharTok{==}\NormalTok{ outcomes\_primary                          }\SpecialCharTok{\textasciitilde{}} \StringTok{"Primary"}\NormalTok{,}
\NormalTok{      outcome }\SpecialCharTok{\%in\%} \FunctionTok{c}\NormalTok{(outcomes\_secondary, outcomes\_binary) }\SpecialCharTok{\textasciitilde{}} \StringTok{"Secondary"}\NormalTok{,}
      \ConstantTok{TRUE}                                                \SpecialCharTok{\textasciitilde{}} \StringTok{"Tertiary"}
\NormalTok{    ),}
    \AttributeTok{outcome\_label =}\NormalTok{ outcome\_label\_map[outcome]}
\NormalTok{  )}

\CommentTok{\# Pooled row for Human 1 vs. Human 2}
\NormalTok{pooled\_row\_h2 }\OtherTok{\textless{}{-}}\NormalTok{ pooled\_ates\_h2 }\SpecialCharTok{|\textgreater{}}
  \FunctionTok{mutate}\NormalTok{(}
    \AttributeTok{delta      =} \FloatTok{0.5} \SpecialCharTok{*} \FunctionTok{abs}\NormalTok{(estimate\_h),}
    \AttributeTok{diff       =}\NormalTok{ estimate\_h }\SpecialCharTok{{-}}\NormalTok{ estimate\_h2,}
    \AttributeTok{se\_diff    =} \FunctionTok{sqrt}\NormalTok{(se\_h}\SpecialCharTok{\^{}}\DecValTok{2} \SpecialCharTok{+}\NormalTok{ se\_h2}\SpecialCharTok{\^{}}\DecValTok{2}\NormalTok{),}
    \AttributeTok{p\_lower    =} \FunctionTok{pnorm}\NormalTok{((diff }\SpecialCharTok{+}\NormalTok{ delta) }\SpecialCharTok{/}\NormalTok{ se\_diff, }\AttributeTok{lower.tail =} \ConstantTok{FALSE}\NormalTok{),}
    \AttributeTok{p\_upper    =} \FunctionTok{pnorm}\NormalTok{((diff }\SpecialCharTok{{-}}\NormalTok{ delta) }\SpecialCharTok{/}\NormalTok{ se\_diff),}
    \AttributeTok{p\_tost     =} \FunctionTok{pmax}\NormalTok{(p\_lower, p\_upper),}
    \AttributeTok{equivalent =}\NormalTok{ p\_tost }\SpecialCharTok{\textless{}} \FloatTok{0.05}\NormalTok{,}
    \AttributeTok{same\_sign  =} \FunctionTok{sign}\NormalTok{(estimate\_h) }\SpecialCharTok{==} \FunctionTok{sign}\NormalTok{(estimate\_h2),}
    \AttributeTok{infer\_h    =} \FunctionTok{case\_when}\NormalTok{(}
      \SpecialCharTok{!}\FunctionTok{is.na}\NormalTok{(p\_adj\_h)  }\SpecialCharTok{\&}\NormalTok{ p\_adj\_h  }\SpecialCharTok{\textless{}} \FloatTok{0.05} \SpecialCharTok{\&}\NormalTok{ estimate\_h  }\SpecialCharTok{\textgreater{}} \DecValTok{0} \SpecialCharTok{\textasciitilde{}} \StringTok{"sig\_positive"}\NormalTok{,}
      \SpecialCharTok{!}\FunctionTok{is.na}\NormalTok{(p\_adj\_h)  }\SpecialCharTok{\&}\NormalTok{ p\_adj\_h  }\SpecialCharTok{\textless{}} \FloatTok{0.05} \SpecialCharTok{\&}\NormalTok{ estimate\_h  }\SpecialCharTok{\textless{}} \DecValTok{0} \SpecialCharTok{\textasciitilde{}} \StringTok{"sig\_negative"}\NormalTok{,}
      \ConstantTok{TRUE} \SpecialCharTok{\textasciitilde{}} \StringTok{"not\_significant"}
\NormalTok{    ),}
    \AttributeTok{infer\_h2   =} \FunctionTok{case\_when}\NormalTok{(}
      \SpecialCharTok{!}\FunctionTok{is.na}\NormalTok{(p\_adj\_h2) }\SpecialCharTok{\&}\NormalTok{ p\_adj\_h2 }\SpecialCharTok{\textless{}} \FloatTok{0.05} \SpecialCharTok{\&}\NormalTok{ estimate\_h2 }\SpecialCharTok{\textgreater{}} \DecValTok{0} \SpecialCharTok{\textasciitilde{}} \StringTok{"sig\_positive"}\NormalTok{,}
      \SpecialCharTok{!}\FunctionTok{is.na}\NormalTok{(p\_adj\_h2) }\SpecialCharTok{\&}\NormalTok{ p\_adj\_h2 }\SpecialCharTok{\textless{}} \FloatTok{0.05} \SpecialCharTok{\&}\NormalTok{ estimate\_h2 }\SpecialCharTok{\textless{}} \DecValTok{0} \SpecialCharTok{\textasciitilde{}} \StringTok{"sig\_negative"}\NormalTok{,}
      \ConstantTok{TRUE} \SpecialCharTok{\textasciitilde{}} \StringTok{"not\_significant"}
\NormalTok{    )}
\NormalTok{  ) }\SpecialCharTok{|\textgreater{}}
  \FunctionTok{summarise}\NormalTok{(}
    \AttributeTok{spearman\_rho    =} \FunctionTok{cor}\NormalTok{(estimate\_h, estimate\_h2, }\AttributeTok{method =} \StringTok{"spearman"}\NormalTok{,}
                          \AttributeTok{use =} \StringTok{"pairwise.complete.obs"}\NormalTok{),}
    \AttributeTok{pearson\_r       =} \FunctionTok{cor}\NormalTok{(estimate\_h, estimate\_h2, }\AttributeTok{use =} \StringTok{"pairwise.complete.obs"}\NormalTok{),}
    \AttributeTok{rmse            =} \ConstantTok{NA\_real\_}\NormalTok{,}
    \AttributeTok{directional\_pct =} \FunctionTok{mean}\NormalTok{(same\_sign,  }\AttributeTok{na.rm =} \ConstantTok{TRUE}\NormalTok{) }\SpecialCharTok{*} \DecValTok{100}\NormalTok{,}
    \AttributeTok{inferential\_pct =} \FunctionTok{mean}\NormalTok{(infer\_h }\SpecialCharTok{==}\NormalTok{ infer\_h2, }\AttributeTok{na.rm =} \ConstantTok{TRUE}\NormalTok{) }\SpecialCharTok{*} \DecValTok{100}\NormalTok{,}
    \AttributeTok{tost\_pct        =} \FunctionTok{mean}\NormalTok{(equivalent, }\AttributeTok{na.rm =} \ConstantTok{TRUE}\NormalTok{) }\SpecialCharTok{*} \DecValTok{100}
\NormalTok{  ) }\SpecialCharTok{|\textgreater{}}
  \FunctionTok{mutate}\NormalTok{(}\AttributeTok{outcome =} \StringTok{"pooled"}\NormalTok{, }\AttributeTok{tier =} \StringTok{"All outcomes"}\NormalTok{, }\AttributeTok{outcome\_label =} \StringTok{"All outcomes"}\NormalTok{)}

\NormalTok{ate\_metrics\_h2 }\OtherTok{\textless{}{-}} \FunctionTok{bind\_rows}\NormalTok{(pooled\_row\_h2, ate\_metrics\_h2)}

\NormalTok{ate\_metrics }\OtherTok{\textless{}{-}} \FunctionTok{bind\_rows}\NormalTok{(pooled\_row, ate\_metrics)}
\end{Highlighting}
\end{Shaded}

\paragraph{All outcomes}\label{all-outcomes}

Figure~\ref{fig-rq1-pooled} shows all 240 estimate pairs --- one per
intervention × outcome combination --- colored by outcome tier. Points
above the diagonal indicate clone overestimation; points below indicate
underestimation.

Table~\ref{tbl-rq1} includes a pooled summary row at the top, followed
by per-outcome rows grouped by tier. RMSE is left blank for the pooled
row because it is not comparable across outcomes on different scales;
all other metrics are scale-invariant. Both comparison pairs (Human 1
vs.~LLM and Human 1 vs.~Human 2) are reported side by side; the Human 2
columns serve as a human-human ceiling for interpreting LLM performance.

Table~\ref{tbl-rq1-calibration} shows calibration regression results for
both comparisons (Human 1 vs.~LLM and Human 1 vs.~Human 2).
Figure~\ref{fig-rq1-all} shows the per-outcome scatter for Human 1
vs.~LLM only.

\begin{Shaded}
\begin{Highlighting}[]
\NormalTok{tier\_colors }\OtherTok{\textless{}{-}} \FunctionTok{c}\NormalTok{(}\StringTok{"Primary"} \OtherTok{=} \StringTok{"\#C0392B"}\NormalTok{, }\StringTok{"Secondary"} \OtherTok{=} \StringTok{"\#2980B9"}\NormalTok{, }\StringTok{"Tertiary"} \OtherTok{=} \StringTok{"\#7F8C8D"}\NormalTok{)}

\FunctionTok{bind\_rows}\NormalTok{(}
\NormalTok{  pooled\_ates    }\SpecialCharTok{|\textgreater{}} \FunctionTok{mutate}\NormalTok{(}\AttributeTok{estimate\_pred =}\NormalTok{ estimate\_l,  }\AttributeTok{comparison =} \StringTok{"Human 1 vs. LLM"}\NormalTok{),}
\NormalTok{  pooled\_ates\_h2 }\SpecialCharTok{|\textgreater{}} \FunctionTok{mutate}\NormalTok{(}\AttributeTok{estimate\_pred =}\NormalTok{ estimate\_h2, }\AttributeTok{comparison =} \StringTok{"Human 1 vs. Human 2"}\NormalTok{)}
\NormalTok{) }\SpecialCharTok{|\textgreater{}}
  \FunctionTok{ggplot}\NormalTok{(}\FunctionTok{aes}\NormalTok{(}\AttributeTok{x =}\NormalTok{ estimate\_h, }\AttributeTok{y =}\NormalTok{ estimate\_pred, }\AttributeTok{color =}\NormalTok{ tier)) }\SpecialCharTok{+}
  \FunctionTok{geom\_abline}\NormalTok{(}\AttributeTok{slope =} \DecValTok{1}\NormalTok{, }\AttributeTok{intercept =} \DecValTok{0}\NormalTok{, }\AttributeTok{linetype =} \StringTok{"dashed"}\NormalTok{,}
              \AttributeTok{color =} \StringTok{"grey60"}\NormalTok{, }\AttributeTok{linewidth =} \FloatTok{0.5}\NormalTok{) }\SpecialCharTok{+}
  \FunctionTok{geom\_hline}\NormalTok{(}\AttributeTok{yintercept =} \DecValTok{0}\NormalTok{, }\AttributeTok{linetype =} \StringTok{"dotted"}\NormalTok{, }\AttributeTok{color =} \StringTok{"grey80"}\NormalTok{) }\SpecialCharTok{+}
  \FunctionTok{geom\_vline}\NormalTok{(}\AttributeTok{xintercept =} \DecValTok{0}\NormalTok{, }\AttributeTok{linetype =} \StringTok{"dotted"}\NormalTok{, }\AttributeTok{color =} \StringTok{"grey80"}\NormalTok{) }\SpecialCharTok{+}
  \FunctionTok{geom\_point}\NormalTok{(}\AttributeTok{size =} \DecValTok{2}\NormalTok{, }\AttributeTok{alpha =} \FloatTok{0.65}\NormalTok{) }\SpecialCharTok{+}
  \FunctionTok{scale\_color\_manual}\NormalTok{(}\AttributeTok{values =}\NormalTok{ tier\_colors, }\AttributeTok{name =} \StringTok{"Outcome tier"}\NormalTok{) }\SpecialCharTok{+}
  \FunctionTok{facet\_wrap}\NormalTok{(}\SpecialCharTok{\textasciitilde{}}\NormalTok{ comparison, }\AttributeTok{ncol =} \DecValTok{2}\NormalTok{) }\SpecialCharTok{+}
  \FunctionTok{labs}\NormalTok{(}\AttributeTok{x =} \StringTok{"Human 1 ATE"}\NormalTok{, }\AttributeTok{y =} \StringTok{"Predicted ATE"}\NormalTok{) }\SpecialCharTok{+}
\NormalTok{  plot\_theme }\SpecialCharTok{+}
  \FunctionTok{theme}\NormalTok{(}\AttributeTok{legend.position =} \StringTok{"top"}\NormalTok{)}
\end{Highlighting}
\end{Shaded}

\begin{figure}[H]

\centering{

\pandocbounded{\includegraphics[keepaspectratio]{preregistration_files/figure-pdf/fig-rq1-pooled-1.pdf}}

}

\caption{\label{fig-rq1-pooled}ATEs pooled across all outcomes for both
comparisons. Left panel: Human 1 vs.~LLM clones. Right panel: Human 1
vs.~Human 2 (baseline ceiling). Each point is one intervention × outcome
pair (N = 240 per panel). The dashed line is perfect agreement; dotted
lines mark zero. Color indicates outcome tier.}

\end{figure}%

\begin{Shaded}
\begin{Highlighting}[]
\NormalTok{ate\_metrics }\SpecialCharTok{|\textgreater{}}
  \FunctionTok{rename\_with}\NormalTok{(}\SpecialCharTok{\textasciitilde{}} \FunctionTok{paste0}\NormalTok{(.x, }\StringTok{"\_llm"}\NormalTok{),}
              \AttributeTok{.cols =} \FunctionTok{c}\NormalTok{(spearman\_rho, pearson\_r, rmse,}
\NormalTok{                        directional\_pct, inferential\_pct, tost\_pct)) }\SpecialCharTok{|\textgreater{}}
  \FunctionTok{left\_join}\NormalTok{(}
\NormalTok{    ate\_metrics\_h2 }\SpecialCharTok{|\textgreater{}}
      \FunctionTok{rename\_with}\NormalTok{(}\SpecialCharTok{\textasciitilde{}} \FunctionTok{paste0}\NormalTok{(.x, }\StringTok{"\_h2"}\NormalTok{),}
                  \AttributeTok{.cols =} \FunctionTok{c}\NormalTok{(spearman\_rho, pearson\_r, rmse,}
\NormalTok{                             directional\_pct, inferential\_pct, tost\_pct)) }\SpecialCharTok{|\textgreater{}}
      \FunctionTok{select}\NormalTok{(outcome, }\FunctionTok{ends\_with}\NormalTok{(}\StringTok{"\_h2"}\NormalTok{)),}
    \AttributeTok{by =} \StringTok{"outcome"}
\NormalTok{  ) }\SpecialCharTok{|\textgreater{}}
  \FunctionTok{arrange}\NormalTok{(}\FunctionTok{match}\NormalTok{(tier, }\FunctionTok{c}\NormalTok{(}\StringTok{"All outcomes"}\NormalTok{, }\StringTok{"Primary"}\NormalTok{, }\StringTok{"Secondary"}\NormalTok{, }\StringTok{"Tertiary"}\NormalTok{)), outcome\_label) }\SpecialCharTok{|\textgreater{}}
  \FunctionTok{select}\NormalTok{(}\AttributeTok{Tier =}\NormalTok{ tier, }\AttributeTok{Outcome =}\NormalTok{ outcome\_label,}
\NormalTok{         spearman\_rho\_llm, pearson\_r\_llm, rmse\_llm,}
\NormalTok{         directional\_pct\_llm, inferential\_pct\_llm, tost\_pct\_llm,}
\NormalTok{         spearman\_rho\_h2, pearson\_r\_h2, rmse\_h2,}
\NormalTok{         directional\_pct\_h2, inferential\_pct\_h2, tost\_pct\_h2) }\SpecialCharTok{|\textgreater{}}
  \FunctionTok{gt}\NormalTok{(}\AttributeTok{groupname\_col =} \StringTok{"Tier"}\NormalTok{) }\SpecialCharTok{|\textgreater{}}
  \FunctionTok{cols\_label}\NormalTok{(}
    \AttributeTok{spearman\_rho\_llm    =} \StringTok{"Spearman ρ"}\NormalTok{, }\AttributeTok{pearson\_r\_llm    =} \StringTok{"Pearson r"}\NormalTok{, }\AttributeTok{rmse\_llm    =} \StringTok{"RMSE"}\NormalTok{,}
    \AttributeTok{directional\_pct\_llm =} \StringTok{"Dir. (\%)"}\NormalTok{,   }\AttributeTok{inferential\_pct\_llm =} \StringTok{"Infer. (\%)"}\NormalTok{, }\AttributeTok{tost\_pct\_llm =} \StringTok{"TOST (\%)"}\NormalTok{,}
    \AttributeTok{spearman\_rho\_h2     =} \StringTok{"Spearman ρ"}\NormalTok{, }\AttributeTok{pearson\_r\_h2     =} \StringTok{"Pearson r"}\NormalTok{, }\AttributeTok{rmse\_h2     =} \StringTok{"RMSE"}\NormalTok{,}
    \AttributeTok{directional\_pct\_h2  =} \StringTok{"Dir. (\%)"}\NormalTok{,   }\AttributeTok{inferential\_pct\_h2  =} \StringTok{"Infer. (\%)"}\NormalTok{, }\AttributeTok{tost\_pct\_h2  =} \StringTok{"TOST (\%)"}
\NormalTok{  ) }\SpecialCharTok{|\textgreater{}}
  \FunctionTok{tab\_spanner}\NormalTok{(}\AttributeTok{label =} \StringTok{"Human 1 vs. LLM"}\NormalTok{, }\AttributeTok{columns =} \FunctionTok{ends\_with}\NormalTok{(}\StringTok{"\_llm"}\NormalTok{)) }\SpecialCharTok{|\textgreater{}}
  \FunctionTok{tab\_spanner}\NormalTok{(}\AttributeTok{label =} \StringTok{"Human 1 vs. Human 2"}\NormalTok{, }\AttributeTok{columns =} \FunctionTok{ends\_with}\NormalTok{(}\StringTok{"\_h2"}\NormalTok{)) }\SpecialCharTok{|\textgreater{}}
  \FunctionTok{fmt\_number}\NormalTok{(}\AttributeTok{columns =} \FunctionTok{c}\NormalTok{(spearman\_rho\_llm, pearson\_r\_llm, rmse\_llm,}
\NormalTok{                          spearman\_rho\_h2,  pearson\_r\_h2,  rmse\_h2), }\AttributeTok{decimals =} \DecValTok{3}\NormalTok{) }\SpecialCharTok{|\textgreater{}}
  \FunctionTok{sub\_missing}\NormalTok{(}\AttributeTok{columns =} \FunctionTok{c}\NormalTok{(rmse\_llm, rmse\_h2), }\AttributeTok{missing\_text =} \StringTok{"—"}\NormalTok{) }\SpecialCharTok{|\textgreater{}}
  \FunctionTok{fmt\_number}\NormalTok{(}\AttributeTok{columns =} \FunctionTok{c}\NormalTok{(directional\_pct\_llm, inferential\_pct\_llm, tost\_pct\_llm,}
\NormalTok{                          directional\_pct\_h2,  inferential\_pct\_h2,  tost\_pct\_h2), }\AttributeTok{decimals =} \DecValTok{1}\NormalTok{) }\SpecialCharTok{|\textgreater{}}
  \FunctionTok{tab\_style}\NormalTok{(}
    \AttributeTok{style     =} \FunctionTok{cell\_text}\NormalTok{(}\AttributeTok{weight =} \StringTok{"bold"}\NormalTok{),}
    \AttributeTok{locations =} \FunctionTok{cells\_row\_groups}\NormalTok{()}
\NormalTok{  )}
\end{Highlighting}
\end{Shaded}

\begin{table}

\caption{\label{tbl-rq1}RQ1 comparison metrics. The `All outcomes' row
pools all estimate pairs across outcomes; RMSE (---) is omitted because
outcomes differ in scale. Per-outcome rows each have N = 20 intervention
pairs. Left columns: Human 1 vs.~LLM clone; right columns: Human 1
vs.~Human 2 (baseline ceiling). TOST: \% of pairs achieving equivalence
(Δ = 50\% of human ATE).}

\centering{

\fontsize{12.0pt}{14.4pt}\selectfont
\begin{tabular*}{\linewidth}{@{\extracolsep{\fill}}lrrrrrrrrrrrr}
\toprule
 & \multicolumn{6}{c}{Human 1 vs. LLM} & \multicolumn{6}{c}{Human 1 vs. Human 2} \\ 
\cmidrule(lr){2-7} \cmidrule(lr){8-13}
Outcome & Spearman ρ & Pearson r & RMSE & Dir. (\%) & Infer. (\%) & TOST (\%) & Spearman ρ & Pearson r & RMSE & Dir. (\%) & Infer. (\%) & TOST (\%) \\ 
\midrule\addlinespace[2.5pt]
\multicolumn{13}{l}{{\bfseries All outcomes}} \\[2.5pt] 
\midrule\addlinespace[2.5pt]
All outcomes & 0.115 & 0.196 & — & 51.2 & 99.2 & 0.0 & -0.087 & -0.056 & — & 48.8 & 97.9 & 0.0 \\ 
\midrule\addlinespace[2.5pt]
\multicolumn{13}{l}{{\bfseries Primary}} \\[2.5pt] 
\midrule\addlinespace[2.5pt]
Multidimensional trust (primary) & -0.208 & -0.187 & 0.556 & 45.0 & 100.0 & 0.0 & -0.289 & -0.236 & 0.746 & 45.0 & 100.0 & 0.0 \\ 
\midrule\addlinespace[2.5pt]
\multicolumn{13}{l}{{\bfseries Secondary}} \\[2.5pt] 
\midrule\addlinespace[2.5pt]
Distrust & -0.024 & -0.229 & 2.284 & 75.0 & 90.0 & 0.0 & -0.167 & -0.076 & 3.017 & 40.0 & 90.0 & 0.0 \\ 
Funding perceptions & 0.086 & 0.064 & 2.254 & 50.0 & 100.0 & 0.0 & 0.153 & 0.281 & 2.588 & 50.0 & 100.0 & 0.0 \\ 
Institutional trust & -0.296 & -0.223 & 1.075 & 50.0 & 100.0 & 0.0 & 0.217 & 0.171 & 0.935 & 60.0 & 100.0 & 0.0 \\ 
Newsletter sign-up & 0.448 & 0.435 & 0.017 & 55.0 & 100.0 & 0.0 & -0.304 & -0.183 & 0.032 & 60.0 & 100.0 & 0.0 \\ 
Overall trust & -0.006 & 0.139 & 1.558 & 40.0 & 100.0 & 0.0 & -0.002 & 0.058 & 1.854 & 50.0 & 100.0 & 0.0 \\ 
Policy role & -0.295 & -0.076 & 0.887 & 30.0 & 100.0 & 0.0 & -0.143 & -0.075 & 0.898 & 45.0 & 100.0 & 0.0 \\ 
\midrule\addlinespace[2.5pt]
\multicolumn{13}{l}{{\bfseries Tertiary}} \\[2.5pt] 
\midrule\addlinespace[2.5pt]
Climate belief & -0.036 & -0.078 & 1.861 & 65.0 & 100.0 & 0.0 & 0.071 & 0.275 & 2.073 & 35.0 & 100.0 & 0.0 \\ 
Climate concern & 0.230 & 0.342 & 1.444 & 40.0 & 100.0 & 0.0 & 0.168 & -0.094 & 1.211 & 55.0 & 100.0 & 0.0 \\ 
General climate policy & -0.050 & 0.090 & 1.592 & 75.0 & 100.0 & 0.0 & 0.173 & 0.266 & 2.814 & 20.0 & 100.0 & 0.0 \\ 
Pro-climate behavior & -0.137 & -0.111 & 1.197 & 35.0 & 100.0 & 0.0 & 0.048 & 0.035 & 0.830 & 85.0 & 85.0 & 0.0 \\ 
Specific climate policy & 0.226 & 0.366 & 0.471 & 55.0 & 100.0 & 0.0 & -0.202 & -0.189 & 0.822 & 40.0 & 100.0 & 0.0 \\ 
\bottomrule
\end{tabular*}

}

\end{table}%

\begin{Shaded}
\begin{Highlighting}[]
\NormalTok{calib\_metrics }\SpecialCharTok{|\textgreater{}}
  \FunctionTok{left\_join}\NormalTok{(}
\NormalTok{    calib\_metrics\_h2 }\SpecialCharTok{|\textgreater{}}
      \FunctionTok{select}\NormalTok{(outcome,}
             \AttributeTok{alpha\_h2 =}\NormalTok{ alpha, }\AttributeTok{alpha\_lo\_h2 =}\NormalTok{ alpha\_lo, }\AttributeTok{alpha\_hi\_h2 =}\NormalTok{ alpha\_hi,}
             \AttributeTok{beta\_h2  =}\NormalTok{ beta,  }\AttributeTok{beta\_lo\_h2  =}\NormalTok{ beta\_lo,  }\AttributeTok{beta\_hi\_h2  =}\NormalTok{ beta\_hi,}
             \AttributeTok{r\_squared\_h2 =}\NormalTok{ r\_squared, }\AttributeTok{p\_beta\_1\_h2 =}\NormalTok{ p\_beta\_1),}
    \AttributeTok{by =} \StringTok{"outcome"}
\NormalTok{  ) }\SpecialCharTok{|\textgreater{}}
  \FunctionTok{arrange}\NormalTok{(}\FunctionTok{match}\NormalTok{(tier, }\FunctionTok{c}\NormalTok{(}\StringTok{"Primary"}\NormalTok{, }\StringTok{"Secondary"}\NormalTok{, }\StringTok{"Tertiary"}\NormalTok{)), outcome\_label) }\SpecialCharTok{|\textgreater{}}
  \FunctionTok{mutate}\NormalTok{(}
    \AttributeTok{beta\_ci      =} \FunctionTok{paste0}\NormalTok{(}\StringTok{"["}\NormalTok{, }\FunctionTok{round}\NormalTok{(beta\_lo,     }\DecValTok{2}\NormalTok{), }\StringTok{", "}\NormalTok{, }\FunctionTok{round}\NormalTok{(beta\_hi,     }\DecValTok{2}\NormalTok{), }\StringTok{"]"}\NormalTok{),}
    \AttributeTok{alpha\_ci     =} \FunctionTok{paste0}\NormalTok{(}\StringTok{"["}\NormalTok{, }\FunctionTok{round}\NormalTok{(alpha\_lo,    }\DecValTok{2}\NormalTok{), }\StringTok{", "}\NormalTok{, }\FunctionTok{round}\NormalTok{(alpha\_hi,    }\DecValTok{2}\NormalTok{), }\StringTok{"]"}\NormalTok{),}
    \AttributeTok{beta\_ci\_h2   =} \FunctionTok{paste0}\NormalTok{(}\StringTok{"["}\NormalTok{, }\FunctionTok{round}\NormalTok{(beta\_lo\_h2,  }\DecValTok{2}\NormalTok{), }\StringTok{", "}\NormalTok{, }\FunctionTok{round}\NormalTok{(beta\_hi\_h2,  }\DecValTok{2}\NormalTok{), }\StringTok{"]"}\NormalTok{),}
    \AttributeTok{alpha\_ci\_h2  =} \FunctionTok{paste0}\NormalTok{(}\StringTok{"["}\NormalTok{, }\FunctionTok{round}\NormalTok{(alpha\_lo\_h2, }\DecValTok{2}\NormalTok{), }\StringTok{", "}\NormalTok{, }\FunctionTok{round}\NormalTok{(alpha\_hi\_h2, }\DecValTok{2}\NormalTok{), }\StringTok{"]"}\NormalTok{)}
\NormalTok{  ) }\SpecialCharTok{|\textgreater{}}
  \FunctionTok{select}\NormalTok{(}\AttributeTok{Tier =}\NormalTok{ tier, }\AttributeTok{Outcome =}\NormalTok{ outcome\_label,}
\NormalTok{         alpha, alpha\_ci, beta, beta\_ci, r\_squared, p\_beta\_1,}
\NormalTok{         alpha\_h2, alpha\_ci\_h2, beta\_h2, beta\_ci\_h2, r\_squared\_h2, p\_beta\_1\_h2) }\SpecialCharTok{|\textgreater{}}
  \FunctionTok{gt}\NormalTok{(}\AttributeTok{groupname\_col =} \StringTok{"Tier"}\NormalTok{) }\SpecialCharTok{|\textgreater{}}
  \FunctionTok{cols\_label}\NormalTok{(}
    \AttributeTok{alpha =} \StringTok{"α"}\NormalTok{, }\AttributeTok{alpha\_ci =} \StringTok{"95\% CI (α)"}\NormalTok{, }\AttributeTok{beta =} \StringTok{"β"}\NormalTok{, }\AttributeTok{beta\_ci =} \StringTok{"95\% CI (β)"}\NormalTok{,}
    \AttributeTok{r\_squared =} \StringTok{"R²"}\NormalTok{, }\AttributeTok{p\_beta\_1 =} \StringTok{"p (β=1)"}\NormalTok{,}
    \AttributeTok{alpha\_h2 =} \StringTok{"α"}\NormalTok{, }\AttributeTok{alpha\_ci\_h2 =} \StringTok{"95\% CI (α)"}\NormalTok{, }\AttributeTok{beta\_h2 =} \StringTok{"β"}\NormalTok{, }\AttributeTok{beta\_ci\_h2 =} \StringTok{"95\% CI (β)"}\NormalTok{,}
    \AttributeTok{r\_squared\_h2 =} \StringTok{"R²"}\NormalTok{, }\AttributeTok{p\_beta\_1\_h2 =} \StringTok{"p (β=1)"}
\NormalTok{  ) }\SpecialCharTok{|\textgreater{}}
  \FunctionTok{tab\_spanner}\NormalTok{(}\AttributeTok{label =} \StringTok{"Human 1 vs. LLM"}\NormalTok{,}
              \AttributeTok{columns =} \FunctionTok{c}\NormalTok{(alpha, alpha\_ci, beta, beta\_ci, r\_squared, p\_beta\_1)) }\SpecialCharTok{|\textgreater{}}
  \FunctionTok{tab\_spanner}\NormalTok{(}\AttributeTok{label =} \StringTok{"Human 1 vs. Human 2"}\NormalTok{,}
              \AttributeTok{columns =} \FunctionTok{c}\NormalTok{(alpha\_h2, alpha\_ci\_h2, beta\_h2, beta\_ci\_h2, r\_squared\_h2, p\_beta\_1\_h2)) }\SpecialCharTok{|\textgreater{}}
  \FunctionTok{fmt\_number}\NormalTok{(}\AttributeTok{columns =} \FunctionTok{c}\NormalTok{(alpha, beta, r\_squared, alpha\_h2, beta\_h2, r\_squared\_h2), }\AttributeTok{decimals =} \DecValTok{3}\NormalTok{) }\SpecialCharTok{|\textgreater{}}
  \FunctionTok{fmt\_number}\NormalTok{(}\AttributeTok{columns =} \FunctionTok{c}\NormalTok{(p\_beta\_1, p\_beta\_1\_h2), }\AttributeTok{decimals =} \DecValTok{3}\NormalTok{) }\SpecialCharTok{|\textgreater{}}
  \FunctionTok{tab\_style}\NormalTok{(}
    \AttributeTok{style     =} \FunctionTok{cell\_text}\NormalTok{(}\AttributeTok{weight =} \StringTok{"bold"}\NormalTok{),}
    \AttributeTok{locations =} \FunctionTok{cells\_row\_groups}\NormalTok{()}
\NormalTok{  )}
\end{Highlighting}
\end{Shaded}

\begin{table}

\caption{\label{tbl-rq1-calibration}Calibration regression by outcome.
α: intercept (systematic bias, with 95\% CI). β: slope (proportionality;
1 = perfect, with 95\% CI). R²: variance explained. p(β=1): F-test of
proportional calibration. Left columns: Human 1 vs.~LLM; right columns:
Human 1 vs.~Human 2 (baseline ceiling).}

\centering{

\fontsize{12.0pt}{14.4pt}\selectfont
\begin{tabular*}{\linewidth}{@{\extracolsep{\fill}}lrlrlrrrlrlrr}
\toprule
 & \multicolumn{6}{c}{Human 1 vs. LLM} & \multicolumn{6}{c}{Human 1 vs. Human 2} \\ 
\cmidrule(lr){2-7} \cmidrule(lr){8-13}
Outcome & α & 95\% CI (α) & β & 95\% CI (β) & R² & p (β=1) & α & 95\% CI (α) & β & 95\% CI (β) & R² & p (β=1) \\ 
\midrule\addlinespace[2.5pt]
\multicolumn{13}{l}{{\bfseries Primary}} \\[2.5pt] 
\midrule\addlinespace[2.5pt]
Multidimensional trust (primary) & 0.263 & [0.1, 0.43] & -0.199 & [-0.72, 0.32] & 0.035 & 0.000 & 0.245 & [0.08, 0.41] & -0.160 & [-0.49, 0.17] & 0.056 & 0.000 \\ 
\midrule\addlinespace[2.5pt]
\multicolumn{13}{l}{{\bfseries Secondary}} \\[2.5pt] 
\midrule\addlinespace[2.5pt]
Distrust & 2.309 & [1.13, 3.49] & -0.449 & [-1.39, 0.49] & 0.053 & 0.005 & 1.888 & [1.05, 2.72] & -0.096 & [-0.72, 0.53] & 0.006 & 0.002 \\ 
Funding perceptions & 1.418 & [0.7, 2.14] & 0.114 & [-0.77, 0.99] & 0.004 & 0.049 & 1.489 & [0.81, 2.17] & 0.235 & [-0.16, 0.63] & 0.079 & 0.001 \\ 
Institutional trust & -0.187 & [-0.73, 0.35] & -0.296 & [-0.94, 0.34] & 0.050 & 0.000 & -0.050 & [-0.4, 0.3] & 0.144 & [-0.27, 0.56] & 0.029 & 0.000 \\ 
Newsletter sign-up & -0.001 & [-0.01, 0.01] & 0.489 & [-0.01, 0.99] & 0.190 & 0.046 & 0.005 & [-0.01, 0.02] & -0.155 & [-0.57, 0.26] & 0.033 & 0.000 \\ 
Overall trust & 0.329 & [-0.38, 1.03] & 0.204 & [-0.52, 0.92] & 0.019 & 0.032 & 0.371 & [-0.32, 1.06] & 0.061 & [-0.46, 0.58] & 0.003 & 0.001 \\ 
Policy role & -0.223 & [-0.54, 0.1] & -0.092 & [-0.69, 0.51] & 0.006 & 0.001 & -0.214 & [-0.53, 0.11] & -0.095 & [-0.72, 0.53] & 0.006 & 0.002 \\ 
\midrule\addlinespace[2.5pt]
\multicolumn{13}{l}{{\bfseries Tertiary}} \\[2.5pt] 
\midrule\addlinespace[2.5pt]
Climate belief & 1.223 & [0.49, 1.95] & -0.135 & [-0.99, 0.72] & 0.006 & 0.012 & 1.206 & [0.55, 1.86] & 0.282 & [-0.21, 0.77] & 0.076 & 0.006 \\ 
Climate concern & -0.967 & [-1.54, -0.39] & 0.554 & [-0.2, 1.31] & 0.117 & 0.229 & -0.647 & [-1.03, -0.26] & -0.106 & [-0.66, 0.45] & 0.009 & 0.001 \\ 
General climate policy & -1.463 & [-2.13, -0.8] & 0.097 & [-0.44, 0.63] & 0.008 & 0.002 & -1.868 & [-2.65, -1.09] & 0.326 & [-0.26, 0.91] & 0.071 & 0.026 \\ 
Pro-climate behavior & 0.883 & [0.6, 1.17] & -0.202 & [-1.09, 0.69] & 0.012 & 0.011 & 0.872 & [0.36, 1.39] & 0.030 & [-0.4, 0.46] & 0.001 & 0.000 \\ 
Specific climate policy & -0.186 & [-0.38, 0.01] & 0.465 & [-0.12, 1.05] & 0.134 & 0.070 & -0.145 & [-0.38, 0.09] & -0.170 & [-0.61, 0.27] & 0.036 & 0.000 \\ 
\bottomrule
\end{tabular*}

}

\end{table}%

\begin{Shaded}
\begin{Highlighting}[]
\FunctionTok{map\_dfr}\NormalTok{(outcomes\_continuous, }\ControlFlowTok{function}\NormalTok{(out) \{}
  \FunctionTok{inner\_join}\NormalTok{(}
\NormalTok{    ate\_model\_results[[out]]}\SpecialCharTok{$}\NormalTok{h }\SpecialCharTok{|\textgreater{}} \FunctionTok{select}\NormalTok{(condition, }\AttributeTok{estimate\_h =}\NormalTok{ estimate),}
\NormalTok{    ate\_model\_results[[out]]}\SpecialCharTok{$}\NormalTok{l }\SpecialCharTok{|\textgreater{}} \FunctionTok{select}\NormalTok{(condition, }\AttributeTok{estimate\_l =}\NormalTok{ estimate),}
    \AttributeTok{by =} \StringTok{"condition"}
\NormalTok{  ) }\SpecialCharTok{|\textgreater{}} \FunctionTok{mutate}\NormalTok{(}\AttributeTok{outcome =}\NormalTok{ outcome\_label\_map[out])}
\NormalTok{\}) }\SpecialCharTok{|\textgreater{}}
  \FunctionTok{ggplot}\NormalTok{(}\FunctionTok{aes}\NormalTok{(}\AttributeTok{x =}\NormalTok{ estimate\_h, }\AttributeTok{y =}\NormalTok{ estimate\_l)) }\SpecialCharTok{+}
  \FunctionTok{geom\_abline}\NormalTok{(}\AttributeTok{slope =} \DecValTok{1}\NormalTok{, }\AttributeTok{intercept =} \DecValTok{0}\NormalTok{, }\AttributeTok{linetype =} \StringTok{"dashed"}\NormalTok{,}
              \AttributeTok{color =} \StringTok{"grey60"}\NormalTok{, }\AttributeTok{linewidth =} \FloatTok{0.4}\NormalTok{, }\AttributeTok{alpha =} \FloatTok{0.35}\NormalTok{) }\SpecialCharTok{+}
  \FunctionTok{geom\_hline}\NormalTok{(}\AttributeTok{yintercept =} \DecValTok{0}\NormalTok{, }\AttributeTok{linetype =} \StringTok{"dotted"}\NormalTok{, }\AttributeTok{color =} \StringTok{"grey80"}\NormalTok{) }\SpecialCharTok{+}
  \FunctionTok{geom\_vline}\NormalTok{(}\AttributeTok{xintercept =} \DecValTok{0}\NormalTok{, }\AttributeTok{linetype =} \StringTok{"dotted"}\NormalTok{, }\AttributeTok{color =} \StringTok{"grey80"}\NormalTok{) }\SpecialCharTok{+}
  \FunctionTok{geom\_smooth}\NormalTok{(}\AttributeTok{method =} \StringTok{"lm"}\NormalTok{, }\AttributeTok{se =} \ConstantTok{FALSE}\NormalTok{, }\AttributeTok{color =} \StringTok{"\#2980B9"}\NormalTok{, }\AttributeTok{linewidth =} \FloatTok{0.6}\NormalTok{) }\SpecialCharTok{+}
  \FunctionTok{geom\_point}\NormalTok{(}\AttributeTok{size =} \FloatTok{1.5}\NormalTok{, }\AttributeTok{alpha =} \FloatTok{0.7}\NormalTok{, }\AttributeTok{color =} \StringTok{"grey20"}\NormalTok{) }\SpecialCharTok{+}
  \FunctionTok{facet\_wrap}\NormalTok{(}\SpecialCharTok{\textasciitilde{}}\NormalTok{ outcome, }\AttributeTok{scales =} \StringTok{"free"}\NormalTok{, }\AttributeTok{ncol =} \DecValTok{3}\NormalTok{) }\SpecialCharTok{+}
  \FunctionTok{labs}\NormalTok{(}\AttributeTok{x =} \StringTok{"Human ATE"}\NormalTok{, }\AttributeTok{y =} \StringTok{"LLM clone ATE"}\NormalTok{) }\SpecialCharTok{+}
\NormalTok{  plot\_theme}
\end{Highlighting}
\end{Shaded}

\begin{figure}[H]

\centering{

\pandocbounded{\includegraphics[keepaspectratio]{preregistration_files/figure-pdf/fig-rq1-all-1.pdf}}

}

\caption{\label{fig-rq1-all}Human vs.~LLM clone ATEs for all continuous
outcomes. Each point is one of the 20 interventions; axes are
free-scaled across panels. Blue line: OLS regression through the panel's
points. Faint dashed line: perfect calibration (slope = 1).}

\end{figure}%

\subsubsection{Response Distributions}\label{response-distributions}

Even if clones match the average treatment effect for each intervention,
they may diverge from humans in how responses are distributed: two
groups with the same mean can differ in spread, skewness, or tail shape.
We ask whether clones reproduce the full response distribution in the
control condition --- not just the mean --- across all continuous
outcomes. (\texttt{newsletter\_signup} is excluded because kernel
density estimates are not meaningful for a binary outcome.)

We compare distributions using the same three metrics defined in the
methods section, reported in increasing order of strictness. The
\textbf{variance ratio} (clone / human variance) is the least demanding
criterion --- it checks only whether spread is calibrated, ignoring
shape. The \textbf{overlapping coefficient (OVL)} raises the bar to
overall shape similarity: it is the proportion of area shared between
the two kernel density estimates, equal to 1 when distributions are
identical and approaching 0 when they barely overlap. Intuitively, OVL
is the overlap area between two smoothed histograms, expressed as a
fraction of their total combined area. The \textbf{Kolmogorov--Smirnov
D-statistic} measures the largest vertical gap between the two
cumulative distributions at any point --- imagine sliding along the
x-axis and tracking how far apart the two curves are; the `D' statistic
is the biggest gap. A value of 0 means the two distributions are
indistinguishable everywhere, while larger values flag that at some
response value, the two groups diverge sharply.

\begin{Shaded}
\begin{Highlighting}[]
\NormalTok{control\_val }\OtherTok{\textless{}{-}} \FunctionTok{levels}\NormalTok{(human\_data}\SpecialCharTok{$}\NormalTok{condition)[}\DecValTok{1}\NormalTok{]}

\NormalTok{dist\_control }\OtherTok{\textless{}{-}} \FunctionTok{map\_dfr}\NormalTok{(outcomes\_illustrative, }\ControlFlowTok{function}\NormalTok{(out) \{}
  \FunctionTok{compare\_distributions}\NormalTok{(human\_data\_1, llm\_data, out, control\_val) }\SpecialCharTok{|\textgreater{}}
    \FunctionTok{mutate}\NormalTok{(}\AttributeTok{outcome =}\NormalTok{ out)}
\NormalTok{\}) }\SpecialCharTok{|\textgreater{}}
  \FunctionTok{mutate}\NormalTok{(}
    \AttributeTok{tier =} \FunctionTok{case\_when}\NormalTok{(}
\NormalTok{      outcome }\SpecialCharTok{==}\NormalTok{ outcomes\_primary     }\SpecialCharTok{\textasciitilde{}} \StringTok{"Primary"}\NormalTok{,}
\NormalTok{      outcome }\SpecialCharTok{\%in\%}\NormalTok{ outcomes\_secondary }\SpecialCharTok{\textasciitilde{}} \StringTok{"Secondary"}\NormalTok{,}
      \ConstantTok{TRUE}                            \SpecialCharTok{\textasciitilde{}} \StringTok{"Tertiary"}
\NormalTok{    ),}
    \AttributeTok{outcome\_label =}\NormalTok{ outcome\_label\_map[outcome]}
\NormalTok{  )}

\CommentTok{\# Human 1 vs. Human 2 baseline ceiling for distribution comparison}
\NormalTok{dist\_control\_h2 }\OtherTok{\textless{}{-}} \FunctionTok{map\_dfr}\NormalTok{(outcomes\_illustrative, }\ControlFlowTok{function}\NormalTok{(out) \{}
  \FunctionTok{compare\_distributions}\NormalTok{(human\_data\_1, human\_data\_2, out, control\_val) }\SpecialCharTok{|\textgreater{}}
    \FunctionTok{mutate}\NormalTok{(}\AttributeTok{outcome =}\NormalTok{ out)}
\NormalTok{\}) }\SpecialCharTok{|\textgreater{}}
  \FunctionTok{mutate}\NormalTok{(}
    \AttributeTok{tier =} \FunctionTok{case\_when}\NormalTok{(}
\NormalTok{      outcome }\SpecialCharTok{==}\NormalTok{ outcomes\_primary     }\SpecialCharTok{\textasciitilde{}} \StringTok{"Primary"}\NormalTok{,}
\NormalTok{      outcome }\SpecialCharTok{\%in\%}\NormalTok{ outcomes\_secondary }\SpecialCharTok{\textasciitilde{}} \StringTok{"Secondary"}\NormalTok{,}
      \ConstantTok{TRUE}                            \SpecialCharTok{\textasciitilde{}} \StringTok{"Tertiary"}
\NormalTok{    ),}
    \AttributeTok{outcome\_label =}\NormalTok{ outcome\_label\_map[outcome]}
\NormalTok{  )}
\end{Highlighting}
\end{Shaded}

\begin{Shaded}
\begin{Highlighting}[]
\NormalTok{source\_colors }\OtherTok{\textless{}{-}} \FunctionTok{c}\NormalTok{(}\StringTok{"Human"} \OtherTok{=} \StringTok{"grey30"}\NormalTok{, }\StringTok{"LLM clone"} \OtherTok{=} \FunctionTok{met.brewer}\NormalTok{(}\StringTok{"Juarez"}\NormalTok{, }\AttributeTok{n =} \DecValTok{10}\NormalTok{)[}\DecValTok{1}\NormalTok{])}

\FunctionTok{map\_dfr}\NormalTok{(outcomes\_illustrative, }\ControlFlowTok{function}\NormalTok{(out) \{}
  \FunctionTok{bind\_rows}\NormalTok{(}
\NormalTok{    human\_data\_1 }\SpecialCharTok{|\textgreater{}} \FunctionTok{filter}\NormalTok{(condition }\SpecialCharTok{==}\NormalTok{ control\_val) }\SpecialCharTok{|\textgreater{}}
      \FunctionTok{select}\NormalTok{(}\AttributeTok{value =} \FunctionTok{all\_of}\NormalTok{(out)) }\SpecialCharTok{|\textgreater{}}
      \FunctionTok{mutate}\NormalTok{(}\AttributeTok{source =} \StringTok{"Human"}\NormalTok{,     }\AttributeTok{outcome =}\NormalTok{ outcome\_label\_map[out]),}
\NormalTok{    llm\_data   }\SpecialCharTok{|\textgreater{}} \FunctionTok{filter}\NormalTok{(condition }\SpecialCharTok{==}\NormalTok{ control\_val) }\SpecialCharTok{|\textgreater{}}
      \FunctionTok{select}\NormalTok{(}\AttributeTok{value =} \FunctionTok{all\_of}\NormalTok{(out)) }\SpecialCharTok{|\textgreater{}}
      \FunctionTok{mutate}\NormalTok{(}\AttributeTok{source =} \StringTok{"LLM clone"}\NormalTok{, }\AttributeTok{outcome =}\NormalTok{ outcome\_label\_map[out])}
\NormalTok{  )}
\NormalTok{\}) }\SpecialCharTok{|\textgreater{}}
  \FunctionTok{filter}\NormalTok{(}\SpecialCharTok{!}\FunctionTok{is.na}\NormalTok{(value)) }\SpecialCharTok{|\textgreater{}}
  \FunctionTok{ggplot}\NormalTok{(}\FunctionTok{aes}\NormalTok{(}\AttributeTok{x =}\NormalTok{ value, }\AttributeTok{fill =}\NormalTok{ source, }\AttributeTok{color =}\NormalTok{ source)) }\SpecialCharTok{+}
  \FunctionTok{geom\_density}\NormalTok{(}\AttributeTok{alpha =} \FloatTok{0.35}\NormalTok{, }\AttributeTok{linewidth =} \FloatTok{0.4}\NormalTok{) }\SpecialCharTok{+}
  \FunctionTok{scale\_fill\_manual}\NormalTok{(}\AttributeTok{values  =}\NormalTok{ source\_colors, }\AttributeTok{name =} \ConstantTok{NULL}\NormalTok{) }\SpecialCharTok{+}
  \FunctionTok{scale\_color\_manual}\NormalTok{(}\AttributeTok{values =}\NormalTok{ source\_colors, }\AttributeTok{name =} \ConstantTok{NULL}\NormalTok{) }\SpecialCharTok{+}
  \FunctionTok{facet\_wrap}\NormalTok{(}\SpecialCharTok{\textasciitilde{}}\NormalTok{ outcome, }\AttributeTok{scales =} \StringTok{"free"}\NormalTok{, }\AttributeTok{ncol =} \DecValTok{2}\NormalTok{) }\SpecialCharTok{+}
  \FunctionTok{labs}\NormalTok{(}\AttributeTok{x =} \StringTok{"Response"}\NormalTok{, }\AttributeTok{y =} \StringTok{"Density"}\NormalTok{) }\SpecialCharTok{+}
\NormalTok{  plot\_theme }\SpecialCharTok{+}
  \FunctionTok{theme}\NormalTok{(}\AttributeTok{legend.position =} \StringTok{"top"}\NormalTok{)}
\end{Highlighting}
\end{Shaded}

\begin{figure}[H]

\centering{

\pandocbounded{\includegraphics[keepaspectratio]{preregistration_files/figure-pdf/fig-rq2-densities-1.pdf}}

}

\caption{\label{fig-rq2-densities}Response distributions in the control
condition, by outcome. Human participants shown in dark grey, LLM clones
in orange. Each panel covers the full range of that outcome's scale.}

\end{figure}%

\begin{Shaded}
\begin{Highlighting}[]
\NormalTok{dist\_control }\SpecialCharTok{|\textgreater{}}
  \FunctionTok{left\_join}\NormalTok{(}
\NormalTok{    dist\_control\_h2 }\SpecialCharTok{|\textgreater{}} \FunctionTok{select}\NormalTok{(outcome, }\AttributeTok{ovl\_h2 =}\NormalTok{ ovl, }\AttributeTok{ks\_d\_h2 =}\NormalTok{ ks\_d, }\AttributeTok{vr\_h2 =}\NormalTok{ variance\_ratio),}
    \AttributeTok{by =} \StringTok{"outcome"}
\NormalTok{  ) }\SpecialCharTok{|\textgreater{}}
  \FunctionTok{arrange}\NormalTok{(}\FunctionTok{match}\NormalTok{(tier, }\FunctionTok{c}\NormalTok{(}\StringTok{"Primary"}\NormalTok{, }\StringTok{"Secondary"}\NormalTok{, }\StringTok{"Tertiary"}\NormalTok{)), outcome\_label) }\SpecialCharTok{|\textgreater{}}
  \FunctionTok{select}\NormalTok{(}\AttributeTok{Tier =}\NormalTok{ tier, }\AttributeTok{Outcome =}\NormalTok{ outcome\_label,}
\NormalTok{         ovl, ks\_d, variance\_ratio, ovl\_h2, ks\_d\_h2, vr\_h2) }\SpecialCharTok{|\textgreater{}}
  \FunctionTok{gt}\NormalTok{(}\AttributeTok{groupname\_col =} \StringTok{"Tier"}\NormalTok{) }\SpecialCharTok{|\textgreater{}}
  \FunctionTok{cols\_label}\NormalTok{(}\AttributeTok{ovl =} \StringTok{"OVL"}\NormalTok{, }\AttributeTok{ks\_d =} \StringTok{"KS D"}\NormalTok{, }\AttributeTok{variance\_ratio =} \StringTok{"Var. ratio"}\NormalTok{,}
             \AttributeTok{ovl\_h2 =} \StringTok{"OVL"}\NormalTok{, }\AttributeTok{ks\_d\_h2 =} \StringTok{"KS D"}\NormalTok{, }\AttributeTok{vr\_h2 =} \StringTok{"Var. ratio"}\NormalTok{) }\SpecialCharTok{|\textgreater{}}
  \FunctionTok{tab\_spanner}\NormalTok{(}\AttributeTok{label =} \StringTok{"Human 1 vs. LLM"}\NormalTok{,       }\AttributeTok{columns =} \FunctionTok{c}\NormalTok{(ovl, ks\_d, variance\_ratio)) }\SpecialCharTok{|\textgreater{}}
  \FunctionTok{tab\_spanner}\NormalTok{(}\AttributeTok{label =} \StringTok{"Human 1 vs. Human 2"}\NormalTok{,   }\AttributeTok{columns =} \FunctionTok{c}\NormalTok{(ovl\_h2, ks\_d\_h2, vr\_h2)) }\SpecialCharTok{|\textgreater{}}
  \FunctionTok{fmt\_number}\NormalTok{(}\AttributeTok{columns =} \FunctionTok{c}\NormalTok{(ovl, ks\_d, variance\_ratio, ovl\_h2, ks\_d\_h2, vr\_h2), }\AttributeTok{decimals =} \DecValTok{3}\NormalTok{) }\SpecialCharTok{|\textgreater{}}
  \FunctionTok{tab\_style}\NormalTok{(}
    \AttributeTok{style     =} \FunctionTok{cell\_text}\NormalTok{(}\AttributeTok{weight =} \StringTok{"bold"}\NormalTok{),}
    \AttributeTok{locations =} \FunctionTok{cells\_row\_groups}\NormalTok{()}
\NormalTok{  )}
\end{Highlighting}
\end{Shaded}

\begin{table}

\caption{\label{tbl-rq2}Response distribution metrics in the control
condition, by outcome. OVL: overlapping coefficient (1 = identical). KS
D: Kolmogorov--Smirnov D-statistic (0 = identical CDFs). Variance ratio:
clone / human variance (1 = equal spread). Left columns: Human 1
vs.~LLM; right columns: Human 1 vs.~Human 2 (baseline ceiling).}

\centering{

\fontsize{12.0pt}{14.4pt}\selectfont
\begin{tabular*}{\linewidth}{@{\extracolsep{\fill}}lrrrrrr}
\toprule
 & \multicolumn{3}{c}{Human 1 vs. LLM} & \multicolumn{3}{c}{Human 1 vs. Human 2} \\ 
\cmidrule(lr){2-4} \cmidrule(lr){5-7}
Outcome & OVL & KS D & Var. ratio & OVL & KS D & Var. ratio \\ 
\midrule\addlinespace[2.5pt]
\multicolumn{7}{l}{{\bfseries Primary}} \\[2.5pt] 
\midrule\addlinespace[2.5pt]
Multidimensional trust (primary) & 0.956 & 0.041 & 1.062 & 0.934 & 0.058 & 1.021 \\ 
\midrule\addlinespace[2.5pt]
\multicolumn{7}{l}{{\bfseries Secondary}} \\[2.5pt] 
\midrule\addlinespace[2.5pt]
Funding perceptions & 0.897 & 0.034 & 1.053 & 0.889 & 0.042 & 1.066 \\ 
\midrule\addlinespace[2.5pt]
\multicolumn{7}{l}{{\bfseries Tertiary}} \\[2.5pt] 
\midrule\addlinespace[2.5pt]
Donation to AMS (\$) & 0.910 & 0.050 & 1.037 & 0.927 & 0.028 & 1.027 \\ 
General climate policy & 0.888 & 0.050 & 0.981 & 0.865 & 0.077 & 0.914 \\ 
\bottomrule
\end{tabular*}

}

\end{table}%

\subsection{2. Subgroup Effects}\label{subgroup-effects}

\subsubsection{Heterogeneity in Treatment
Effects}\label{heterogeneity-in-treatment-effects}

Interventions may work differently for different demographic groups ---
for example, a message about scientific consensus might shift trust more
strongly among Republicans than Democrats, or more among younger than
older respondents. We ask whether clones reproduce these subgroup
differences, or whether they produce a one-size-fits-all response that
averages across demographic groups.

For each of six categorical moderators (gender, age band, race,
education, income, and partisan identity), we estimate condition ×
moderator interactions using \texttt{run\_moderator\_model()} separately
on human and clone data for all continuous outcomes. Each interaction
term captures how much more (or less) effective a given intervention is
for one demographic group compared to the reference group. We compare
the full set of interaction estimates between humans and clones using
\texttt{compare\_estimates()}, reporting the six metrics in the same
least-to-most-strict order as in Section 1 (directional agreement
through TOST equivalence). We first report metrics pooled across all
illustrative outcomes and moderators, then break down by outcome.

\begin{Shaded}
\begin{Highlighting}[]
\CommentTok{\# Nested: outcome → moderator → model (illustrative subset)}
\NormalTok{mod\_results\_all\_h }\OtherTok{\textless{}{-}} \FunctionTok{map}\NormalTok{(}\FunctionTok{set\_names}\NormalTok{(outcomes\_illustrative), }\ControlFlowTok{function}\NormalTok{(out) \{}
  \FunctionTok{map}\NormalTok{(}\FunctionTok{set\_names}\NormalTok{(moderators\_cat), }\ControlFlowTok{function}\NormalTok{(mod) \{}
    \FunctionTok{run\_moderator\_model}\NormalTok{(human\_data\_1, }\AttributeTok{outcome =}\NormalTok{ out, }\AttributeTok{moderator =}\NormalTok{ mod)}
\NormalTok{  \})}
\NormalTok{\})}
\NormalTok{mod\_results\_all\_l }\OtherTok{\textless{}{-}} \FunctionTok{map}\NormalTok{(}\FunctionTok{set\_names}\NormalTok{(outcomes\_illustrative), }\ControlFlowTok{function}\NormalTok{(out) \{}
  \FunctionTok{map}\NormalTok{(}\FunctionTok{set\_names}\NormalTok{(moderators\_cat), }\ControlFlowTok{function}\NormalTok{(mod) \{}
    \FunctionTok{run\_moderator\_model}\NormalTok{(llm\_data, }\AttributeTok{outcome =}\NormalTok{ out, }\AttributeTok{moderator =}\NormalTok{ mod)}
\NormalTok{  \})}
\NormalTok{\})}
\NormalTok{mod\_results\_all\_h2 }\OtherTok{\textless{}{-}} \FunctionTok{map}\NormalTok{(}\FunctionTok{set\_names}\NormalTok{(outcomes\_illustrative), }\ControlFlowTok{function}\NormalTok{(out) \{}
  \FunctionTok{map}\NormalTok{(}\FunctionTok{set\_names}\NormalTok{(moderators\_cat), }\ControlFlowTok{function}\NormalTok{(mod) \{}
    \FunctionTok{run\_moderator\_model}\NormalTok{(human\_data\_2, }\AttributeTok{outcome =}\NormalTok{ out, }\AttributeTok{moderator =}\NormalTok{ mod)}
\NormalTok{  \})}
\NormalTok{\})}

\CommentTok{\# Primary{-}outcome references for the scatter plot below}
\NormalTok{mod\_results\_h  }\OtherTok{\textless{}{-}}\NormalTok{ mod\_results\_all\_h[[}\StringTok{"trust\_multidimensional"}\NormalTok{]]}
\NormalTok{mod\_results\_l  }\OtherTok{\textless{}{-}}\NormalTok{ mod\_results\_all\_l[[}\StringTok{"trust\_multidimensional"}\NormalTok{]]}
\NormalTok{mod\_results\_h2 }\OtherTok{\textless{}{-}}\NormalTok{ mod\_results\_all\_h2[[}\StringTok{"trust\_multidimensional"}\NormalTok{]]}
\end{Highlighting}
\end{Shaded}

\begin{Shaded}
\begin{Highlighting}[]
\CommentTok{\# All interaction estimate pairs (outcome × moderator × condition × level)}
\NormalTok{rq3\_pooled\_interactions }\OtherTok{\textless{}{-}} \FunctionTok{map\_dfr}\NormalTok{(outcomes\_illustrative, }\ControlFlowTok{function}\NormalTok{(out) \{}
  \FunctionTok{map\_dfr}\NormalTok{(moderators\_cat, }\ControlFlowTok{function}\NormalTok{(mod) \{}
    \FunctionTok{inner\_join}\NormalTok{(}
\NormalTok{      mod\_results\_all\_h[[out]][[mod]]}\SpecialCharTok{$}\NormalTok{interaction\_effects }\SpecialCharTok{|\textgreater{}}
        \FunctionTok{select}\NormalTok{(condition, moderator\_level, }\AttributeTok{estimate\_h =}\NormalTok{ estimate),}
\NormalTok{      mod\_results\_all\_l[[out]][[mod]]}\SpecialCharTok{$}\NormalTok{interaction\_effects }\SpecialCharTok{|\textgreater{}}
        \FunctionTok{select}\NormalTok{(condition, moderator\_level, }\AttributeTok{estimate\_l =}\NormalTok{ estimate),}
      \AttributeTok{by =} \FunctionTok{c}\NormalTok{(}\StringTok{"condition"}\NormalTok{, }\StringTok{"moderator\_level"}\NormalTok{)}
\NormalTok{    ) }\SpecialCharTok{|\textgreater{}} \FunctionTok{mutate}\NormalTok{(}\AttributeTok{outcome =}\NormalTok{ out, }\AttributeTok{moderator =}\NormalTok{ mod)}
\NormalTok{  \})}
\NormalTok{\})}

\CommentTok{\# Pooled headline (scale{-}invariant metrics only; RMSE omitted across mixed scales)}
\NormalTok{rq3\_pooled\_row }\OtherTok{\textless{}{-}}\NormalTok{ rq3\_pooled\_interactions }\SpecialCharTok{|\textgreater{}}
  \FunctionTok{summarise}\NormalTok{(}
    \AttributeTok{spearman\_rho    =} \FunctionTok{cor}\NormalTok{(estimate\_h, estimate\_l, }\AttributeTok{method =} \StringTok{"spearman"}\NormalTok{,}
                          \AttributeTok{use =} \StringTok{"pairwise.complete.obs"}\NormalTok{),}
    \AttributeTok{pearson\_r       =} \FunctionTok{cor}\NormalTok{(estimate\_h, estimate\_l, }\AttributeTok{use =} \StringTok{"pairwise.complete.obs"}\NormalTok{),}
    \AttributeTok{directional\_pct =} \FunctionTok{mean}\NormalTok{(}\FunctionTok{sign}\NormalTok{(estimate\_h) }\SpecialCharTok{==} \FunctionTok{sign}\NormalTok{(estimate\_l), }\AttributeTok{na.rm =} \ConstantTok{TRUE}\NormalTok{) }\SpecialCharTok{*} \DecValTok{100}
\NormalTok{  ) }\SpecialCharTok{|\textgreater{}}
  \FunctionTok{mutate}\NormalTok{(}\AttributeTok{outcome\_label =} \StringTok{"All outcomes"}\NormalTok{, }\AttributeTok{tier =} \StringTok{"All outcomes"}\NormalTok{)}

\CommentTok{\# Per{-}outcome × per{-}moderator metrics}
\NormalTok{rq3\_metrics\_by\_outcome }\OtherTok{\textless{}{-}} \FunctionTok{map\_dfr}\NormalTok{(outcomes\_illustrative, }\ControlFlowTok{function}\NormalTok{(out) \{}
  \FunctionTok{map\_dfr}\NormalTok{(moderators\_cat, }\ControlFlowTok{function}\NormalTok{(mod) \{}
    \FunctionTok{compare\_estimates}\NormalTok{(}
\NormalTok{      mod\_results\_all\_h[[out]][[mod]]}\SpecialCharTok{$}\NormalTok{interaction\_effects,}
\NormalTok{      mod\_results\_all\_l[[out]][[mod]]}\SpecialCharTok{$}\NormalTok{interaction\_effects,}
      \AttributeTok{join\_by =} \FunctionTok{c}\NormalTok{(}\StringTok{"condition"}\NormalTok{, }\StringTok{"moderator\_level"}\NormalTok{)}
\NormalTok{    ) }\SpecialCharTok{|\textgreater{}} \FunctionTok{mutate}\NormalTok{(}
      \AttributeTok{outcome       =}\NormalTok{ out,}
      \AttributeTok{moderator     =}\NormalTok{ mod,}
      \AttributeTok{outcome\_label =}\NormalTok{ outcome\_label\_map[out],}
      \AttributeTok{tier =} \FunctionTok{case\_when}\NormalTok{(}
\NormalTok{        out }\SpecialCharTok{==}\NormalTok{ outcomes\_primary     }\SpecialCharTok{\textasciitilde{}} \StringTok{"Primary"}\NormalTok{,}
\NormalTok{        out }\SpecialCharTok{\%in\%}\NormalTok{ outcomes\_secondary }\SpecialCharTok{\textasciitilde{}} \StringTok{"Secondary"}\NormalTok{,}
        \ConstantTok{TRUE}                        \SpecialCharTok{\textasciitilde{}} \StringTok{"Tertiary"}
\NormalTok{      )}
\NormalTok{    )}
\NormalTok{  \})}
\NormalTok{\})}

\CommentTok{\# Human 1 vs. Human 2 baseline ceiling for subgroup effects}
\NormalTok{rq3\_pooled\_interactions\_h2 }\OtherTok{\textless{}{-}} \FunctionTok{map\_dfr}\NormalTok{(outcomes\_illustrative, }\ControlFlowTok{function}\NormalTok{(out) \{}
  \FunctionTok{map\_dfr}\NormalTok{(moderators\_cat, }\ControlFlowTok{function}\NormalTok{(mod) \{}
    \FunctionTok{inner\_join}\NormalTok{(}
\NormalTok{      mod\_results\_all\_h[[out]][[mod]]}\SpecialCharTok{$}\NormalTok{interaction\_effects }\SpecialCharTok{|\textgreater{}}
        \FunctionTok{select}\NormalTok{(condition, moderator\_level, }\AttributeTok{estimate\_h =}\NormalTok{ estimate),}
\NormalTok{      mod\_results\_all\_h2[[out]][[mod]]}\SpecialCharTok{$}\NormalTok{interaction\_effects }\SpecialCharTok{|\textgreater{}}
        \FunctionTok{select}\NormalTok{(condition, moderator\_level, }\AttributeTok{estimate\_h2 =}\NormalTok{ estimate),}
      \AttributeTok{by =} \FunctionTok{c}\NormalTok{(}\StringTok{"condition"}\NormalTok{, }\StringTok{"moderator\_level"}\NormalTok{)}
\NormalTok{    ) }\SpecialCharTok{|\textgreater{}} \FunctionTok{mutate}\NormalTok{(}\AttributeTok{outcome =}\NormalTok{ out, }\AttributeTok{moderator =}\NormalTok{ mod)}
\NormalTok{  \})}
\NormalTok{\})}

\NormalTok{rq3\_pooled\_row\_h2 }\OtherTok{\textless{}{-}}\NormalTok{ rq3\_pooled\_interactions\_h2 }\SpecialCharTok{|\textgreater{}}
  \FunctionTok{summarise}\NormalTok{(}
    \AttributeTok{spearman\_rho    =} \FunctionTok{cor}\NormalTok{(estimate\_h, estimate\_h2, }\AttributeTok{method =} \StringTok{"spearman"}\NormalTok{,}
                          \AttributeTok{use =} \StringTok{"pairwise.complete.obs"}\NormalTok{),}
    \AttributeTok{pearson\_r       =} \FunctionTok{cor}\NormalTok{(estimate\_h, estimate\_h2, }\AttributeTok{use =} \StringTok{"pairwise.complete.obs"}\NormalTok{),}
    \AttributeTok{directional\_pct =} \FunctionTok{mean}\NormalTok{(}\FunctionTok{sign}\NormalTok{(estimate\_h) }\SpecialCharTok{==} \FunctionTok{sign}\NormalTok{(estimate\_h2), }\AttributeTok{na.rm =} \ConstantTok{TRUE}\NormalTok{) }\SpecialCharTok{*} \DecValTok{100}
\NormalTok{  ) }\SpecialCharTok{|\textgreater{}}
  \FunctionTok{mutate}\NormalTok{(}\AttributeTok{outcome\_label =} \StringTok{"All outcomes"}\NormalTok{, }\AttributeTok{tier =} \StringTok{"All outcomes"}\NormalTok{)}

\NormalTok{rq3\_metrics\_by\_outcome\_h2 }\OtherTok{\textless{}{-}} \FunctionTok{map\_dfr}\NormalTok{(outcomes\_illustrative, }\ControlFlowTok{function}\NormalTok{(out) \{}
  \FunctionTok{map\_dfr}\NormalTok{(moderators\_cat, }\ControlFlowTok{function}\NormalTok{(mod) \{}
    \FunctionTok{compare\_estimates}\NormalTok{(}
\NormalTok{      mod\_results\_all\_h[[out]][[mod]]}\SpecialCharTok{$}\NormalTok{interaction\_effects,}
\NormalTok{      mod\_results\_all\_h2[[out]][[mod]]}\SpecialCharTok{$}\NormalTok{interaction\_effects,}
      \AttributeTok{join\_by =} \FunctionTok{c}\NormalTok{(}\StringTok{"condition"}\NormalTok{, }\StringTok{"moderator\_level"}\NormalTok{)}
\NormalTok{    ) }\SpecialCharTok{|\textgreater{}} \FunctionTok{mutate}\NormalTok{(}
      \AttributeTok{outcome       =}\NormalTok{ out,}
      \AttributeTok{moderator     =}\NormalTok{ mod,}
      \AttributeTok{outcome\_label =}\NormalTok{ outcome\_label\_map[out],}
      \AttributeTok{tier =} \FunctionTok{case\_when}\NormalTok{(}
\NormalTok{        out }\SpecialCharTok{==}\NormalTok{ outcomes\_primary     }\SpecialCharTok{\textasciitilde{}} \StringTok{"Primary"}\NormalTok{,}
\NormalTok{        out }\SpecialCharTok{\%in\%}\NormalTok{ outcomes\_secondary }\SpecialCharTok{\textasciitilde{}} \StringTok{"Secondary"}\NormalTok{,}
        \ConstantTok{TRUE}                        \SpecialCharTok{\textasciitilde{}} \StringTok{"Tertiary"}
\NormalTok{      )}
\NormalTok{    )}
\NormalTok{  \})}
\NormalTok{\})}
\end{Highlighting}
\end{Shaded}

Pooled across all outcomes and moderators, Human 1 vs.~LLM: Spearman ρ =
-0.069, Pearson r = -0.116, directional agreement = 48.8\%. Human 1
vs.~Human 2 (baseline ceiling): Spearman ρ = 0.034, Pearson r = 0.022,
directional agreement = 51.6\%.

\begin{Shaded}
\begin{Highlighting}[]
\NormalTok{rq3\_metrics\_by\_outcome }\SpecialCharTok{|\textgreater{}}
  \FunctionTok{rename\_with}\NormalTok{(}\SpecialCharTok{\textasciitilde{}} \FunctionTok{paste0}\NormalTok{(.x, }\StringTok{"\_llm"}\NormalTok{),}
              \AttributeTok{.cols =} \FunctionTok{c}\NormalTok{(spearman\_rho, pearson\_r, rmse,}
\NormalTok{                        directional\_pct, inferential\_pct, tost\_pct)) }\SpecialCharTok{|\textgreater{}}
  \FunctionTok{left\_join}\NormalTok{(}
\NormalTok{    rq3\_metrics\_by\_outcome\_h2 }\SpecialCharTok{|\textgreater{}}
      \FunctionTok{rename\_with}\NormalTok{(}\SpecialCharTok{\textasciitilde{}} \FunctionTok{paste0}\NormalTok{(.x, }\StringTok{"\_h2"}\NormalTok{),}
                  \AttributeTok{.cols =} \FunctionTok{c}\NormalTok{(spearman\_rho, pearson\_r, rmse,}
\NormalTok{                             directional\_pct, inferential\_pct, tost\_pct)) }\SpecialCharTok{|\textgreater{}}
      \FunctionTok{select}\NormalTok{(outcome, moderator, }\FunctionTok{ends\_with}\NormalTok{(}\StringTok{"\_h2"}\NormalTok{)),}
    \AttributeTok{by =} \FunctionTok{c}\NormalTok{(}\StringTok{"outcome"}\NormalTok{, }\StringTok{"moderator"}\NormalTok{)}
\NormalTok{  ) }\SpecialCharTok{|\textgreater{}}
  \FunctionTok{arrange}\NormalTok{(}\FunctionTok{match}\NormalTok{(tier, }\FunctionTok{c}\NormalTok{(}\StringTok{"Primary"}\NormalTok{, }\StringTok{"Secondary"}\NormalTok{, }\StringTok{"Tertiary"}\NormalTok{)), outcome\_label, moderator) }\SpecialCharTok{|\textgreater{}}
  \FunctionTok{select}\NormalTok{(}\AttributeTok{Tier =}\NormalTok{ tier, }\AttributeTok{Outcome =}\NormalTok{ outcome\_label, }\AttributeTok{Moderator =}\NormalTok{ moderator,}
\NormalTok{         spearman\_rho\_llm, pearson\_r\_llm, rmse\_llm,}
\NormalTok{         directional\_pct\_llm, inferential\_pct\_llm, tost\_pct\_llm,}
\NormalTok{         spearman\_rho\_h2, pearson\_r\_h2, rmse\_h2,}
\NormalTok{         directional\_pct\_h2, inferential\_pct\_h2, tost\_pct\_h2) }\SpecialCharTok{|\textgreater{}}
  \FunctionTok{gt}\NormalTok{(}\AttributeTok{groupname\_col =} \StringTok{"Tier"}\NormalTok{) }\SpecialCharTok{|\textgreater{}}
  \FunctionTok{cols\_label}\NormalTok{(}
    \AttributeTok{spearman\_rho\_llm    =} \StringTok{"Spearman ρ"}\NormalTok{, }\AttributeTok{pearson\_r\_llm    =} \StringTok{"Pearson r"}\NormalTok{, }\AttributeTok{rmse\_llm    =} \StringTok{"RMSE"}\NormalTok{,}
    \AttributeTok{directional\_pct\_llm =} \StringTok{"Dir. (\%)"}\NormalTok{,   }\AttributeTok{inferential\_pct\_llm =} \StringTok{"Infer. (\%)"}\NormalTok{, }\AttributeTok{tost\_pct\_llm =} \StringTok{"TOST (\%)"}\NormalTok{,}
    \AttributeTok{spearman\_rho\_h2     =} \StringTok{"Spearman ρ"}\NormalTok{, }\AttributeTok{pearson\_r\_h2     =} \StringTok{"Pearson r"}\NormalTok{, }\AttributeTok{rmse\_h2     =} \StringTok{"RMSE"}\NormalTok{,}
    \AttributeTok{directional\_pct\_h2  =} \StringTok{"Dir. (\%)"}\NormalTok{,   }\AttributeTok{inferential\_pct\_h2  =} \StringTok{"Infer. (\%)"}\NormalTok{, }\AttributeTok{tost\_pct\_h2  =} \StringTok{"TOST (\%)"}
\NormalTok{  ) }\SpecialCharTok{|\textgreater{}}
  \FunctionTok{tab\_spanner}\NormalTok{(}\AttributeTok{label =} \StringTok{"Human 1 vs. LLM"}\NormalTok{,     }\AttributeTok{columns =} \FunctionTok{ends\_with}\NormalTok{(}\StringTok{"\_llm"}\NormalTok{)) }\SpecialCharTok{|\textgreater{}}
  \FunctionTok{tab\_spanner}\NormalTok{(}\AttributeTok{label =} \StringTok{"Human 1 vs. Human 2"}\NormalTok{, }\AttributeTok{columns =} \FunctionTok{ends\_with}\NormalTok{(}\StringTok{"\_h2"}\NormalTok{)) }\SpecialCharTok{|\textgreater{}}
  \FunctionTok{fmt\_number}\NormalTok{(}\AttributeTok{columns =} \FunctionTok{c}\NormalTok{(spearman\_rho\_llm, pearson\_r\_llm, rmse\_llm,}
\NormalTok{                          spearman\_rho\_h2,  pearson\_r\_h2,  rmse\_h2), }\AttributeTok{decimals =} \DecValTok{3}\NormalTok{) }\SpecialCharTok{|\textgreater{}}
  \FunctionTok{fmt\_number}\NormalTok{(}\AttributeTok{columns =} \FunctionTok{c}\NormalTok{(directional\_pct\_llm, inferential\_pct\_llm, tost\_pct\_llm,}
\NormalTok{                          directional\_pct\_h2,  inferential\_pct\_h2,  tost\_pct\_h2), }\AttributeTok{decimals =} \DecValTok{1}\NormalTok{) }\SpecialCharTok{|\textgreater{}}
  \FunctionTok{tab\_style}\NormalTok{(}
    \AttributeTok{style     =} \FunctionTok{cell\_text}\NormalTok{(}\AttributeTok{weight =} \StringTok{"bold"}\NormalTok{),}
    \AttributeTok{locations =} \FunctionTok{cells\_row\_groups}\NormalTok{()}
\NormalTok{  )}
\end{Highlighting}
\end{Shaded}

\begin{table}

\caption{\label{tbl-rq3}Interaction estimate comparison metrics by
outcome and moderator. Outcomes are grouped by tier; rows are
moderators. Spearman ρ: rank correlation across all condition ×
moderator level pairs. RMSE is in outcome units. Left columns: Human 1
vs.~LLM; right columns: Human 1 vs.~Human 2 (baseline ceiling).}

\centering{

\fontsize{12.0pt}{14.4pt}\selectfont
\begin{tabular*}{\linewidth}{@{\extracolsep{\fill}}llrrrrrrrrrrrr}
\toprule
 &  & \multicolumn{6}{c}{Human 1 vs. LLM} & \multicolumn{6}{c}{Human 1 vs. Human 2} \\ 
\cmidrule(lr){3-8} \cmidrule(lr){9-14}
Outcome & Moderator & Spearman ρ & Pearson r & RMSE & Dir. (\%) & Infer. (\%) & TOST (\%) & Spearman ρ & Pearson r & RMSE & Dir. (\%) & Infer. (\%) & TOST (\%) \\ 
\midrule\addlinespace[2.5pt]
\multicolumn{14}{l}{{\bfseries Primary}} \\[2.5pt] 
\midrule\addlinespace[2.5pt]
Multidimensional trust (primary) & age\_band & 0.018 & 0.029 & 1.744 & 62.0 & 100.0 & 0.0 & 0.050 & 0.064 & 2.385 & 44.0 & 100.0 & 0.0 \\ 
Multidimensional trust (primary) & education & -0.192 & -0.168 & 1.868 & 40.0 & 100.0 & 0.0 & 0.136 & 0.135 & 1.957 & 53.0 & 100.0 & 0.0 \\ 
Multidimensional trust (primary) & gender & -0.173 & -0.135 & 1.349 & 47.5 & 100.0 & 0.0 & 0.016 & 0.026 & 1.279 & 55.0 & 100.0 & 0.0 \\ 
Multidimensional trust (primary) & income & 0.035 & 0.029 & 1.553 & 56.2 & 100.0 & 0.0 & 0.081 & 0.062 & 1.783 & 62.5 & 100.0 & 0.0 \\ 
Multidimensional trust (primary) & party & -0.394 & -0.356 & 1.698 & 51.7 & 100.0 & 0.0 & -0.344 & -0.348 & 2.140 & 36.7 & 100.0 & 0.0 \\ 
Multidimensional trust (primary) & race & 0.019 & 0.054 & 2.061 & 52.5 & 100.0 & 0.0 & -0.026 & 0.011 & 1.909 & 66.2 & 100.0 & 0.0 \\ 
\midrule\addlinespace[2.5pt]
\multicolumn{14}{l}{{\bfseries Secondary}} \\[2.5pt] 
\midrule\addlinespace[2.5pt]
Funding perceptions & age\_band & -0.209 & -0.214 & 6.749 & 42.0 & 100.0 & 0.0 & 0.023 & 0.041 & 6.435 & 49.0 & 100.0 & 0.0 \\ 
Funding perceptions & education & 0.155 & 0.100 & 5.853 & 58.0 & 100.0 & 0.0 & 0.102 & 0.107 & 7.627 & 48.0 & 100.0 & 0.0 \\ 
Funding perceptions & gender & -0.116 & -0.060 & 6.234 & 32.5 & 100.0 & 0.0 & -0.278 & -0.315 & 5.897 & 52.5 & 100.0 & 0.0 \\ 
Funding perceptions & income & -0.076 & -0.075 & 5.318 & 53.8 & 98.8 & 0.0 & 0.052 & 0.035 & 7.827 & 43.8 & 100.0 & 0.0 \\ 
Funding perceptions & party & 0.243 & 0.277 & 6.138 & 50.0 & 100.0 & 0.0 & -0.109 & -0.168 & 6.456 & 50.0 & 100.0 & 0.0 \\ 
Funding perceptions & race & -0.057 & -0.067 & 6.505 & 47.5 & 100.0 & 0.0 & -0.030 & -0.051 & 5.912 & 52.5 & 100.0 & 0.0 \\ 
\midrule\addlinespace[2.5pt]
\multicolumn{14}{l}{{\bfseries Tertiary}} \\[2.5pt] 
\midrule\addlinespace[2.5pt]
Donation to AMS (\$) & age\_band & -0.067 & -0.091 & 0.689 & 46.0 & 100.0 & 0.0 & -0.025 & -0.048 & 0.813 & 46.0 & 100.0 & 0.0 \\ 
Donation to AMS (\$) & education & 0.015 & 0.038 & 0.746 & 50.0 & 100.0 & 0.0 & -0.252 & -0.236 & 1.091 & 50.0 & 100.0 & 0.0 \\ 
Donation to AMS (\$) & gender & -0.049 & -0.070 & 0.463 & 47.5 & 100.0 & 0.0 & -0.034 & -0.054 & 0.508 & 55.0 & 100.0 & 0.0 \\ 
Donation to AMS (\$) & income & 0.196 & 0.181 & 0.633 & 52.5 & 100.0 & 0.0 & -0.109 & -0.133 & 0.911 & 60.0 & 100.0 & 0.0 \\ 
Donation to AMS (\$) & party & 0.224 & 0.239 & 0.639 & 43.3 & 100.0 & 0.0 & 0.098 & 0.138 & 0.640 & 50.0 & 100.0 & 0.0 \\ 
Donation to AMS (\$) & race & -0.087 & -0.055 & 0.587 & 46.2 & 100.0 & 0.0 & 0.125 & 0.182 & 0.606 & 48.8 & 100.0 & 0.0 \\ 
General climate policy & age\_band & -0.155 & -0.131 & 8.293 & 39.0 & 100.0 & 0.0 & 0.191 & 0.156 & 7.166 & 68.0 & 100.0 & 0.0 \\ 
General climate policy & education & 0.072 & 0.121 & 7.538 & 48.0 & 100.0 & 0.0 & 0.096 & 0.124 & 7.698 & 45.0 & 100.0 & 0.0 \\ 
General climate policy & gender & 0.174 & 0.186 & 3.313 & 55.0 & 100.0 & 0.0 & 0.008 & 0.030 & 4.845 & 50.0 & 100.0 & 0.0 \\ 
General climate policy & income & -0.240 & -0.208 & 6.008 & 40.0 & 100.0 & 0.0 & -0.130 & -0.115 & 7.394 & 47.5 & 100.0 & 0.0 \\ 
General climate policy & party & -0.009 & 0.023 & 6.723 & 48.3 & 100.0 & 0.0 & 0.060 & 0.071 & 7.394 & 51.7 & 100.0 & 0.0 \\ 
General climate policy & race & 0.255 & 0.260 & 4.828 & 56.2 & 100.0 & 0.0 & 0.126 & 0.074 & 6.766 & 53.8 & 100.0 & 0.0 \\ 
\bottomrule
\end{tabular*}

}

\end{table}%

\begin{Shaded}
\begin{Highlighting}[]
\FunctionTok{map\_dfr}\NormalTok{(moderators\_cat, }\ControlFlowTok{function}\NormalTok{(mod) \{}
  \FunctionTok{inner\_join}\NormalTok{(}
\NormalTok{    mod\_results\_h[[mod]]}\SpecialCharTok{$}\NormalTok{interaction\_effects }\SpecialCharTok{|\textgreater{}}
      \FunctionTok{select}\NormalTok{(condition, moderator\_level, }\AttributeTok{estimate\_h =}\NormalTok{ estimate),}
\NormalTok{    mod\_results\_l[[mod]]}\SpecialCharTok{$}\NormalTok{interaction\_effects }\SpecialCharTok{|\textgreater{}}
      \FunctionTok{select}\NormalTok{(condition, moderator\_level, }\AttributeTok{estimate\_l =}\NormalTok{ estimate),}
    \AttributeTok{by =} \FunctionTok{c}\NormalTok{(}\StringTok{"condition"}\NormalTok{, }\StringTok{"moderator\_level"}\NormalTok{)}
\NormalTok{  ) }\SpecialCharTok{|\textgreater{}} \FunctionTok{mutate}\NormalTok{(}\AttributeTok{moderator =}\NormalTok{ mod)}
\NormalTok{\}) }\SpecialCharTok{|\textgreater{}}
  \FunctionTok{ggplot}\NormalTok{(}\FunctionTok{aes}\NormalTok{(}\AttributeTok{x =}\NormalTok{ estimate\_h, }\AttributeTok{y =}\NormalTok{ estimate\_l)) }\SpecialCharTok{+}
  \FunctionTok{geom\_abline}\NormalTok{(}\AttributeTok{slope =} \DecValTok{1}\NormalTok{, }\AttributeTok{intercept =} \DecValTok{0}\NormalTok{, }\AttributeTok{linetype =} \StringTok{"dashed"}\NormalTok{,}
              \AttributeTok{color =} \StringTok{"grey60"}\NormalTok{, }\AttributeTok{linewidth =} \FloatTok{0.4}\NormalTok{, }\AttributeTok{alpha =} \FloatTok{0.35}\NormalTok{) }\SpecialCharTok{+}
  \FunctionTok{geom\_hline}\NormalTok{(}\AttributeTok{yintercept =} \DecValTok{0}\NormalTok{, }\AttributeTok{linetype =} \StringTok{"dotted"}\NormalTok{, }\AttributeTok{color =} \StringTok{"grey80"}\NormalTok{) }\SpecialCharTok{+}
  \FunctionTok{geom\_vline}\NormalTok{(}\AttributeTok{xintercept =} \DecValTok{0}\NormalTok{, }\AttributeTok{linetype =} \StringTok{"dotted"}\NormalTok{, }\AttributeTok{color =} \StringTok{"grey80"}\NormalTok{) }\SpecialCharTok{+}
  \FunctionTok{geom\_smooth}\NormalTok{(}\AttributeTok{method =} \StringTok{"lm"}\NormalTok{, }\AttributeTok{se =} \ConstantTok{FALSE}\NormalTok{, }\AttributeTok{color =} \StringTok{"\#2980B9"}\NormalTok{, }\AttributeTok{linewidth =} \FloatTok{0.6}\NormalTok{) }\SpecialCharTok{+}
  \FunctionTok{geom\_point}\NormalTok{(}\AttributeTok{size =} \FloatTok{1.5}\NormalTok{, }\AttributeTok{alpha =} \FloatTok{0.65}\NormalTok{, }\AttributeTok{color =} \StringTok{"grey20"}\NormalTok{) }\SpecialCharTok{+}
  \FunctionTok{facet\_wrap}\NormalTok{(}\SpecialCharTok{\textasciitilde{}}\NormalTok{ moderator, }\AttributeTok{scales =} \StringTok{"free"}\NormalTok{, }\AttributeTok{ncol =} \DecValTok{2}\NormalTok{) }\SpecialCharTok{+}
  \FunctionTok{labs}\NormalTok{(}\AttributeTok{x =} \StringTok{"Human interaction estimate"}\NormalTok{, }\AttributeTok{y =} \StringTok{"LLM clone interaction estimate"}\NormalTok{) }\SpecialCharTok{+}
\NormalTok{  plot\_theme}
\end{Highlighting}
\end{Shaded}

\begin{figure}[H]

\centering{

\pandocbounded{\includegraphics[keepaspectratio]{preregistration_files/figure-pdf/fig-rq3-1.pdf}}

}

\caption{\label{fig-rq3}Human vs.~LLM clone condition × moderator
interaction estimates for the primary outcome (multidimensional trust),
by moderator. Each point is one condition × moderator level combination.
Blue line: OLS regression. Faint dashed line: perfect calibration (slope
= 1).}

\end{figure}%

\subsubsection{Response Distributions by
Subgroup}\label{response-distributions-by-subgroup}

Mirroring the distribution analysis in Section 1, we compare response
distributions within each demographic subgroup in the control condition.
A clone dataset could reproduce the overall response distribution while
having miscalibrated within-group spreads. We use the same OVL, KS D,
and variance ratio metrics, in the same increasing order of strictness.
These analyses are run separately for all outcomes; the preregistration
demonstrates them on the primary outcome
(\texttt{trust\_multidimensional}). Minimum group size of 10 is required
for stable kernel density estimation.

\begin{Shaded}
\begin{Highlighting}[]
\CommentTok{\# For each moderator × level, compare human and clone distributions in control}
\NormalTok{subgroup\_dist }\OtherTok{\textless{}{-}} \FunctionTok{map\_dfr}\NormalTok{(moderators\_cat, }\ControlFlowTok{function}\NormalTok{(mod) \{}
\NormalTok{  groups }\OtherTok{\textless{}{-}} \FunctionTok{levels}\NormalTok{(human\_data[[mod]])}
  \FunctionTok{map\_dfr}\NormalTok{(groups, }\ControlFlowTok{function}\NormalTok{(grp) \{}
\NormalTok{    x }\OtherTok{\textless{}{-}}\NormalTok{ human\_data\_1 }\SpecialCharTok{|\textgreater{}}
      \FunctionTok{filter}\NormalTok{(condition }\SpecialCharTok{==}\NormalTok{ control\_val, .data[[mod]] }\SpecialCharTok{==}\NormalTok{ grp) }\SpecialCharTok{|\textgreater{}}
      \FunctionTok{pull}\NormalTok{(trust\_multidimensional) }\SpecialCharTok{|\textgreater{}} \FunctionTok{na.omit}\NormalTok{()}
\NormalTok{    y }\OtherTok{\textless{}{-}}\NormalTok{ llm\_data }\SpecialCharTok{|\textgreater{}}
      \FunctionTok{filter}\NormalTok{(condition }\SpecialCharTok{==}\NormalTok{ control\_val, .data[[mod]] }\SpecialCharTok{==}\NormalTok{ grp) }\SpecialCharTok{|\textgreater{}}
      \FunctionTok{pull}\NormalTok{(trust\_multidimensional) }\SpecialCharTok{|\textgreater{}} \FunctionTok{na.omit}\NormalTok{()}
    \ControlFlowTok{if}\NormalTok{ (}\FunctionTok{length}\NormalTok{(x) }\SpecialCharTok{\textless{}} \DecValTok{10} \SpecialCharTok{|} \FunctionTok{length}\NormalTok{(y) }\SpecialCharTok{\textless{}} \DecValTok{10}\NormalTok{) }\FunctionTok{return}\NormalTok{(}\ConstantTok{NULL}\NormalTok{)}
    \FunctionTok{tibble}\NormalTok{(}\AttributeTok{moderator =}\NormalTok{ mod, }\AttributeTok{group =}\NormalTok{ grp,}
           \AttributeTok{ovl =} \FunctionTok{compute\_ovl}\NormalTok{(x, y),}
           \AttributeTok{ks\_d =} \FunctionTok{suppressWarnings}\NormalTok{(}\FunctionTok{ks.test}\NormalTok{(x, y)}\SpecialCharTok{$}\NormalTok{statistic),}
           \AttributeTok{variance\_ratio =} \FunctionTok{var}\NormalTok{(y) }\SpecialCharTok{/} \FunctionTok{var}\NormalTok{(x))}
\NormalTok{  \})}
\NormalTok{\})}

\CommentTok{\# Human 1 vs. Human 2 baseline ceiling for within{-}subgroup distributions}
\NormalTok{subgroup\_dist\_h2 }\OtherTok{\textless{}{-}} \FunctionTok{map\_dfr}\NormalTok{(moderators\_cat, }\ControlFlowTok{function}\NormalTok{(mod) \{}
\NormalTok{  groups }\OtherTok{\textless{}{-}} \FunctionTok{levels}\NormalTok{(human\_data[[mod]])}
  \FunctionTok{map\_dfr}\NormalTok{(groups, }\ControlFlowTok{function}\NormalTok{(grp) \{}
\NormalTok{    x }\OtherTok{\textless{}{-}}\NormalTok{ human\_data\_1 }\SpecialCharTok{|\textgreater{}}
      \FunctionTok{filter}\NormalTok{(condition }\SpecialCharTok{==}\NormalTok{ control\_val, .data[[mod]] }\SpecialCharTok{==}\NormalTok{ grp) }\SpecialCharTok{|\textgreater{}}
      \FunctionTok{pull}\NormalTok{(trust\_multidimensional) }\SpecialCharTok{|\textgreater{}} \FunctionTok{na.omit}\NormalTok{()}
\NormalTok{    y }\OtherTok{\textless{}{-}}\NormalTok{ human\_data\_2 }\SpecialCharTok{|\textgreater{}}
      \FunctionTok{filter}\NormalTok{(condition }\SpecialCharTok{==}\NormalTok{ control\_val, .data[[mod]] }\SpecialCharTok{==}\NormalTok{ grp) }\SpecialCharTok{|\textgreater{}}
      \FunctionTok{pull}\NormalTok{(trust\_multidimensional) }\SpecialCharTok{|\textgreater{}} \FunctionTok{na.omit}\NormalTok{()}
    \ControlFlowTok{if}\NormalTok{ (}\FunctionTok{length}\NormalTok{(x) }\SpecialCharTok{\textless{}} \DecValTok{10} \SpecialCharTok{|} \FunctionTok{length}\NormalTok{(y) }\SpecialCharTok{\textless{}} \DecValTok{10}\NormalTok{) }\FunctionTok{return}\NormalTok{(}\ConstantTok{NULL}\NormalTok{)}
    \FunctionTok{tibble}\NormalTok{(}\AttributeTok{moderator =}\NormalTok{ mod, }\AttributeTok{group =}\NormalTok{ grp,}
           \AttributeTok{ovl =} \FunctionTok{compute\_ovl}\NormalTok{(x, y),}
           \AttributeTok{ks\_d =} \FunctionTok{suppressWarnings}\NormalTok{(}\FunctionTok{ks.test}\NormalTok{(x, y)}\SpecialCharTok{$}\NormalTok{statistic),}
           \AttributeTok{variance\_ratio =} \FunctionTok{var}\NormalTok{(y) }\SpecialCharTok{/} \FunctionTok{var}\NormalTok{(x))}
\NormalTok{  \})}
\NormalTok{\})}
\end{Highlighting}
\end{Shaded}

\begin{Shaded}
\begin{Highlighting}[]
\FunctionTok{map\_dfr}\NormalTok{(moderators\_cat, }\ControlFlowTok{function}\NormalTok{(mod) \{}
\NormalTok{  groups }\OtherTok{\textless{}{-}} \FunctionTok{levels}\NormalTok{(human\_data[[mod]])}
  \FunctionTok{map\_dfr}\NormalTok{(groups, }\ControlFlowTok{function}\NormalTok{(grp) \{}
    \FunctionTok{bind\_rows}\NormalTok{(}
\NormalTok{      human\_data\_1 }\SpecialCharTok{|\textgreater{}} \FunctionTok{filter}\NormalTok{(condition }\SpecialCharTok{==}\NormalTok{ control\_val, .data[[mod]] }\SpecialCharTok{==}\NormalTok{ grp) }\SpecialCharTok{|\textgreater{}}
        \FunctionTok{select}\NormalTok{(}\AttributeTok{value =}\NormalTok{ trust\_multidimensional) }\SpecialCharTok{|\textgreater{}}
        \FunctionTok{mutate}\NormalTok{(}\AttributeTok{source =} \StringTok{"Human"}\NormalTok{,     }\AttributeTok{panel =} \FunctionTok{paste0}\NormalTok{(mod, }\StringTok{": "}\NormalTok{, grp)),}
\NormalTok{      llm\_data     }\SpecialCharTok{|\textgreater{}} \FunctionTok{filter}\NormalTok{(condition }\SpecialCharTok{==}\NormalTok{ control\_val, .data[[mod]] }\SpecialCharTok{==}\NormalTok{ grp) }\SpecialCharTok{|\textgreater{}}
        \FunctionTok{select}\NormalTok{(}\AttributeTok{value =}\NormalTok{ trust\_multidimensional) }\SpecialCharTok{|\textgreater{}}
        \FunctionTok{mutate}\NormalTok{(}\AttributeTok{source =} \StringTok{"LLM clone"}\NormalTok{, }\AttributeTok{panel =} \FunctionTok{paste0}\NormalTok{(mod, }\StringTok{": "}\NormalTok{, grp))}
\NormalTok{    )}
\NormalTok{  \})}
\NormalTok{\}) }\SpecialCharTok{|\textgreater{}}
  \FunctionTok{filter}\NormalTok{(}\SpecialCharTok{!}\FunctionTok{is.na}\NormalTok{(value)) }\SpecialCharTok{|\textgreater{}}
  \FunctionTok{ggplot}\NormalTok{(}\FunctionTok{aes}\NormalTok{(}\AttributeTok{x =}\NormalTok{ value, }\AttributeTok{fill =}\NormalTok{ source, }\AttributeTok{color =}\NormalTok{ source)) }\SpecialCharTok{+}
  \FunctionTok{geom\_density}\NormalTok{(}\AttributeTok{alpha =} \FloatTok{0.35}\NormalTok{, }\AttributeTok{linewidth =} \FloatTok{0.4}\NormalTok{) }\SpecialCharTok{+}
  \FunctionTok{scale\_fill\_manual}\NormalTok{(}\AttributeTok{values  =}\NormalTok{ source\_colors, }\AttributeTok{name =} \ConstantTok{NULL}\NormalTok{) }\SpecialCharTok{+}
  \FunctionTok{scale\_color\_manual}\NormalTok{(}\AttributeTok{values =}\NormalTok{ source\_colors, }\AttributeTok{name =} \ConstantTok{NULL}\NormalTok{) }\SpecialCharTok{+}
  \FunctionTok{facet\_wrap}\NormalTok{(}\SpecialCharTok{\textasciitilde{}}\NormalTok{ panel, }\AttributeTok{scales =} \StringTok{"free\_y"}\NormalTok{, }\AttributeTok{ncol =} \DecValTok{4}\NormalTok{) }\SpecialCharTok{+}
  \FunctionTok{labs}\NormalTok{(}\AttributeTok{x =} \StringTok{"Multidimensional trust (0–100)"}\NormalTok{, }\AttributeTok{y =} \StringTok{"Density"}\NormalTok{) }\SpecialCharTok{+}
\NormalTok{  plot\_theme }\SpecialCharTok{+}
  \FunctionTok{theme}\NormalTok{(}\AttributeTok{legend.position =} \StringTok{"top"}\NormalTok{)}
\end{Highlighting}
\end{Shaded}

\begin{figure}[H]

\centering{

\pandocbounded{\includegraphics[keepaspectratio]{preregistration_files/figure-pdf/fig-rq4c-densities-1.pdf}}

}

\caption{\label{fig-rq4c-densities}Response distributions within
demographic subgroups in the control condition (primary outcome:
multidimensional trust). Human participants in dark grey, LLM clones in
orange.}

\end{figure}%

\begin{Shaded}
\begin{Highlighting}[]
\NormalTok{subgroup\_dist }\SpecialCharTok{|\textgreater{}}
  \FunctionTok{left\_join}\NormalTok{(}
\NormalTok{    subgroup\_dist\_h2 }\SpecialCharTok{|\textgreater{}} \FunctionTok{select}\NormalTok{(moderator, group, }\AttributeTok{ovl\_h2 =}\NormalTok{ ovl, }\AttributeTok{ks\_d\_h2 =}\NormalTok{ ks\_d, }\AttributeTok{vr\_h2 =}\NormalTok{ variance\_ratio),}
    \AttributeTok{by =} \FunctionTok{c}\NormalTok{(}\StringTok{"moderator"}\NormalTok{, }\StringTok{"group"}\NormalTok{)}
\NormalTok{  ) }\SpecialCharTok{|\textgreater{}}
  \FunctionTok{arrange}\NormalTok{(moderator, group) }\SpecialCharTok{|\textgreater{}}
  \FunctionTok{gt}\NormalTok{(}\AttributeTok{groupname\_col =} \StringTok{"moderator"}\NormalTok{) }\SpecialCharTok{|\textgreater{}}
  \FunctionTok{cols\_label}\NormalTok{(}\AttributeTok{group =} \StringTok{"Group"}\NormalTok{,}
             \AttributeTok{ovl =} \StringTok{"OVL"}\NormalTok{, }\AttributeTok{ks\_d =} \StringTok{"KS D"}\NormalTok{, }\AttributeTok{variance\_ratio =} \StringTok{"Var. ratio"}\NormalTok{,}
             \AttributeTok{ovl\_h2 =} \StringTok{"OVL"}\NormalTok{, }\AttributeTok{ks\_d\_h2 =} \StringTok{"KS D"}\NormalTok{, }\AttributeTok{vr\_h2 =} \StringTok{"Var. ratio"}\NormalTok{) }\SpecialCharTok{|\textgreater{}}
  \FunctionTok{tab\_spanner}\NormalTok{(}\AttributeTok{label =} \StringTok{"Human 1 vs. LLM"}\NormalTok{,     }\AttributeTok{columns =} \FunctionTok{c}\NormalTok{(ovl, ks\_d, variance\_ratio)) }\SpecialCharTok{|\textgreater{}}
  \FunctionTok{tab\_spanner}\NormalTok{(}\AttributeTok{label =} \StringTok{"Human 1 vs. Human 2"}\NormalTok{, }\AttributeTok{columns =} \FunctionTok{c}\NormalTok{(ovl\_h2, ks\_d\_h2, vr\_h2)) }\SpecialCharTok{|\textgreater{}}
  \FunctionTok{fmt\_number}\NormalTok{(}\AttributeTok{columns =} \FunctionTok{c}\NormalTok{(ovl, ks\_d, variance\_ratio, ovl\_h2, ks\_d\_h2, vr\_h2), }\AttributeTok{decimals =} \DecValTok{3}\NormalTok{) }\SpecialCharTok{|\textgreater{}}
  \FunctionTok{tab\_style}\NormalTok{(}\AttributeTok{style =} \FunctionTok{cell\_text}\NormalTok{(}\AttributeTok{weight =} \StringTok{"bold"}\NormalTok{),}
            \AttributeTok{locations =} \FunctionTok{cells\_row\_groups}\NormalTok{())}
\end{Highlighting}
\end{Shaded}

\begin{table}

\caption{\label{tbl-rq4c}Within-subgroup distribution comparison in the
control condition. OVL: overlapping coefficient (1 = identical). KS D:
Kolmogorov--Smirnov statistic. Variance ratio: second group / Human 1
variance. Left columns: Human 1 vs.~LLM; right columns: Human 1
vs.~Human 2 (baseline ceiling).}

\centering{

\fontsize{12.0pt}{14.4pt}\selectfont
\begin{tabular*}{\linewidth}{@{\extracolsep{\fill}}lrrrrrr}
\toprule
 & \multicolumn{3}{c}{Human 1 vs. LLM} & \multicolumn{3}{c}{Human 1 vs. Human 2} \\ 
\cmidrule(lr){2-4} \cmidrule(lr){5-7}
Group & OVL & KS D & Var. ratio & OVL & KS D & Var. ratio \\ 
\midrule\addlinespace[2.5pt]
\multicolumn{7}{l}{{\bfseries age\_band}} \\[2.5pt] 
\midrule\addlinespace[2.5pt]
18-29 & 0.916 & 0.070 & 1.336 & 0.903 & 0.119 & 1.135 \\ 
30-39 & 0.938 & 0.081 & 1.195 & 0.877 & 0.114 & 1.387 \\ 
40-49 & 0.931 & 0.072 & 0.918 & 0.861 & 0.167 & 0.886 \\ 
50-59 & 0.928 & 0.083 & 1.049 & 0.865 & 0.131 & 0.867 \\ 
60-69 & 0.909 & 0.091 & 0.917 & 0.931 & 0.073 & 0.958 \\ 
70+ & 0.920 & 0.077 & 0.992 & 0.929 & 0.081 & 1.094 \\ 
\midrule\addlinespace[2.5pt]
\multicolumn{7}{l}{{\bfseries education}} \\[2.5pt] 
\midrule\addlinespace[2.5pt]
Bachelor's degree & 0.890 & 0.120 & 1.314 & 0.868 & 0.149 & 1.196 \\ 
Doctorate degree / Ph.D. & 0.898 & 0.128 & 1.268 & 0.876 & 0.123 & 1.231 \\ 
High school diploma / GED & 0.925 & 0.059 & 1.198 & 0.850 & 0.163 & 1.247 \\ 
Less than high school & 0.949 & 0.079 & 1.102 & 0.917 & 0.099 & 0.905 \\ 
Master's degree / Professional degree & 0.930 & 0.072 & 0.883 & 0.924 & 0.083 & 0.834 \\ 
Some college or Associate's degree & 0.903 & 0.111 & 0.813 & 0.896 & 0.101 & 0.843 \\ 
\midrule\addlinespace[2.5pt]
\multicolumn{7}{l}{{\bfseries gender}} \\[2.5pt] 
\midrule\addlinespace[2.5pt]
Female & 0.916 & 0.070 & 1.319 & 0.894 & 0.124 & 1.149 \\ 
Male & 0.954 & 0.058 & 0.993 & 0.946 & 0.056 & 1.037 \\ 
Other & 0.932 & 0.044 & 0.908 & 0.941 & 0.053 & 0.892 \\ 
\midrule\addlinespace[2.5pt]
\multicolumn{7}{l}{{\bfseries income}} \\[2.5pt] 
\midrule\addlinespace[2.5pt]
\$100,000 to \$167,999 & 0.944 & 0.058 & 1.262 & 0.883 & 0.116 & 1.138 \\ 
\$168,000 or more & 0.908 & 0.118 & 1.026 & 0.915 & 0.100 & 1.105 \\ 
\$30,000 to \$55,999 & 0.954 & 0.061 & 1.145 & 0.876 & 0.125 & 1.223 \\ 
\$56,000 to \$99,999 & 0.951 & 0.051 & 0.966 & 0.918 & 0.071 & 0.874 \\ 
Less than \$30,000 & 0.936 & 0.069 & 1.006 & 0.940 & 0.061 & 0.828 \\ 
\midrule\addlinespace[2.5pt]
\multicolumn{7}{l}{{\bfseries party}} \\[2.5pt] 
\midrule\addlinespace[2.5pt]
Democrat & 0.932 & 0.083 & 0.954 & 0.933 & 0.100 & 1.002 \\ 
Independent & 0.945 & 0.047 & 1.011 & 0.929 & 0.079 & 0.935 \\ 
Other & 0.929 & 0.064 & 1.060 & 0.917 & 0.072 & 1.032 \\ 
Republican & 0.909 & 0.078 & 1.241 & 0.895 & 0.113 & 1.145 \\ 
\midrule\addlinespace[2.5pt]
\multicolumn{7}{l}{{\bfseries race}} \\[2.5pt] 
\midrule\addlinespace[2.5pt]
Asian / Asian American & 0.930 & 0.089 & 1.206 & 0.897 & 0.127 & 0.983 \\ 
Black / African American & 0.908 & 0.068 & 0.743 & 0.908 & 0.092 & 0.844 \\ 
Hispanic / Latino & 0.873 & 0.152 & 1.420 & 0.880 & 0.107 & 1.264 \\ 
Other & 0.947 & 0.066 & 1.017 & 0.913 & 0.102 & 0.937 \\ 
White / Caucasian & 0.935 & 0.081 & 1.018 & 0.915 & 0.074 & 1.107 \\ 
\bottomrule
\end{tabular*}

}

\end{table}%

\subsection{3. Demographic Baseline
Calibration}\label{demographic-baseline-calibration}

Section 2 already provides a partial test of demographic stereotyping:
if clones exaggerate heterogeneity in treatment effects across
subgroups, that is evidence of over-reliance on demographic cues. But
there is a second, orthogonal question namely whether clones get the
baseline right: For example, the clones could still assign far too low
trust to, say, Republicans or far too high trust to Democats before any
treatment is applied. This woud be a distinct failure mode ---
demographic baseline stereotyping rather than effect stereotyping.

Section 3 tests this directly, using only the control condition. It asks
two questions: (a) do clone group means match human group means for each
demographic group? and (b) do demographic variables predict clone
outcomes too strongly --- does the demographic profile explain more
variance in clone responses than in human responses, indicating that
clones map demographics to outcomes too mechanically? Both analyses are
run separately for all outcomes; the preregistration illustrates them
with the primary outcome (\texttt{trust\_multidimensional}).

\subsubsection{Demographic Baseline
Calibration}\label{demographic-baseline-calibration-1}

For each moderator, we compute mean outcomes within each demographic
group in the control condition separately for humans and clones, then
correlate and compare those cell means. High correlation and low RMSE
indicate that clones start from the right baseline response level for
each demographic type --- before any intervention is applied.

\begin{Shaded}
\begin{Highlighting}[]
\CommentTok{\# Cell{-}mean comparison in the control condition, one moderator at a time}
\NormalTok{baseline\_cal }\OtherTok{\textless{}{-}} \FunctionTok{compare\_demographic\_baselines}\NormalTok{(}
\NormalTok{  human\_data\_1, llm\_data,}
  \AttributeTok{outcome    =} \StringTok{"trust\_multidimensional"}\NormalTok{,}
  \AttributeTok{moderators =}\NormalTok{ moderators\_cat}
\NormalTok{)}

\CommentTok{\# Human 1 vs. Human 2 baseline ceiling for demographic calibration}
\NormalTok{baseline\_cal\_h2 }\OtherTok{\textless{}{-}} \FunctionTok{compare\_demographic\_baselines}\NormalTok{(}
\NormalTok{  human\_data\_1, human\_data\_2,}
  \AttributeTok{outcome    =} \StringTok{"trust\_multidimensional"}\NormalTok{,}
  \AttributeTok{moderators =}\NormalTok{ moderators\_cat}
\NormalTok{)}
\end{Highlighting}
\end{Shaded}

\begin{Shaded}
\begin{Highlighting}[]
\NormalTok{baseline\_cal }\SpecialCharTok{|\textgreater{}}
  \FunctionTok{left\_join}\NormalTok{(}
\NormalTok{    baseline\_cal\_h2 }\SpecialCharTok{|\textgreater{}} \FunctionTok{select}\NormalTok{(moderator, }\AttributeTok{r\_h2 =}\NormalTok{ r, }\AttributeTok{rmse\_h2 =}\NormalTok{ rmse),}
    \AttributeTok{by =} \StringTok{"moderator"}
\NormalTok{  ) }\SpecialCharTok{|\textgreater{}}
  \FunctionTok{select}\NormalTok{(}\AttributeTok{Moderator =}\NormalTok{ moderator, r, }\AttributeTok{RMSE =}\NormalTok{ rmse, }\StringTok{\textasciigrave{}}\AttributeTok{N cells}\StringTok{\textasciigrave{}} \OtherTok{=}\NormalTok{ n\_cells, r\_h2, rmse\_h2) }\SpecialCharTok{|\textgreater{}}
  \FunctionTok{gt}\NormalTok{() }\SpecialCharTok{|\textgreater{}}
  \FunctionTok{cols\_label}\NormalTok{(}\AttributeTok{r\_h2 =} \StringTok{"r"}\NormalTok{, }\AttributeTok{rmse\_h2 =} \StringTok{"RMSE"}\NormalTok{) }\SpecialCharTok{|\textgreater{}}
  \FunctionTok{tab\_spanner}\NormalTok{(}\AttributeTok{label =} \StringTok{"Human 1 vs. LLM"}\NormalTok{,     }\AttributeTok{columns =} \FunctionTok{c}\NormalTok{(r, RMSE)) }\SpecialCharTok{|\textgreater{}}
  \FunctionTok{tab\_spanner}\NormalTok{(}\AttributeTok{label =} \StringTok{"Human 1 vs. Human 2"}\NormalTok{, }\AttributeTok{columns =} \FunctionTok{c}\NormalTok{(r\_h2, rmse\_h2)) }\SpecialCharTok{|\textgreater{}}
  \FunctionTok{fmt\_number}\NormalTok{(}\AttributeTok{columns =} \FunctionTok{c}\NormalTok{(r, RMSE, r\_h2, rmse\_h2), }\AttributeTok{decimals =} \DecValTok{3}\NormalTok{)}
\end{Highlighting}
\end{Shaded}

\begin{table}

\caption{\label{tbl-rq4a}Demographic baseline calibration. For each
moderator, cell means in the control condition are compared. r: Pearson
correlation of cell means. RMSE: root mean squared error. N cells:
number of demographic groups. Left columns: Human 1 vs.~LLM; right
columns: Human 1 vs.~Human 2 (baseline ceiling).}

\centering{

\fontsize{12.0pt}{14.4pt}\selectfont
\begin{tabular*}{\linewidth}{@{\extracolsep{\fill}}lrrrrr}
\toprule
 & \multicolumn{2}{c}{Human 1 vs. LLM} &  & \multicolumn{2}{c}{Human 1 vs. Human 2} \\ 
\cmidrule(lr){2-3} \cmidrule(lr){5-6}
Moderator & r & RMSE & N cells & r & RMSE \\ 
\midrule\addlinespace[2.5pt]
gender & -0.981 & 0.438 & 3 & 0.020 & 0.553 \\ 
age\_band & 0.564 & 0.520 & 6 & -0.296 & 0.878 \\ 
race & -0.264 & 0.894 & 5 & 0.446 & 0.912 \\ 
education & -0.714 & 0.898 & 6 & -0.149 & 1.083 \\ 
income & -0.528 & 0.636 & 5 & 0.481 & 0.630 \\ 
party & -0.726 & 0.855 & 4 & -0.852 & 0.904 \\ 
\bottomrule
\end{tabular*}

}

\end{table}%

\subsubsection{Demographic
Predictability}\label{demographic-predictability}

For each of the six demographic moderators we fit a separate OLS
regression of the primary outcome on that moderator plus condition fixed
effects, on human and clone data independently. Running these
regressions separately avoids conflating correlated predictors (race,
income, and education are themselves related) and gives a clean
per-moderator answer: how much does, say, partisan identity alone
explain outcome variance in humans versus clones? If clones over-rely on
demographic stereotypes, their R² will be substantially higher than
humans' for one or more moderators. We also compare the individual
dummy-coded coefficients to identify which groups are over- or
underweighted.

\begin{Shaded}
\begin{Highlighting}[]
\NormalTok{pred\_test }\OtherTok{\textless{}{-}} \FunctionTok{compare\_demographic\_predictability}\NormalTok{(}
\NormalTok{  human\_data\_1, llm\_data,}
  \AttributeTok{outcome    =} \StringTok{"trust\_multidimensional"}\NormalTok{,}
  \AttributeTok{predictors =}\NormalTok{ moderators\_cat}
\NormalTok{)}

\CommentTok{\# Human 1 vs. Human 2 baseline ceiling for demographic predictability}
\NormalTok{pred\_test\_h2 }\OtherTok{\textless{}{-}} \FunctionTok{compare\_demographic\_predictability}\NormalTok{(}
\NormalTok{  human\_data\_1, human\_data\_2,}
  \AttributeTok{outcome    =} \StringTok{"trust\_multidimensional"}\NormalTok{,}
  \AttributeTok{predictors =}\NormalTok{ moderators\_cat}
\NormalTok{)}
\end{Highlighting}
\end{Shaded}

\begin{Shaded}
\begin{Highlighting}[]
\NormalTok{pred\_test}\SpecialCharTok{$}\NormalTok{r\_squared }\SpecialCharTok{|\textgreater{}}
  \FunctionTok{pivot\_wider}\NormalTok{(}\AttributeTok{names\_from =}\NormalTok{ source, }\AttributeTok{values\_from =}\NormalTok{ r\_squared) }\SpecialCharTok{|\textgreater{}}
  \FunctionTok{left\_join}\NormalTok{(}
\NormalTok{    pred\_test\_h2}\SpecialCharTok{$}\NormalTok{r\_squared }\SpecialCharTok{|\textgreater{}}
      \FunctionTok{filter}\NormalTok{(source }\SpecialCharTok{==} \StringTok{"llm"}\NormalTok{) }\SpecialCharTok{|\textgreater{}}
      \FunctionTok{select}\NormalTok{(moderator, }\StringTok{\textasciigrave{}}\AttributeTok{Human 2 R²}\StringTok{\textasciigrave{}} \OtherTok{=}\NormalTok{ r\_squared),}
    \AttributeTok{by =} \StringTok{"moderator"}
\NormalTok{  ) }\SpecialCharTok{|\textgreater{}}
  \FunctionTok{rename}\NormalTok{(}\AttributeTok{Moderator =}\NormalTok{ moderator, }\StringTok{\textasciigrave{}}\AttributeTok{Human 1 R²}\StringTok{\textasciigrave{}} \OtherTok{=}\NormalTok{ human, }\StringTok{\textasciigrave{}}\AttributeTok{LLM R²}\StringTok{\textasciigrave{}} \OtherTok{=}\NormalTok{ llm) }\SpecialCharTok{|\textgreater{}}
  \FunctionTok{gt}\NormalTok{() }\SpecialCharTok{|\textgreater{}}
  \FunctionTok{fmt\_number}\NormalTok{(}\AttributeTok{columns =} \FunctionTok{c}\NormalTok{(}\StringTok{\textasciigrave{}}\AttributeTok{Human 1 R²}\StringTok{\textasciigrave{}}\NormalTok{, }\StringTok{\textasciigrave{}}\AttributeTok{LLM R²}\StringTok{\textasciigrave{}}\NormalTok{, }\StringTok{\textasciigrave{}}\AttributeTok{Human 2 R²}\StringTok{\textasciigrave{}}\NormalTok{), }\AttributeTok{decimals =} \DecValTok{3}\NormalTok{)}
\end{Highlighting}
\end{Shaded}

\begin{table}

\caption{\label{tbl-rq4b-rsq}R² from regressing the primary outcome on
each demographic moderator + condition FE. Clone R²
\textgreater\textgreater{} Human R² signals over-reliance on demographic
cues. Human 2 R² provides the baseline ceiling (sampling variation in
human predictability).}

\centering{

\fontsize{12.0pt}{14.4pt}\selectfont
\begin{tabular*}{\linewidth}{@{\extracolsep{\fill}}lrrr}
\toprule
Moderator & Human 1 R² & LLM R² & Human 2 R² \\ 
\midrule\addlinespace[2.5pt]
gender & 0.002 & 0.001 & 0.004 \\ 
age\_band & 0.002 & 0.002 & 0.004 \\ 
race & 0.002 & 0.002 & 0.004 \\ 
education & 0.002 & 0.002 & 0.004 \\ 
income & 0.002 & 0.002 & 0.004 \\ 
party & 0.002 & 0.002 & 0.004 \\ 
\bottomrule
\end{tabular*}

}

\end{table}%

\begin{Shaded}
\begin{Highlighting}[]
\NormalTok{pred\_test}\SpecialCharTok{$}\NormalTok{coefficients }\SpecialCharTok{|\textgreater{}}
  \FunctionTok{ggplot}\NormalTok{(}\FunctionTok{aes}\NormalTok{(}\AttributeTok{x =}\NormalTok{ est\_h, }\AttributeTok{y =}\NormalTok{ est\_l)) }\SpecialCharTok{+}
  \FunctionTok{geom\_abline}\NormalTok{(}\AttributeTok{slope =} \DecValTok{1}\NormalTok{, }\AttributeTok{intercept =} \DecValTok{0}\NormalTok{, }\AttributeTok{linetype =} \StringTok{"dashed"}\NormalTok{,}
              \AttributeTok{color =} \StringTok{"grey60"}\NormalTok{, }\AttributeTok{linewidth =} \FloatTok{0.5}\NormalTok{) }\SpecialCharTok{+}
  \FunctionTok{geom\_hline}\NormalTok{(}\AttributeTok{yintercept =} \DecValTok{0}\NormalTok{, }\AttributeTok{linetype =} \StringTok{"dotted"}\NormalTok{, }\AttributeTok{color =} \StringTok{"grey80"}\NormalTok{) }\SpecialCharTok{+}
  \FunctionTok{geom\_vline}\NormalTok{(}\AttributeTok{xintercept =} \DecValTok{0}\NormalTok{, }\AttributeTok{linetype =} \StringTok{"dotted"}\NormalTok{, }\AttributeTok{color =} \StringTok{"grey80"}\NormalTok{) }\SpecialCharTok{+}
  \FunctionTok{geom\_point}\NormalTok{(}\AttributeTok{size =} \DecValTok{2}\NormalTok{, }\AttributeTok{alpha =} \FloatTok{0.8}\NormalTok{, }\AttributeTok{color =} \StringTok{"grey20"}\NormalTok{) }\SpecialCharTok{+}
  \FunctionTok{geom\_text}\NormalTok{(}\FunctionTok{aes}\NormalTok{(}\AttributeTok{label =}\NormalTok{ term), }\AttributeTok{size =} \FloatTok{2.2}\NormalTok{, }\AttributeTok{vjust =} \SpecialCharTok{{-}}\FloatTok{0.8}\NormalTok{, }\AttributeTok{color =} \StringTok{"grey40"}\NormalTok{) }\SpecialCharTok{+}
  \FunctionTok{facet\_wrap}\NormalTok{(}\SpecialCharTok{\textasciitilde{}}\NormalTok{ moderator, }\AttributeTok{scales =} \StringTok{"free"}\NormalTok{) }\SpecialCharTok{+}
  \FunctionTok{labs}\NormalTok{(}\AttributeTok{x =} \StringTok{"Human coefficient"}\NormalTok{, }\AttributeTok{y =} \StringTok{"LLM clone coefficient"}\NormalTok{) }\SpecialCharTok{+}
\NormalTok{  plot\_theme}
\end{Highlighting}
\end{Shaded}

\begin{figure}[H]

\centering{

\pandocbounded{\includegraphics[keepaspectratio]{preregistration_files/figure-pdf/fig-rq4b-coefs-1.pdf}}

}

\caption{\label{fig-rq4b-coefs}Demographic regression coefficients for
the primary outcome in humans (x-axis) vs.~LLM clones (y-axis), faceted
by moderator. Each point is one dummy-coded level. The dashed line marks
perfect agreement (slope = 1).}

\end{figure}%

\phantomsection\label{refs}
\begin{CSLReferences}{1}{0}
\bibitem[\citeproctext]{ref-aherUsingLargeLanguage2023}
Aher, Gati, Rosa I. Arriaga, and Adam Tauman Kalai. 2023. {``Using Large
Language Models to Simulate Multiple Humans and Replicate Human Subject
Studies,''} July. \url{https://doi.org/10.48550/arXiv.2208.10264}.

\bibitem[\citeproctext]{ref-almaatouqPlaying20Questions2022a}
Almaatouq, Abdullah, Thomas L. Griffiths, Jordan W. Suchow, Mark E.
Whiting, James Evans, and Duncan J. Watts. 2022. {``Beyond Playing 20
Questions with Nature: Integrative Experiment Design in the Social and
Behavioral Sciences.''} \emph{Behavioral and Brain Sciences}, December,
1--55. \url{https://doi.org/10.1017/S0140525X22002874}.

\bibitem[\citeproctext]{ref-argyleOutOneMany2023}
Argyle, Lisa P., Ethan C. Busby, Nancy Fulda, Joshua R. Gubler,
Christopher Rytting, and David Wingate. 2023. {``Out of One, Many: Using
Language Models to Simulate Human Samples.''} \emph{Political Analysis}
31 (3): 337--51. \url{https://doi.org/10.1017/pan.2023.2}.

\bibitem[\citeproctext]{ref-binzUsingCognitivePsychology2023}
Binz, Marcel, and Eric Schulz. 2023. {``Using Cognitive Psychology to
Understand GPT-3.''} \emph{Proceedings of the National Academy of
Sciences} 120 (6): e2218523120.
\url{https://doi.org/10.1073/pnas.2218523120}.

\bibitem[\citeproctext]{ref-bisbeeSyntheticReplacementsHuman2024}
Bisbee, James, Joshua D. Clinton, Cassy Dorff, Brenton Kenkel, and
Jennifer M. Larson. 2024. {``Synthetic Replacements for Human Survey
Data? The Perils of Large Language Models.''} \emph{Political Analysis}
32 (4): 401--16. \url{https://doi.org/10.1017/pan.2024.5}.

\bibitem[\citeproctext]{ref-chenPredictingFieldExperiments2025}
Chen, Yaoyu, Yuheng Hu, and Yingda Lu. 2025. {``Predicting Field
Experiments with Large Language Models,''} May.
\url{https://doi.org/10.48550/arXiv.2504.01167}.

\bibitem[\citeproctext]{ref-cuiCanAIReplace2024}
Cui, Ziyan, Ning Li, and Huaikang Zhou. 2024. {``Can AI Replace Human
Subjects? A Large-Scale Replication of Psychological Experiments with
LLMs.''} \url{https://doi.org/10.2139/ssrn.4940173}.

\bibitem[\citeproctext]{ref-doudkinAIPersuadingAI2025}
Doudkin, Alexander, Pat Pataranutaporn, and Pattie Maes. 2025. {``AI
Persuading AI Vs AI Persuading Humans: LLMs' Differential Effectiveness
in Promoting Pro-Environmental Behavior,''} March.
\url{https://doi.org/10.48550/arXiv.2503.02067}.

\bibitem[\citeproctext]{ref-durmusMeasuringRepresentationSubjective2024}
Durmus, Esin, Karina Nguyen, Thomas I. Liao, Nicholas Schiefer, Amanda
Askell, Anton Bakhtin, Carol Chen, et al. 2024. {``Towards Measuring the
Representation of Subjective Global Opinions in Language Models,''}
April. \url{https://doi.org/10.48550/arXiv.2306.16388}.

\bibitem[\citeproctext]{ref-guoEstimatingCausalEffects2025}
Guo, Siyi, Myrl G. Marmarelis, Fred Morstatter, and Kristina Lerman.
2025. {``Estimating Causal Effects of Text Interventions Leveraging
LLMs,''} March. \url{https://doi.org/10.48550/arXiv.2410.21474}.

\bibitem[\citeproctext]{ref-hewittPredictingResultsSocial2025}
Hewitt, Luke, Ashwini Ashokkumar, Isaias Ghezae, and Robb Willer. 2025.
{``Predicting Results of Social Science Experiments Using Large Language
Models.''}

\bibitem[\citeproctext]{ref-huSimBenchBenchmarkingAbility2025}
Hu, Tiancheng, Joachim Baumann, Lorenzo Lupo, Nigel Collier, Dirk Hovy,
and Paul Röttger. 2025. {``SimBench: Benchmarking the Ability of Large
Language Models to Simulate Human Behaviors,''} October.
\url{https://doi.org/10.48550/arXiv.2510.17516}.

\bibitem[\citeproctext]{ref-jonesPeopleCannotDistinguish2025}
Jones, Cameron Robert, Ishika Rathi, Sydney Taylor, and Benjamin K.
Bergen. 2025. {``FAccT '25: The 2025 ACM Conference on Fairness,
Accountability, and Transparency.''} In, 1615--39. Athens Greece: ACM.
\url{https://doi.org/10.1145/3715275.3732108}.

\bibitem[\citeproctext]{ref-kaiserSimulatingHumanOpinions2025}
Kaiser, Carolin, Jakob Kaiser, Vladimir Manewitsch, Lea Rau, and Rene
Schallner. 2025. {``Simulating Human Opinions with Large Language
Models: Opportunities and Challenges for Personalized Survey Data
Modeling.''} In, 8286. UMAP Adjunct '25. New York, NY, USA: Association
for Computing Machinery. \url{https://doi.org/10.1145/3708319.3733685}.

\bibitem[\citeproctext]{ref-lippertCanLargeLanguage2024}
Lippert, Steffen, Anna Dreber, Magnus Johannesson, Warren Tierney,
Wilson Cyrus-Lai, Eric Luis Uhlmann, null null, and Thomas Pfeiffer.
2024. {``Can Large Language Models Help Predict Results from a Complex
Behavioural Science Study?''} \emph{Royal Society Open Science} 11 (9):
240682. \url{https://doi.org/10.1098/rsos.240682}.

\bibitem[\citeproctext]{ref-nosekScientificUtopiaII2012}
Nosek, Brian A., Jeffrey R. Spies, and Matt Motyl. 2012. {``Scientific
Utopia: II. Restructuring Incentives and Practices to Promote Truth Over
Publishability.''} \emph{Perspectives on Psychological Science} 7 (6):
615--31. \url{https://doi.org/10.1177/1745691612459058}.

\bibitem[\citeproctext]{ref-paganComputationalTuringTest2025}
Pagan, Nicolò, Petter Törnberg, Christopher A. Bail, Anikó Hannák, and
Christopher Barrie. 2025. {``Computational Turing Test Reveals
Systematic Differences Between Human and AI Language,''} November.
\url{https://doi.org/10.48550/arXiv.2511.04195}.

\bibitem[\citeproctext]{ref-santurkarWhoseOpinionsLanguage2023}
Santurkar, Shibani, Esin Durmus, Faisal Ladhak, Cinoo Lee, Percy Liang,
and Tatsunori Hashimoto. 2023. {``Whose Opinions Do Language Models
Reflect?''} March. \url{https://doi.org/10.48550/arXiv.2303.17548}.

\bibitem[\citeproctext]{ref-schroderLargeLanguageModels2025a}
Schröder, Sarah, Thekla Morgenroth, Ulrike Kuhl, Valerie Vaquet, and
Benjamin Paaßen. 2025. {``Large Language Models Do Not Simulate Human
Psychology,''} August. \url{https://doi.org/10.48550/arXiv.2508.06950}.

\bibitem[\citeproctext]{ref-shresthaWEIRDCanSynthetic2024}
Shrestha, Pujen, Dario Krpan, Fatima Koaik, Robin Schnider, Dima Sayess,
and May Saad Binbaz. 2024. {``Beyond WEIRD: Can Synthetic Survey
Participants Substitute for Humans in Global Policy Research?''}
\emph{Behavioral Science \& Policy} 10 (2): 26--45.
\url{https://doi.org/10.1177/23794607241311793}.

\bibitem[\citeproctext]{ref-smaldinoNaturalSelectionBad2016}
Smaldino, Paul E., and Richard McElreath. 2016. {``The Natural Selection
of Bad Science.''} \emph{Royal Society Open Science} 3 (9): 160384.
\url{https://doi.org/10.1098/rsos.160384}.

\bibitem[\citeproctext]{ref-wangLargeLanguageModels2025}
Wang, Angelina, Jamie Morgenstern, and John P. Dickerson. 2025. {``Large
Language Models That Replace Human Participants Can Harmfully Misportray
and Flatten Identity Groups,''} February.
\url{https://doi.org/10.48550/arXiv.2402.01908}.

\end{CSLReferences}




\end{document}
